% before: cv-hermitian-norms
The covolume bound compares lattices of forms over a CM field whose metrics come from
Petersson norms at the complex places. All comparisons in this chapter are made through one
normalized slope, and the bound itself is stated for the normalized log-covolume. The two
are the same number up to the factor $-\operatorname{rk}L$.

% before: cv-rescaling
The lattice whose covolume is bounded is written in coordinates, while the lower bounds for
slopes are proved for the same forms inside a larger space. Passing between the two
loses nothing when the coordinate metric is a contraction of the ambient one. This is the
first step of the covolume bound.

% before: cv-slope-base-change
The lifted forms have coefficients in the coefficient field of the point, while the automorphic
forms they come from have coefficients in the smaller field of the Shimura variety. Extending
scalars with the same norms does not change slopes, and the lemma below measures the loss when
the norms agree only up to a bounded factor.

% before: cv-cut
The lattice of lifted forms of the prescribed type contains a small part on which the
level group acts through characters that are not separated. Cutting it off by an idempotent of
the group algebra costs a little slope, controlled by the rank of the discarded part and by
an upper bound for slopes of sublattices. The cost below is linear in the proportion of
discarded rank.

% before: cv-cut-index
The inequality between the slopes of $L$, $L_g$ and $L_b$ in the proof above compares covolumes
through the image of $g(1-e)$. The same inequality has a shorter derivation from the index of
$L_g\oplus L_b$ in $L$, since the two summands are orthogonal.

% before: cv-slope-descent
Lower bounds for slopes of automorphic forms of one weight are transferred to another
weight by multiplying with sections. The transfer needs a slope inequality for several maps
at once, with a loss measured by the norms of the maps, by the size of a spanning set of
the target, and by how far the ranks are from balanced.

% before: cv-minkowski-slope
An upper bound for the maximal slope of a lattice is needed to cut off the part of the lifted
forms that is not separated. Minkowski's theorem turns a lower bound for the norms of all nonzero
vectors into such an upper bound, with a loss of the logarithm of the rank.

% before: cv-mass
On the definite side the forms live on finite sets of lattice classes, and their number is
controlled by masses. The mass weights each class in a genus by the inverse order of its
automorphism group.

% before: cv-siegel-mass-ratio
The chain of definite unitary groups grows in rank by one at each step of the telescoping
below, so only the ratio of consecutive masses of standard lattices is needed. For a CM
field unramified over its real subfield this ratio is an explicit product of factorials,
powers of the root discriminant and a special value of an $L$-function.

% before: cv-transference
The upper bound for the mass counts classes of self-dual lattices, and each class is found from a
basis of short vectors. For the trace form, a self-dual lattice has no short vectors in its dual
lattice. A transference theorem turns this into short vectors in the lattice itself.

% before: cv-short-basis
A self-dual lattice over a CM field has a basis whose Gram matrix is bounded in terms of the rank
and the root discriminant alone. The reason is the different: the Hermitian norm of a vector in
the trace dual lies in the inverse square of the different, so that vector cannot be short.

% before: cv-class-count
A lattice is recovered, up to isometry, from the Gram matrix of a basis and the module of the
values of the form against that basis. Once the Gram matrix is bounded, both run over finite sets
whose sizes are explicit, and this bounds the number of classes.

% before: cv-mass-bounds
The volume of a unitary group over a number field is the mass of a standard lattice up to
explicit archimedean factors. Only crude bounds for that mass are needed: the class of the
standard lattice gives a lower bound, and the count of classes above gives an upper bound.

% before: cv-volume-mass-haar
The measure on $U(P,Q)$ in the volume identity below is fixed by its density at the identity, a
normalization that makes sense in every signature. The compact group $U(N)$ also carries its
probability measure, and the identity is proved with the measures on $U(P,Q)$ that have the
density of that probability measure. The two normalizations differ by the volume of $U(N)$.

% before: cv-volume-mass-normalized
The stabilizer of a self-dual lattice at the finite places and these measures at the real places
determine a Haar measure on the adelic unitary group. The measure it gives a fundamental domain of
the rational points is the volume that the identity compares with a mass.

% before: cv-volume-mass-local-lattices
At a finite place where $E/F$ does not ramify, every self-dual lattice looks like the standard one.
It has an orthonormal basis, any two such lattices are conjugate, and the stabilizer of one acts
transitively on its vectors of a given unit norm.

% before: cv-volume-mass-invariance
The volume does not depend on the choices made in defining it. Another self-dual lattice has a
conjugate stabilizer, by the local facts just stated, and an isometry of spaces carries every
ingredient of the definition to the other space.

% before: cv-volume-mass-landherr
The comparison of signatures builds auxiliary spaces and needs them to be isometric to the given
ones. Over a CM field unramified over its real subfield, a self-dual lattice makes the local
determinants norms, so rank and signatures decide the isometry class.

% before: cv-volume-mass-unit-vector
A step of the comparison splits off a line on which the form is $\pm1$. In the diagonal basis just
found, a basis vector spans such a line, and the span of the basis is a self-dual lattice that
splits along it.

% before: cv-volume-mass-definite
In a definite space the adelic quotient is a finite union of pieces indexed by the classes in the
genus of the lattice. Each piece has volume the reciprocal of the order of an automorphism group,
so the volume is the mass.

% before: cv-volume-mass-unfolding
In an indefinite space the volume is compared with that of the orthogonal complement of a vector of
norm $\pm1$. The theta series of the self-dual lattice for that norm, integrated over a
fundamental domain and unfolded along the orbit of the vector, gives the volume of the complement
times an archimedean integral.

% before: cv-volume-mass-siegel-weil
The same integral of the theta series is computed by Siegel's formula as a product of local
densities. This is where the arithmetic of the lattice enters the comparison of volumes.

% before: cv-volume-mass-local-density
At a finite place the local density of a self-dual lattice for a unit norm comes from a count of
solutions modulo the place. It depends only on whether the place splits in $E$.

% before: cv-volume-mass-euler-product
The local densities multiply to a power of the discriminant divided by the $L$-value that appears
in the mass ratio. The discriminant of $F$ enters through the conductor of the additive character.

% before: cv-volume-mass-density-gauges
The constant in the integration formula at a real place is found by comparing densities at the
identity. Near the identity the group is measured by the sum of the absolute values of the entries of
$g-1$, and near a basis vector the quadric by a similar sum of coordinates. These gauges have to be
compared with each other and with the Lebesgue measures that define the densities.

% before: cv-volume-mass-archimedean
At a real place the unfolded integral is an integral over the quadric of vectors of norm $\pm1$.
With the measures of the compact dual, the constant relating the group, the stabilizer and the
quadric is the same in every signature, so it can be read off from the volume of a sphere.

% before: cv-volume-mass-ratio
The unfolding and Siegel's formula compute the same integral, and their comparison gives the ratio
of the volumes of a space and of the complement of a vector of norm $\pm1$. The ratio is the right
side of the mass ratio formula, in every signature.

% before: cv-volume-mass
The norms of lifted forms are Petersson norms rescaled by the volume of a fundamental domain
of the indefinite group, normalized by its level subgroup. Theorem~\ref{cv-volume-mass}
expresses that volume through the mass of the standard lattice, in every signature.

% before: cv-volume-mass-literature
In the definite case the identity just proved computes a mass as a volume. Mass formulas of this
kind are usually obtained from volume formulas for arithmetic groups, and the comparison names them.

% before: cv-level-sizes
The upper bound for the volume below holds under a hypothesis on one real number $s$. It bounds the
ranks and the weights of the datum by powers of $e^s$, and the logarithm of the level primes, the
logarithmic root discriminant and a logarithmic derivative of an $L$-function by multiples of $s$.
At the points of a family, $s$ is twice $x^\varepsilon$.

% before: cv-level-volume
The covolume bound needs the logarithm of that volume to be of the size of the error term,
above and below. The level subgroups of the data have index bounded by a power of the level
primes in the stabilizer of a self-dual lattice, and the mass bounds do the rest.

% before: cv-class-number-formula
The rank of the space of algebraic modular forms at the bottom of a chain is the number of
lattice classes of the level, weighted by the dimensions of the archimedean types. When a
level prime is large, no class has a nontrivial automorphism of that level, and the count
becomes the mass times an explicit index.

% before: cv-where-mass-enters
The overview of the argument lists the mass formula among the inputs to the upper bound for
the auxiliary lattice. That bound combines the covolume bound of this chapter with the norms
of lifts and the counts of classes on the definite side. Masses enter both, in different
forms.

% before: cv-chowla-selberg
The main term of the covolume bound is a multiple of the logarithm of a period. That period
is the Chowla--Selberg period of the imaginary quadratic field $\mathbf Q(\sqrt{-d})$
contained in the host field. Its logarithm is an average over ideal classes, which the
Chowla--Selberg formula evaluates in terms of $L'/L(1,\chi)$.

% before: cv-log-omega
The covolume bound and the flat degree formula use the same normalization of the
Chowla--Selberg period. It replaces $\sqrt d$ in the formula by the root discriminant of
the host field.

% before: cv-constant-sections
The slopes of automorphic forms of one weight are compared with those of another weight by
multiplying with sections whose norms are constant on each component. Such a section has a
degree, the archimedean and finite logarithms of its norms summed with the product formula in
mind.

% before: cv-constant-sections-proportional
A degree attached to a component is useful only if it does not depend on the section chosen.
Two constant sections of the same weight and defect differ by a constant, because a function
with constant absolute value at every complex point is algebraic over the field of definition.

% before: cv-gaga
The flat sections of the next theorem are first found as holomorphic functions on the complex
points of a moduli space. A holomorphic section of a line bundle on a proper smooth complex
variety is algebraic, and this is how such a function becomes a rational section defined over the
coefficient field.

% before: cv-resolution-coordinates
The algebraicity statement needs a proper smooth variety, while the moduli spaces are only open in
their closures. Resolving the closure and making the boundary a divisor with normal crossings gives
such a variety, together with local coordinates in which the boundary is a union of coordinate
hyperplanes; near the boundary a holomorphic function is then studied one coordinate at a time.

% before: cv-cm-point-periods
The degree of a flat section is computed at a single point of the component, where its norms can
be read off from periods. Over the complex numbers the component is parametrized by the periods of
a Hermitian lattice, and the points used are those at which the lattice splits into two pieces whose
periods lie in complementary sets of coordinates. The Gram matrices of the lattice and the Hodge
Gram matrices of tangent vectors carry the archimedean and the finite norms.

% before: cv-cm-point-gram
The norms at such a point are expressed through Gram determinants of families in the lattice. Their
size at a prime is governed by the different of $E$, because the $E$-valued form is the trace dual
of a perfect alternating form, and the two pieces of a split lattice are orthogonal.

% before: cv-cm-point-curve
The split point is assembled from two elliptic curves with complex multiplication by the ring of
integers of $\mathbf Q(\sqrt{-d})$, one for each square root of $-d$. The archimedean norms of their
tangent vectors are the periods to which the Chowla--Selberg formula applies.

% before: cv-cm-point-models
The CM curves are built from elliptic curves over a ring of integers whose complex points are the
lattices of the ideal classes of $\mathbf Q(\sqrt{-d})$. Taking one curve for each class, together with
the homomorphisms between them, allows the mean of the Hodge norms over all embeddings to be computed
one class at a time.

% before: cv-cm-point-area
The squared Hodge norm of a tangent vector of an elliptic curve is $\pi$ divided by the area of its
period lattice. The discriminants of the curves are units and their homomorphisms are integral, so
summing over the embeddings turns the mean of the logarithms of these areas into a mean over the ideal
classes, to which the Chowla--Selberg formula applies.

% before: cv-cm-point-action
The curve obtained so far carries an endomorphism of its generic fibre with derivative
$(1+i\sqrt d)/2$. It becomes a CM curve once this endomorphism generates an action of $\mathbf Z[w_0]$
over the ring of integers whose Rosati involution is complex conjugation, and conjugating the action
gives the curve for the other root.

% before: cv-cm-point-split
At a split point the lattice has a local basis whose first $p$ vectors have periods in the coordinates
of $\theta$ and whose last $q$ vectors have periods in the others. Such a basis is found by perturbing a
local basis $\ell$-adically until its first $p$ vectors have positive definite Gram matrices at every
embedding over $\theta$, and then splitting off their orthogonal complement.

% before: cv-cm-point-descent
At the split point the abelian variety receives homomorphisms from the two CM curves, one for each
vector of the pure sublattices. They show that the point is defined over a number field and that the
abelian variety has good reduction everywhere, and at a prime $\ell$ they identify its Lie algebra with
that of a product of curves.

% before: cv-cm-point-trace
Tangent vectors at the split point are built from the images of the tangent vectors of the CM curves,
summed against the trace-dual basis of $\mathcal O_E$. This sum projects onto one eigenspace of
$\mathcal O_E$, and it turns the Hermitian forms pulled back from the curves into the $E$-valued form of
the lattice.

% before: cv-cm-point-eigen
At a prime $\ell$ the Lie algebra of the model has a basis of images of tangent vectors of curves, and
the eigenspace of $\tau$ cuts a sublattice out of it. Projecting onto that eigenspace divides by the
relative different of $E$ over $\mathbf Q(\sqrt{-d})$, so multiplying by a generator of that different
gives a basis of the sublattice.

% before: cv-cm-point
At a point where the lattice splits, the abelian variety is built from the two CM curves, and frames
of its Lie algebra come from their tangent vectors. The Hodge Gram matrices of these frames at the
complex places and their integral structure at a prime are then explicit, and the curves can be
chosen with the mean period of Theorem~\ref{cv-chowla-selberg}(1). Step 2 of the next proof computes
the local norms of a flat section from these data.

% before: cv-function-situation
The flat sections of Theorem~\ref{cv-flat-degree} are built on one connected component of the moduli
scheme at a time, from the Lie algebra of the universal abelian scheme. The action of
$\mathcal O_E$ cuts out the eigenspaces of an embedding $\tau$ of the CM type and of its conjugate,
and the weight $\lambda_\tau$ is carried by their determinants.

% before: cv-function-framed
A section of the determinant line of the two eigenspaces is recorded by its values in chosen bases at
every point. Under a change of bases these values change as a section of the determinant of the first
eigenspace divided by one of the second, and they are compatible with enlarging the field of
definition of the point.

% before: cv-model
The norms of a flat section at a finite place are read at integral points, so the universal abelian
scheme has to extend over a model of the moduli scheme. The extension is a semi-abelian scheme on a
normal projective scheme over a discrete valuation ring, abelian over the points of the moduli scheme
and over no other point of the generic fibre.

% before: cv-semistable-reduction
The models of the moduli schemes over discrete valuation rings are built from the reduction of the
abelian varieties they parametrize. Over a finite extension of $\mathbf Q_p$, an abelian variety
need not extend to an abelian scheme over the ring of integers, but after a finite extension it
extends to a semi-abelian scheme.

% before: cv-model-toroidal-semiabelian
The models are built in stages, over rings of integers in which fewer and fewer primes are inverted
and over moduli schemes at several levels. At each stage the object is a normal projective scheme over
a Dedekind domain with a semi-abelian scheme that carries the structures of the datum, and three
properties of that semi-abelian scheme are carried from stage to stage.

% before: cv-model-toroidal-auxiliary
The first stage inverts an auxiliary prime and uses a full level structure at a power of it. The prime
is chosen inert in $\mathbf Q(\sqrt{-d})$ and unramified in $E$, so that it is prime to the level of
the datum, and two such primes are needed so that one of them differs from the prime at which the
model is wanted.

% before: cv-model-toroidal-extension
Homomorphisms between semi-abelian schemes over a discrete valuation ring are determined by their
restriction to the generic fibre, and every homomorphism of generic fibres extends. This is how group
actions and changes of level structure pass from the moduli schemes to their compactifications.

% before: cv-toroidal-family
The first stage is a compactification of the moduli schemes whose full level is a power of an
auxiliary prime $n$, over a ring of integers in which $n$ is inverted. It is built over one scheme
$\mathcal M$ that carries an abelian scheme with the structures of the datum and a full level
structure, and that classifies such abelian schemes over reduced flat bases. The moduli scheme at
every level of this full level maps to $\mathcal M$.

% before: cv-toroidal-theta
The scheme $\mathcal M$ is built from theta constants, as in Mumford's construction of moduli of
abelian varieties. A principally polarized abelian scheme with a symplectic basis of its points of
order $N'^2$ has, locally on the base, a family of theta constants of level $N'$. The theta family is
cut out by the polynomial relations among the theta constants of the complex points of the datum,
together with a condition that keeps enough of them away from zero, and its points are determined by
their theta constants.

% before: cv-toroidal-mumford
Over a complete discrete valuation ring, a principally polarized abelian variety whose semi-abelian
reduction has a split torus is given by Mumford's construction from degeneration data. The
construction describes its theta functions and its torsion points near the closed point. Combined
with the semistable reduction theorem, it controls the theta constants of points of the theta family
over $p$-adic fields.

% before: cv-toroidal-reduction-situation
Lemma~\ref{cv-toroidal-theta-reduction} below is proved in three stages: a finite extension over which
the abelian variety has semistable reduction with a split torus, a proof with Mumford's construction
that this reduction is good, and the extension of the theta-level object and its theta constants over
the larger ring. All of its intermediate statements concern theta-level objects for one choice of the
data of Definition~\ref{cv-toroidal-theta}.

% before: cv-toroidal-reduction-semistable
The first stage is the semistable reduction theorem in the form that Theorem~\ref{cv-toroidal-mumford}
needs: over a complete discrete valuation ring with finite residue field, and with a split torus in the
closed fibre. The torus splits over a finite extension of the residue field, which comes from a finite
extension of the fraction field.

% before: cv-toroidal-reduction-base-change
After the extension, the theta-level object and its theta constants must be carried to the larger
field. Theta-level objects pull back along any morphism of the base and theta constants along any ring
homomorphism, and the conditions on $\alpha_w$ and $\alpha_{w'}$ in
Lemma~\ref{cv-toroidal-theta-reduction} survive the pullback.

% before: cv-toroidal-reduction-local
Theta constants exist locally on the base of any theta-level object of odd level. The data of
Definition~\ref{cv-toroidal-theta}(2) come from the theorem of the cube, from a generator of the
sections of the principal line bundle, and from symmetric lifts of points of odd order.

% before: cv-toroidal-reduction-coset
At no prime do the theta constants vanish on a whole coset of the subgroup
$(N'/n)(\mathbf Z/N'^2)^{2g}$. Over an algebraically closed field this says that the zero locus of a
nonzero section of the principal line bundle contains no translate of the points of order dividing $n$.

% before: cv-toroidal-reduction-unit
Two families of theta constants at the same point differ by a unit as soon as one of them generates the
unit ideal. Locally the data of Definition~\ref{cv-toroidal-theta}(2) are unique up to units of the
base, because a point of odd order has only one symmetric lift of that order.

% before: cv-toroidal-reduction-generic
The last three lemmas combine to extend theta constants over a discrete valuation ring. If a
theta-level object over the fraction field extends over the ring, then integral theta constants at the
generic point that include a unit are theta constants over the ring; this is the last step in the proof
of Lemma~\ref{cv-toroidal-theta-reduction}.

% before: cv-toroidal-reduction-model
A theta-level object over the fraction field of a discrete valuation ring extends over the ring once its
principally polarized abelian variety does. The torsion points extend by properness and the
endomorphisms by Theorem~\ref{cv-model-toroidal-extension}; the conditions on them spread from the
generic point, because $N'$ is invertible and the base is connected.

% before: cv-toroidal-reduction-abelian
Good reduction is read off from degeneration data. When the free $\mathbf Z$-module $X$ of the data is
zero, the torus of Mumford's construction is trivial, and the abelian variety is the generic fibre of the
abelian scheme of the data.

% before: cv-toroidal-torus-valuation
The proof that $X$ is zero compares valuations of theta constants at torsion points of $A$. These points
are defined only over an algebraic closure of the fraction field, so the valuation of the discrete
valuation ring has to be extended there, with rational values.

% before: cv-toroidal-torus-situation
The proof that $X$ is zero takes place over a complete discrete valuation ring with degeneration data and
a uniformization, and over an algebraically closed field containing its fraction field, with a valuation.
On the side of the degeneration data it uses the valuations of the trivializations $\psi(\chi)$; on the
side of the abelian variety it uses the theta constants of $L$.

% before: cv-toroidal-torus-quadratic
The valuations of the trivializations $\psi(\chi)$ form an integer-valued function $Q$ on $X$. It vanishes
at $0$, by the bimultiplicativity of $\tau$, and it is positive elsewhere, because $\psi(\chi)$ lies in
the maximal ideal for $\chi\ne0$.

% before: cv-toroidal-torus-points
Torsion points of $A$ over an algebraically closed field are described by division points of the
degeneration data (Theorem~\ref{cv-toroidal-mumford}(2)), which are points of the Raynaud extension, the
extension of $B$ by the torus with character group $X$. Sections of powers of $L$ are described by Fourier
coefficients indexed by $X$, whose valuations are measured against the integral sections over $B$.

% before: cv-toroidal-torus-extension
Every element $\chi$ of $X$ has to be the first component of a division point of each order. This needs
the divisibility of the group of points of the Raynaud extension over an algebraically closed field, which
comes from the divisibility of $B(F)$ and of the torus.

% before: cv-toroidal-torus-expansion
The uniformization identifies $L$ with the line bundle of $G$ on the generic fibre only up to a constant.
Fixing that constant at the origin makes the Fourier expansion of a section of $L^k$ over $A_F$ well
defined, and comparable with theta constants normalized at the origin.

% before: cv-toroidal-torus-normalization
A normalization exists because two trivializations of one line over $K$ differ by a constant. All Fourier
expansions in the proof that $X$ is zero are taken with respect to one normalization.

% before: cv-toroidal-torus-sections
Valuations of Fourier coefficients are measured against the sections of the line bundles $M_{k,\chi}$ over
$B$. These sections form a free module of rank $k^g$ over $W$ whose bases stay bases after every base
change, so a section over $B_F$ has coordinates in $F$ with respect to sections over $B$.

% before: cv-toroidal-torus-coordinates
In these coordinates the valuation bound of Definition~\ref{cv-toroidal-torus-points}(4) becomes a bound on
each coordinate. A section over $B$ has coordinates in $W$, which have nonnegative valuation, so writing a
section with other integral sections cannot improve the bound.

% before: cv-toroidal-torus-translation
The functional equation of a family of Fourier coefficients goes through isomorphisms between translates of
the line bundles $M_{k,\chi}$. They are built from the cubical structure of $M$ and fixed by their values at
the origin. Up to a power of $\varpi$ they are defined over $B$, and that power measures how the
valuations of the coefficients change along a class modulo $kX$.

% before: cv-toroidal-torus-classes
On each class of $X$ modulo $kX$, a family of Fourier coefficients is determined by any one of its
components. Every section at one point of a class extends to a family supported on that class, so families
can be built one class at a time.

% before: cv-toroidal-torus-shift
Moving along a class modulo $kX$ changes the valuation bound of a coefficient by the change of $Q/k$. So a
bound of the form $\nu(\sigma_\chi)\ge r+Q(\chi)/k$ for all $\chi$ needs to be checked only at one
representative of each class.

% before: cv-toroidal-torus-rational
The valuation on the algebraically closed field $F$ takes every rational value. So the elements of
valuation at least $q$ are the multiples of one element by elements of nonnegative valuation.

% before: cv-toroidal-torus-ball
The families with $\nu(\sigma_\chi)\ge Q(\chi)/k$ for all $\chi$ form a free module over the elements of $F$
of nonnegative valuation. A basis comes from the coordinates of the components at representatives of the
classes modulo $kX$, each scaled by an element of the right valuation.

% before: cv-toroidal-torus-base-change
Sections over $A_F$ are compared with sections over $A$ by extending scalars from $K$ to $F$. Both the
sections of a line bundle and the families of Fourier coefficients over $F$ are obtained from those over $K$
in this way; for the families this holds because they are determined by finitely many components.

% before: cv-toroidal-torus-bijection
The uniformization identifies the sections of $L^k$ over the generic fibre with families of Fourier
coefficients over $K$. Extending scalars to $F$ keeps this identification, so every section of $L^k$ over
$A_F$ has one Fourier expansion, and every family over $F$ is the expansion of one section.

% before: cv-toroidal-torus-primitive
The theta constants are computed from a generator of the sections of $L$ over $A_F$. The generator whose
Fourier expansion is primitive has a fixed scale: its Fourier coefficients have valuation at least
$Q(\chi)$, and no uniform positive improvement of these bounds holds.

% before: cv-toroidal-torus-normalized
Theta constants in the sense of Definition~\ref{cv-toroidal-theta}(2) are determined only up to a common
unit. Normalizing the isomorphism $[m]^*L\cong L^{m^2}$ and the trivialization at the origin by one
trivialization $\rho$ of $L$ removes this unit, so that a single theta constant has a valuation.

% before: cv-toroidal-torus-constant
For odd $m$ the normalized theta constant is well defined. The normalized isomorphism
$[m]^*L\cong L^{m^2}$ is unique because the global functions on an abelian variety are constant, and a
point of odd order has only one symmetric lift of that order.

% before: cv-toroidal-torus-family
Normalized theta constants at several torsion points use one isomorphism $[m]^*L\cong L^{m^2}$. So, for a
generator of the sections of $L$, they form a family of theta constants in the sense of
Definition~\ref{cv-toroidal-theta}(2), and Lemma~\ref{cv-toroidal-reduction-unit} compares them with any
other such family.

% before: cv-toroidal-torus-terms
The normalized theta constant at a torsion point is computed as a series indexed by $X$. Its terms pair a
Fourier coefficient at a point of $B$ with the torus coordinates of a point of $\tilde G(F)$ and with a
trivialization at that point. The valuations of these trivializations are measured against trivializations
over the ring of elements of nonnegative valuation.

% before: cv-toroidal-torus-sums
A series that converges for the valuation has valuation at least the minimum of the valuations of its
terms, and valuation $0$ when one term has valuation $0$ and all others have positive valuation. The series
for theta constants converge $\mathfrak m$-adically in a finitely generated module, which implies
convergence for the valuation.

% before: cv-toroidal-torus-symmetric
The trivialization at the point of $B$ under an $N$-division point is fixed by two conditions, one on its
square and one on its $N$-th power. Both use isomorphisms $[N]^*M\cong M^{N^2}$ normalized at the origin.

% before: cv-toroidal-torus-symmetric-frame
For odd $N$ the two conditions determine the trivialization. Its existence comes from identities of the
cubical structure along the multiples of the division point, and its uniqueness from the fact that an
element whose square and odd power are $1$ is $1$.

% before: cv-toroidal-torus-theta-sum-phi
The theta constant is the value at the origin of a section $\theta s_1$, where $\theta$ lies in the theta
group over a torsion point. On Fourier expansions such an element acts by translating each coefficient by
the point of $B$ under a division point and shifting its index by $y$. The isomorphism that carries the
translated coefficient to its new place is fixed by its value at the origin.

% before: cv-toroidal-torus-theta-sum-operator
The isomorphisms in this operator come from the cubical structure, as for the translations by points
$c(\mu)$. At the zero point of $\tilde G(F)$ the theta terms of the transformed family are then the theta terms
of the original family at the division point, with shifted index.

% before: cv-toroidal-torus-theta-sum-origin
The value of a section at the origin is recovered from its Fourier expansion. Modulo each power of
$\mathfrak m$ the uniformization writes a section over $G$ as a finite sum of coefficients times characters,
and the characters are $1$ at the origin.

% before: cv-toroidal-torus-expansion-coefficients
Modulo each power of $\mathfrak m$, the uniformization makes $G$ an extension of $B$ by a split torus, with a
trivialization for each character. A section of $L_G^k$ over $G$ is then a finite sum of pulled back sections of
the bundles $M_{k,\chi}$ times these trivializations, with one family of coefficients for all powers of
$\mathfrak m$. These families are the Fourier coefficients from which the expansions over $F$ are built.

% before: cv-toroidal-torus-expansion-endomorphism
The pullback along multiplication by $m$ is computed through endomorphisms of all the data of the uniformization
at once. Such an endomorphism acts on $G$, $A$ and $B$, on the character group $X$ and on the trivializations of
the characters, compatibly with the extensions modulo every power of $\mathfrak m$. It acts on the points of
$\tilde G(F)$ by the same formulas.

% before: cv-toroidal-torus-expansion-multiples
Multiplication by an integer on all the data is such an endomorphism. Its effect on the trivializations of the
characters is a product of values of the cubical structure, and it acts on $\tilde G(F)$ by multiplication by the
integer. The cases $m\ge0$ and $-1$ are the ones used below.

% before: cv-toroidal-torus-descent-structures
The descent below compares trivializations over $G$ with trivializations over $B$ through the isomorphisms $w_n$,
and it uses more of the uniformization than Definition~\ref{cv-toroidal-torus-expansion-coefficients} lists. The
line bundles $L_G$ and $M$ carry cubical structures and symmetries, and each $w_n$ carries those of $L_G$ to those
of $M$. The quotient maps $q_n$ are flat and affine, which is what descent along them needs.

% before: cv-toroidal-torus-descent-functions
Two trivializations of one line bundle on $G$ that agree at the origin agree everywhere when every global function
on $G$ is constant. The generic fibre of $G$ is an abelian variety, and as $G$ is flat over $W$ its global
functions, and those of $G\times_WG\times_WG$, are determined there.

% before: cv-toroidal-torus-descent-cube
A trivialization of a combination of pullbacks of $L_G$ along endomorphisms is controlled by the cubical structure
and the symmetry of $L_G$: both sides of each identity below trivialize one line bundle and agree at the origin.
Through the isomorphisms $w_n$, the same identities hold modulo each power of $\mathfrak m$ with the cubical
structure and the symmetry of $M$ at points of $B$.

% before: cv-toroidal-torus-descent-character
Under translation by the torus, such a trivialization changes by an invertible function on the product of the
torus with a connected scheme. A function of this kind that is multiplicative in the torus variable is a character,
by the duality between diagonalizable groups and their character groups.

% before: cv-toroidal-torus-descent-ratio
The criterion below shows that the change under translation by the torus is trivial. Its three conditions say that
the cube expression of the values does not see points of the torus and that the values at a point of the torus and
at its inverse agree; the cubical structure and the symmetry give them. The proof applies the previous lemma twice,
first to the function divided by its restriction to the origin and then to the function itself.

% before: cv-toroidal-torus-descent-thickenings
The trivializations descended to the reductions $B_n$ have to come from one trivialization over $B$. Normalized
along the zero section they are compatible, and the comparison theorem for the proper scheme $B$ over the complete
ring $W$ gives a section over $B$. It generates the line bundle along the closed fibre, hence everywhere, as $B$ is
proper.

% before: cv-toroidal-torus-expansion-descent
A trivialization of $L_G^{k'}\otimes a_G^*L_G^{-k}$ over $G$ becomes, modulo each power of $\mathfrak m$, a
trivialization of a line bundle pulled back from $B$. It descends to $B$, and the descended trivializations come
from one trivialization over $B$. This trivialization is the factor by which the coefficients of a pulled back
section change.

% before: cv-toroidal-torus-expansion-integral
For sections over $G$, pulling back along an endomorphism and twisting by a trivialization of
$L_G^{k'}\otimes a_G^*L_G^{-k}$ acts on the Fourier coefficients one by one. The coefficient of index $\chi$ moves to
the index $h(\chi)$, pulled back along $a_B$ and multiplied by the descended trivialization and by $u(\chi)$.

% before: cv-toroidal-torus-expansion-twist
Over $F$, the twisted pullback on $A$ has to be compared with the one on $G$. The two twists are trivializations
of the same line bundle on $G_F\cong A_F$ that agree at the origin, so they agree everywhere.

% before: cv-toroidal-torus-expansion-pullback
Together, the last three lemmas compute the Fourier expansion of a twisted pullback along an endomorphism that is
injective on $X$. In terms of theta terms, the term of index $h(\chi)$ at a point $g$ is the term of index $\chi$
of the original expansion at the image of $g$, and the other terms vanish.

% before: cv-toroidal-torus-expansion-phi-map
The operator $\Phi_{p,w}$ of Definition~\ref{cv-toroidal-torus-theta-sum-phi} is a linear map on families.
Changing the frame multiplies it by a constant, and this is how the frame that comes from the uniformization is
compared with the symmetric frame.

% before: cv-toroidal-torus-expansion-phi-order
For a symmetric frame, applying $\Phi_{p,w}$ $N$ times translates the coefficients by $Nb=c(y)$, which the
functional equation of the families also does. The second condition on the symmetric frame makes the two
isomorphisms agree at the origin.

% before: cv-toroidal-torus-expansion-assigned
The point of $A(F)$ that the uniformization gives for a division point is characterized by the action of the
theta group on Fourier expansions. Each element of the theta group over the point acts through the operator
$\Phi_{p,w}$ for some frame $w$, which is not specified.

% before: cv-toroidal-torus-expansion-frame-square
A symmetric element of the theta group is inverted by conjugation with the inversion of $A$. Comparing the two
ways in which the inversion acts on Fourier expansions, one of them through the endomorphism $-1$ of the
uniformization, determines the square of the frame of its operator.

% before: cv-toroidal-torus-expansion-symmetric-lift
For a symmetric element of order dividing $N$ over the point that the uniformization assigns to a division point,
the frame of its operator is the symmetric frame. Its square and its $N$-th power satisfy the two conditions that
determine the symmetric frame, unless all families vanish.

% before: cv-toroidal-torus-theta-sum-expansion
The section of $L^{m^2}$ that gives the theta constant is obtained from the generator by pulling back along
multiplication by $m$ and acting by the symmetric lift. On Fourier expansions the first step spreads the
coefficients of the generator over $mX$, and the second shifts them by $y$.

% before: cv-toroidal-torus-theta-sum
The normalized theta constant at the torsion point given by a division point $(y,g)$ is the sum of the theta
terms of the Fourier expansion of the generator at the point $mg$. This comes from expanding the twisted
section of $L^{m^2}$ that defines the theta constant, and evaluating it at the origin.

% before: cv-toroidal-torus-term-rules
The valuations of the theta terms are computed from valuations of trivializations of fibres of $M$ at
points of $B(F)$. Every such point extends uniquely to the valuation ring of $F$, so a trivialization is a
scalar times one over the valuation ring, and the valuation of the scalar is well defined. Trivializations
defined over $B$ or over $B^3$, such as the cubical structure, have valuation $0$ at every point.

% before: cv-toroidal-torus-term-points
The trivializations $s(\chi)$ of a point of $\tilde G(F)$ are linked by the cubical structure, whose values
have valuation $0$. So their valuations form an additive function of $\chi$, and this function is additive
in the point as well.

% before: cv-toroidal-torus-term-division
For a $k$-division point $(y,g)$ the $k$-th power of $s(\alpha)$ is $\tau(y,\alpha)$ up to values of the
cubical structure along the multiples of $g$. Since $\tau(y,\alpha)$ has valuation $b(y,\alpha)$, the point
$g$ has valuations $b(y,\cdot)/k$.

% before: cv-toroidal-torus-term-bound
A theta term is linear in the Fourier coefficient. For a coefficient defined over $B$ its valuation is at
least the sum of the valuations of $s(\chi)$ and of the frame, which bounds the terms of any family by the
valuations of its coefficients.

% before: cv-toroidal-torus-term-valuations
The valuations of the theta terms are controlled by those of $\psi$ and $\tau$. For a family whose
coefficients have valuations at least $Q$, the term indexed by $\chi$ has valuation at least
$Q(m\chi+y)/m^2$.

% before: cv-toroidal-torus-series
The sum formula and the valuations together give a series for the theta constant whose terms have bounded
valuations. This is the form in which both parts of Lemma~\ref{cv-toroidal-torus-valuations} use it.

% before: cv-toroidal-torus-exact-section
An exact valuation for one theta term needs a Fourier coefficient that does not vanish at the point where
it is evaluated. For weight $1$ the sections of $M_{1,\chi}$ over $B$ have one generator, the translate of a
generator of the sections of $M$. Its zero locus on a geometric fibre contains no translate of the points of
order dividing $n$, by the corresponding statement for a nonzero section of $M$.

% before: cv-toroidal-torus-exact-lift
The point of order dividing $n$ is found over the residue field of the valuation ring of $F$, and it has to
be lifted to the valuation ring. The kernel of multiplication by $n$ is finite and flat, and over a
valuation ring of an algebraically closed field every point of a finite flat scheme over the residue field
lifts.

% before: cv-toroidal-torus-exact-value
At a point whose integral extension avoids the zeros of the generator, the valuation of a theta term can be
read off. The values of the generator, of the trivialization $s(\chi)$ and of the frame are scalars times
trivializations over the valuation ring, and the product of these trivializations is a unit times the
trivialization of $M$ at the origin.

% before: cv-toroidal-torus-exact
A theta constant of valuation $0$ needs one term of the series with exact valuation. Adding a point of order
dividing $n$ to the point of $\tilde G(F)$ moves the Fourier coefficient to a point at whose reduction its
integral section does not vanish.

% before: cv-toroidal-torus-valuations
The Fourier expansion of the primitive generator bounds the valuations of its normalized theta constants.
At the point given by a division point $(y,g)$ of order $m^2$, the valuation is at least the minimum of
$Q(m\chi+y)/m^2$, which is positive unless $y\in mX$; and some division point with first component $0$
gives a theta constant of valuation $0$.

% before: cv-toroidal-torus-free
The unit condition on the theta constants makes every element of $X$ divisible by $n^{j-1}$. A basis
element of a nonzero free $\mathbf Z$-module is divisible by no integer $d\ge2$, so this forces $X=0$.

% before: cv-toroidal-reduction-torus
The theta constants force $X$ to be zero. The uniformization expresses the theta functions as Fourier
series indexed by $X$, so the valuations of the theta constants at the torsion points are given by the
degeneration data. A unit in every coset of $n^{j-1}(\mathbf Z/N'^2)^{2g}$ then makes every element of
$X$ divisible by $n^{j-1}$, and a free $\mathbf Z$-module of finite rank with this property is zero.

% before: cv-toroidal-theta-reduction
The construction of the theta family uses a criterion for good reduction in terms of theta constants.
A theta-level object over a $p$-adic field whose theta constants are integral and include a unit in
every coset of the subgroup $N'_n(\mathbf Z/N'^2)^{2g}$ extends, after a finite extension, to a
theta-level object over the ring of integers with the same theta constants.

% before: cv-toroidal-theta-family-good
The theta family is built for any class of complex theta-level objects that is good in the following
sense. Its members act on their torsion points through fixed integer matrices and satisfy the conditions
on $w$ and $w'$ of Lemma~\ref{cv-toroidal-theta-reduction}, and their theta constants fill out finitely
many holomorphic families with $\delta$ parameters.

% before: cv-toroidal-theta-family-pel-good
The PEL class of the datum is good once the level is a large power of $n$. The matrices record the action
of the generators of $\mathcal O_E$ on the symplectic basis, and the holomorphic families come from the
complex uniformization of the members by lattices with Hermitian forms.

% before: cv-toroidal-theta-family-exists
A good class has a theta family once the level is a large power of $n$. It is cut out by the polynomial
relations among the theta constants of the members of the class, and the criterion for good reduction of
Lemma~\ref{cv-toroidal-theta-reduction} is used in its construction.

% before: cv-toroidal-theta-family-pel-fibre-line
A PEL structure over a field is descended along a field extension by descending its line bundle and the isomorphisms
that define its structures. For smooth, proper and geometrically connected schemes such as an abelian variety and its
powers, triviality of a line bundle can be checked after the extension, because the global functions are the constants.

% before: cv-toroidal-theta-family-pel-fibre-descent
At a point where a theta-level object comes from a complex member of the PEL class, the PEL structure is known on the
complex member and must be moved to the field of definition. The action descends because the image of $a$ generates
$\mathcal O_E$, and the conditions at geometric points can be compared in a common algebraically closed field.

% before: cv-toroidal-theta-family-pel-fibre-pull
A PEL structure on a theta-level object over a field must then be carried to its base change. Pulling back the
$\mathcal O_E$-semi-abelian scheme gives the structure, and the conditions at geometric points are those of the original
object at the image points.

% before: cv-toroidal-theta-family-pel-fibre
A PEL structure on a theta-level object is spread from its fibres at the points where it comes from a complex member of
the PEL class. At such a point the fibre has a PEL structure over a field of characteristic zero, descended from the
complex member.

% before: cv-toroidal-theta-family-pel-action
The action of $\mathcal O_E$ on the abelian scheme is determined by the endomorphisms given by a generating set. They
satisfy the relations of $\mathcal O_E$ on a dense set of fibres, and that is enough over a reduced base.

% before: cv-toroidal-theta-family-pel-bundles-seesaw
Isomorphisms of line bundles on an abelian scheme are spread from a dense set of fibres by a seesaw argument. Over a
reduced base, triviality on a dense set of fibres gives triviality locally on the base, and triviality along a section
removes the line bundle pulled back from the base.

% before: cv-toroidal-theta-family-pel-bundles
The Rosati condition and the normalization of the line bundle are isomorphisms of line bundles. Over a reduced base
they spread from a dense set of fibres.

% before: cv-toroidal-theta-family-pel-cubical
An $\mathcal O_E$-semi-abelian scheme carries a symmetric cubical line bundle. Once the Rosati condition holds, the
normalized bundle $\Lambda_M(y,w_0y)$ has these structures and the same polarization as $M$.

% before: cv-toroidal-theta-family-pel-lie
The signature condition on the Lie algebra is a condition at the geometric points of characteristic zero. On a smooth
group scheme over a reduced base it spreads from a dense set of such points.

% before: cv-toroidal-theta-family-pel-level
The Tate systems of a PEL structure are built level by level. At each level, an isometry from the lattice to the torsion
points that holds over a dense set of points holds at every geometric point where the level is invertible.

% before: cv-toroidal-theta-family-pel-tate
Isometries at each finite level assemble into a Tate system. An isometry at one level is bijective onto the torsion
because the form on the lattice is perfect, and the finiteness of the torsion lets compatible choices be made over all
levels.

% before: cv-toroidal-theta-family-pel
The universal theta-level object of a theta family for the PEL class carries a PEL structure. The
structure is spread over the reduced, locally noetherian base from the dense set of points at which the
object comes from a complex member of the class.

% before: cv-toroidal-theta-family-trdeg-osgood
The lower bound for the dimension of a theta family comes from holomorphic families of theta constants, and its
proof uses the identity theorem for them. A complex differentiable function of several variables is analytic, so the
identity theorem applies to it.

% before: cv-toroidal-theta-family-trdeg-general
A polynomial in the values of a holomorphic family either vanishes on the whole family or vanishes on a closed set
with empty interior. Over a countable ring there are only countably many polynomials, so by the Baire category theorem
one point of the family sees all of them.

% before: cv-toroidal-theta-family-trdeg-rank
A holomorphic map whose fibres are countable cannot lose rank everywhere. Near a point of largest rank, a map with
non-injective derivative would be constant along a complex submanifold of positive dimension, an uncountable fibre.

% before: cv-toroidal-theta-family-trdeg-ratios
For a holomorphic family of theta constants whose proportional values are countable, the projective coordinates
depend on all the parameters. So some $d$ ratios of coordinates have an invertible derivative at a point, and no
nonzero polynomial relation holds among them on the family.

% before: cv-toroidal-theta-family-trdeg-cone
The polynomial relations that hold on a holomorphic family and on all its nonzero multiples form an ideal. Multiples
are allowed because the theta constants of a complex member of a class are defined only up to a nonzero factor.

% before: cv-toroidal-theta-family-trdeg-cone-ideal
This ideal is prime for a connected family, and $d$ ratios of the coordinates stay algebraically independent in any
field in which it is the ideal of relations. These are the two properties of it that the transcendence bound uses.

% before: cv-toroidal-theta-family-trdeg-prime
The theta ideal of a good class is the intersection of finitely many such ideals, one for each holomorphic family of
condition (3). So a prime that contains the theta ideal contains one of them, which carries $\delta$ independent ratios.

% before: cv-toroidal-theta-family-trdeg
The holomorphic families of theta constants of a good class bound the dimension of the theta family from
below. Their sets of proportional values are countable, so each component of the theta family has
transcendence degree at least $\delta$ over the base.

% before: cv-toroidal-theta-family-criterion
The smoothness of the generic fibre comes from two bounds: a lower bound for the dimensions of its
components and an upper bound for its tangent spaces at complex points. Over a field of characteristic
zero they make every local ring at a closed point regular of dimension $\delta$.

% before: cv-toroidal-theta-family-smooth
For a theta family of a good class, the lower bound is
Lemma~\ref{cv-toroidal-theta-family-trdeg}. What remains for the smoothness of the generic fibre is the
bound for the tangent spaces at complex points.

% before: cv-toroidal-theta-family-tangent-good
The tangent bound is computed on the complex uniformization of a fibre. It needs a uniformization whose differentials
describe the tangent vectors, through which endomorphisms lift to linear maps, and whose theta functions come from
sections.

% before: cv-toroidal-theta-family-tangent-ks
A derivation of the base deforms the complex fibre, and the deformation is recorded by a matrix acting antilinearly on
the universal cover. The hypothesis on points with values in $\mathbf C[\epsilon]$ makes this Kodaira--Spencer map
injective, and the polarization and the endomorphisms constrain its values.

% before: cv-toroidal-theta-family-tangent-fibre
The constraints from $\mathcal O_E$ need its action on the complex fibre, with the signature of its action on the Lie
algebra. Both come from the PEL structure through its isomorphism with the abelian scheme of the theta-level object.

% before: cv-toroidal-theta-family-tangent-lift
On the universal cover the action of $\mathcal O_E$ becomes linear. The eigenspaces of the linear action have the
dimensions of the eigenspaces in the Lie algebra, as the uniformization identifies the Lie algebra with the cover.

% before: cv-toroidal-theta-family-tangent-dimension
What remains is linear algebra. A symmetric antilinear map that commutes with $\mathcal O_E$ exchanges the eigenspaces
of conjugate embeddings, and positivity shows that it is determined by its values on the eigenspaces of the embeddings
in the CM type.

% before: cv-toroidal-theta-family-tangent
The tangent spaces of a family with a PEL structure are bounded through deformations. When distinct
tangent vectors carry distinct deformations, they inject into the symmetric, antilinear maps between the
eigenspaces of the Lie algebra that commute with $\mathcal O_E$, a space of dimension $\delta$.

% before: cv-toroidal-theta-family
For the class of complex points of the datum, the theta family exists once the level is a large power
of $n$. Its universal abelian scheme carries the structures of the datum, and its generic fibre is
smooth of the dimension of the moduli schemes.

% before: cv-toroidal-universal-cubical
A morphism of $\mathcal O_E$-semi-abelian schemes must respect the cubical structures of the line
bundles. Over a reduced base whose generic points lie in the abelian locus, this follows from an
isomorphism of the line bundles alone, because two cubical structures differ by a unit on the triple
product that is constant on the abelian fibres.

% before: cv-toroidal-universal-matched
A theta-level object and a full-level object can describe the same abelian scheme, with the theta
level points the images of a basis of the lattice under the full level structure. An isomorphism that
identifies the two descriptions lets the classifying properties of the theta family be transferred to
full-level objects.

% before: cv-toroidal-universal-structures
The PEL structure of a theta family carries a full level structure, and a full-level object over a
flat base is a theta-level object with a PEL structure. Morphisms of full-level objects and of the
matched theta-level objects correspond, once the actions of the generators of $\mathcal O_E$ determine
the whole action.

% before: cv-toroidal-universal-complex
To apply the classifying property of the theta family, the complex points of a theta-level object must
be dense. A theta-level object with a PEL structure over a field of characteristic zero descends to a
countable field, which embeds into $\mathbf C$, and its PEL structure descends with it.

% before: cv-toroidal-theta-universal
Over $R_{n,N}$, the PEL structure of the theta family gives its universal abelian scheme a full level
structure. With it, the theta family classifies the full-level objects over reduced, locally
noetherian, flat bases, so it is a full-level family.

% before: cv-toroidal-envelope
The compactification of the theta family is a normal projective scheme that contains a dense open part
of it, with a semi-abelian scheme extending the universal one there. It carries an ample line bundle
presented by bases of theta sections of the semi-abelian scheme. Because the line bundle is tied to
the semi-abelian scheme in this way, automorphisms of the theta family extend to the compactification.

% before: cv-toroidal-theta-envelope-chart
The compactification is built over the generic points of the theta family, one integral chart at a
time. A generic point of a reduced scheme of finite type and flat over a noetherian domain lies in an
affine open subscheme whose ring is a finitely generated domain.

% before: cv-toroidal-theta-envelope-homogeneous
Theta constants are defined up to a unit, so the relations among them are homogeneous. The theta ideal
contains the homogeneous components of its elements.

% before: cv-toroidal-theta-envelope-ratios
On an integral chart of a theta family, the theta constants determine the point. So the ratios of the
theta constants, with the constants of the base, generate the function field of the chart.

% before: cv-toroidal-theta-envelope-pull
Over a chart, the compactification is built from the restrictions of the theta-level object and of its
PEL structure. Pulling back a theta-level object and an $\mathcal O_E$-semi-abelian scheme identified with
it keeps the identification.

% before: cv-toroidal-theta-envelope-component
Over an integral chart whose function field is generated by ratios of theta constants, the
compactification is built from theta functions: its semi-abelian scheme is an open part of the closure of
the abelian scheme in a projective space of theta functions. The semistable reduction theorem and
Mumford's construction enter through bounds for the theta constants near the boundary.

% before: cv-toroidal-theta-envelope-glue-dense
A theta envelope only sees its base through a dense open subscheme. So an envelope of a dense open part of the base is
an envelope of the whole base, and an envelope of the base restricts to one of each dense open part.

% before: cv-toroidal-theta-envelope-glue-rank
Envelopes over different charts can be put side by side only if their presentations have the same number of sections.
At a point with values in a complete discrete valuation ring the sections form a basis of the sections of a power of a
principal polarization, so their number is fixed by Riemann--Roch; without such points it can be changed freely.

% before: cv-toroidal-theta-envelope-glue-triple
Two $\mathcal O_E$-semi-abelian schemes over the two pieces of an open cover, identified over the overlap, glue to one
over the whole scheme. Reducedness and abelian generic points make the identification respect the cubical structures.

% before: cv-toroidal-theta-envelope-glue-sum
Envelopes of two disjoint open parts of the base whose presentations have the same shape combine into one envelope.
Its scheme is the disjoint union of the two, and the semi-abelian schemes are glued in two steps so that the
semi-abelian scheme of the base maps to the result.

% before: cv-toroidal-theta-envelope-glue
The envelopes over the integral charts through the generic points are glued into one. The charts are
disjoint, and once their presentations share the same integers, their disjoint union is an envelope over
a dense open part of the base.

% before: cv-toroidal-theta-envelope
The theta envelope of a PEL theta family is glued from envelopes over integral charts through its
generic points. On each chart the ratios of the theta constants generate the function field, which is
what the construction from theta functions requires.

% before: cv-toroidal-abelian-pullback
The map from the abelian locus to the full-level family is built by extending a full level structure
from a dense open part, and that extension is checked after base change to points and to discrete
valuation rings. All the conditions defining a full-level object are compatible with base change, so a
full-level object restricts to one over any scheme mapping to its base.

% before: cv-toroidal-abelian-torsion
A full level structure consists of torsion sections, and multiplication by an invertible integer on a
proper group scheme is finite. Over a normal base, a lift through a finite morphism that is given on a
dense open subscheme extends to the whole base, so torsion sections extend as well.

% before: cv-toroidal-abelian-cover
Ampleness, the principal polarization and the conditions at geometric points of characteristic zero can
each be tested on morphisms whose images cover the base. For ampleness and principality this comes from
the fibres, through the openness of both conditions; for the geometric conditions, from extending the
algebraically closed field.

% before: cv-toroidal-abelian-dvr-closed
Over a discrete valuation ring, a principal polarization is extended from the generic fibre by studying the
group $K(\mathcal L)$ of translations that preserve the line bundle up to a line bundle from the base. This
group is represented by a closed subscheme of the abelian scheme, and it is finite when the line bundle is
relatively ample.

% before: cv-toroidal-abelian-dvr-flat-criterion
Flatness over a discrete valuation ring means that a uniformizer kills no nonzero element of a local ring. A
closed subscheme is flat once, near each point of its special fibre, the morphisms from local rings that factor
through it are cut out by an ideal modulo which the uniformizer is a nonzerodivisor.

% before: cv-toroidal-abelian-dvr-flat-stalk
The ideals are found in the local rings of a smooth scheme over the discrete valuation ring, at closed points of
the special fibre. Modulo the uniformizer these local rings are those of a smooth variety, which are regular of
the relative dimension.

% before: cv-toroidal-abelian-dvr-flat-complex-sections
Whether a line bundle in a family is trivial on a fibre is read off from its sections there. A complex of free
modules over the base that computes these sections after every base change, together with their functions in
fixed charts, turns this into linear algebra over the base.

% before: cv-toroidal-abelian-dvr-flat-complex-model
The \v{C}ech complex of a line bundle for an affine cover computes its cohomology after every base change, but its
terms are not finitely generated. The theorem on cohomology and base change replaces it by a complex of finite
projective modules with the same property.

% before: cv-toroidal-abelian-dvr-flat-complex-matrices
Finite projective modules are direct summands of free ones. Adding complements to the first three terms turns the
complex into matrices without changing the kernel in degree $0$ or the exactness in degree $1$ after any base change.

% before: cv-toroidal-abelian-dvr-flat-complex
Such a complex exists for every line bundle on a smooth projective scheme. Its matrices come from a finite projective
model of the \v{C}ech complex, and its charts from the cover that the projective embedding provides.

% before: cv-toroidal-abelian-dvr-flat-cyclic-minor
On the fibre where the line bundle is trivial, the kernel of $D_0$ is a line. A nonvanishing minor spreads this to
every base change on which one element of the base becomes a unit, with a fixed generator of the kernel.

% before: cv-toroidal-abelian-dvr-flat-cyclic-coker
A matrix whose kernel is spanned by one vector after every base change has a cyclic dual cokernel. Testing the
kernel on a square-zero extension identifies that cokernel with a quotient of the base.

% before: cv-toroidal-abelian-dvr-flat-cyclic
When the functions on a fibre are constant and the line bundle is trivial there, the sections over a local base
are governed by a single ideal. Its vanishing locus is the set of base changes on which the line bundle is
trivial.

% before: cv-toroidal-abelian-dvr-flat-free-model
Over a noetherian base, a bounded complex of flat modules with finitely generated cohomology has a finite projective
model. The model computes the same cohomology after every base change, in particular on every fibre.

% before: cv-toroidal-abelian-dvr-flat-free-yoneda
Two matrices whose kernels correspond after every base change have isomorphic dual cokernels. A square-zero
extension of the base by a module turns homomorphisms out of the dual cokernel into kernel vectors.

% before: cv-toroidal-abelian-dvr-flat-free
For a line bundle $\mathcal F$ on an abelian scheme with $\Lambda(\mathcal F)$ trivial, the complex can be replaced by a longer
complex of free modules whose fibres are exact away from the points where the line bundle is trivial. This
exactness is what the depth argument of the next lemma needs.

% before: cv-toroidal-abelian-dvr-flat-regular-split
The depth argument starts in the top degree, where the fibre at the closed point is exact. Over a local ring,
exactness of the reduction of a complex of free modules at one place splits the complex there.

% before: cv-toroidal-abelian-dvr-flat-regular-homotopy
Away from the closed point the fibres are exact in every degree. There the splitting holds after inverting an
element outside the prime, and it is built degree by degree from the top.

% before: cv-toroidal-abelian-dvr-flat-regular-zero
An element of a module that dies at every prime other than the maximal ideal is supported at the closed point. If
the module has no submodule isomorphic to the residue field, such an element is zero.

% before: cv-toroidal-abelian-dvr-flat-regular-step
The dual complex is treated one degree at a time, passing from degree $i+1$ to degree $i$. Each step trades one
degree of vanishing of $\operatorname{Ext}$ from the residue field for the injectivity of $\varpi$ on the next
cokernel.

% before: cv-toroidal-abelian-dvr-flat-regular
A complex of free modules whose fibres are exact except at the closed point, in degrees up to the dimension, has
a dual cokernel on which a regular parameter acts injectively. This is the commutative algebra behind the
flatness of $K(\mathcal L)$.

% before: cv-toroidal-abelian-dvr-flat-ideal
These lemmas combine at a point of an abelian scheme. The translates of $\mathcal L$ by the point form a line
bundle in a family over the local ring, and the morphisms that factor through the subscheme are those over
which this line bundle is trivial.

% before: cv-toroidal-abelian-dvr-flat
A finite subscheme over a discrete valuation ring is determined by its generic fibre once it is flat, since
its ring of functions then embeds in that of the generic fibre. The subscheme representing $K(\mathcal L)$
has this flatness.

% before: cv-toroidal-abelian-dvr-kernel
With flatness, the subscheme representing $K(\mathcal L)$ is the closure of its generic fibre. If the
polarization of the generic fibre is principal, that generic fibre is the zero section, and so is the whole
subscheme.

% before: cv-toroidal-abelian-dvr-specialize
For a principally polarized abelian scheme over an integral base, a line bundle whose restriction to the
generic fibre is ample is ample on the whole scheme. This is the input that lets ampleness pass along a
finite flat surjection from such an abelian scheme.

% before: cv-toroidal-abelian-dvr-isogenous
Ampleness of a line bundle on an abelian scheme can be tested after pulling it back along a finite surjective
morphism. If the source of that morphism is principally polarized, the generic fibre then controls ampleness
over a Dedekind base.

% before: cv-toroidal-abelian-dvr-cubical-prime
Extending a generator of a line bundle across a divisor comes down to valuations at the generic point of the
divisor. When the divisor is cut out by a function that is prime near its generic point, the local ring there
is a discrete valuation ring, and a section off the divisor extends after multiplication by a unique power of
that function.

% before: cv-toroidal-abelian-dvr-cubical-dvr
Over a discrete valuation ring with integral fibres, the special fibre is such a divisor, cut out by the
uniformizer. A section of the base then fixes the power of the uniformizer, so a generator on the generic fibre
that extends along the section extends everywhere.

% before: cv-toroidal-abelian-dvr-cubical-normal
Over a normal base the same extension holds, since a function on a normal scheme is determined by its values at
the points of codimension one. Those points lie over points of codimension one of the base, where the local ring
is a discrete valuation ring.

% before: cv-toroidal-abelian-dvr-cubical-iso
Applied to the quotient of two line bundles, the extension of generators shows that line bundles that agree on
the generic part and along a section agree everywhere. This is the form used for the Rosati condition and the
normalization.

% before: cv-toroidal-abelian-dvr-cubical-extension
A cubical structure and a symmetry are generators of line bundles on powers of the group. They extend from the
generic part, and the identities they satisfy continue to hold because the powers of the group are integral.

% before: cv-toroidal-abelian-dvr-cubical
To extend an $\mathcal O_E$-semi-abelian scheme from the generic fibre, the group, the action and a line
bundle are extended first. The remaining structure, a cubical structure, its symmetry, the Rosati condition
and the normalization, then extends over a normal base.

% before: cv-toroidal-abelian-dvr-ample-square
Over a valuation ring with algebraically closed fraction field, points of the closed fibre lift to sections of
the abelian scheme. Translating a section of a line bundle by a lifted point and its inverse then produces
sections of the square of the line bundle that do not vanish at any given point.

% before: cv-toroidal-abelian-dvr-ample-descent
The sections produced over such a valuation ring must come back to the discrete valuation ring. Sections are
compatible with flat base change, so sections that generate after the base change can be replaced by sections
before it.

% before: cv-toroidal-abelian-dvr-ample-generating
Together with the extension of sections from the generic fibre, these give sections of a power of the line
bundle that generate it on the whole abelian scheme. Some of them do not vanish at the origin of the special
fibre, which is what the ampleness criterion there requires.

% before: cv-toroidal-abelian-dvr-ample-effective-finite
A closed subgroup of an abelian variety that is infinite contains an abelian variety of positive dimension, which
has many points of every prime order. So the absence of such points bounds the subgroup.

% before: cv-toroidal-abelian-dvr-ample-effective-rigid
A morphism from an abelian variety that contracts a curve is invariant under translations along it, by the
rigidity lemma. If only finitely many translations preserve the morphism, its fibres are finite.

% before: cv-toroidal-abelian-dvr-ample-effective-quasifinite
For the morphism given by sections of a power of $\mathcal L$, a translation that preserves it also preserves that
power of $\mathcal L$. The theorem of the square and the hypothesis on torsion points then leave only finitely many
such translations.

% before: cv-toroidal-abelian-dvr-ample-effective-cover
The ampleness criterion starts from one nonzero section of a power of $\mathcal L$ on a geometric fibre, and the map
to projective space needs sections that generate a power everywhere. The translates of the section by a point and
by its inverse, combined by the theorem of the square, give a section of the doubled power that does not vanish at
any given point.

% before: cv-toroidal-abelian-dvr-ample-effective
On an abelian variety, a line bundle with a nonzero section of some power is ample once the translations that
preserve it have no nonzero torsion of a suitable prime order. This criterion is applied on a geometric special
fibre.

% before: cv-toroidal-abelian-dvr-ample-torsion-weil
A cubical structure gives, for two points killed by $N$, two generators of the same line bundle. Their ratio is a
unit on the base, and it is the pairing that carries the torsion condition between fibres.

% before: cv-toroidal-abelian-dvr-ample-torsion-basic
The pairing is computed from products of values of the cubical structure. The cocycle identity makes it
multiplicative, and so its values are $N$-th roots of unity, compatible with base change.

% before: cv-toroidal-abelian-dvr-ample-torsion-kernel
A point whose translate of the line bundle differs from it by a trivial line bundle pairs trivially with every
torsion point. A trivialization of the translate gives values that turn the two generators into each other.

% before: cv-toroidal-abelian-dvr-ample-torsion-descent
The converse needs descent along multiplication by $N$ on a geometric fibre. That morphism is finite and \'etale,
and functions invariant under the translations by $N$-torsion points descend along it.

% before: cv-toroidal-abelian-dvr-ample-torsion-nondegenerate
On a geometric fibre the pairing detects the triviality of translates. A point that pairs trivially with all
torsion points gives an invariant generator upstairs, which descends.

% before: cv-toroidal-abelian-dvr-ample-torsion
The torsion condition on the special fibre is checked on the generic fibre. Torsion points of order prime to the
residue characteristic lift over a valuation ring, and the Weil pairing defined by the cubical structure carries
the condition from the special fibre to the generic one.

% before: cv-toroidal-abelian-dvr-ample
The extended line bundle must also be relatively ample. Over a discrete valuation ring this follows from
ampleness and principality on the generic fibre, by passing to a geometric special fibre.

% before: cv-toroidal-abelian-dvr-structure
These pieces build, over a discrete valuation ring, an $\mathcal O_E$-semi-abelian structure with an ample
principal polarization on an abelian scheme whose generic fibre carries one. The action extends because
homomorphisms of abelian schemes extend from the generic fibre, and the line bundle because the abelian
scheme is regular.

% before: cv-toroidal-abelian-dvr-eigen
The conditions at geometric points spread from a dense set of points to all of them. For the Lie algebra
condition, the eigenspace dimensions of an action of $\mathcal O_E$ by matrices are determined by traces,
which are functions on the base.

% before: cv-toroidal-abelian-dvr-determinant
Locally on a reduced base, the Lie algebra of a smooth group scheme is free, and the action of
$\mathcal O_E$ on it is given by matrices. So the signature condition on the Lie algebra holds everywhere
once it holds on a dense set, and it fixes the relative dimension.

% before: cv-toroidal-abelian-dvr-pairing
The torsion condition involves the commutator pairing of the theta group. Between two torsion points of order
invertible on the base, this pairing is a root of unity, so the locus where it takes a given value is open and
closed.

% before: cv-toroidal-abelian-dvr-isometry
Level structures of a fixed level, together with a root of unity, are parametrized by a finite scheme over
the base. The conditions defining an isometry cut out a closed subset of it, so isometries at a dense set of
points give isometries at every geometric point.

% before: cv-toroidal-abelian-dvr-tate
Condition (b) asks for compatible isometries at all levels at once. Since there are finitely many isometries
at each level, compatible choices can be extracted from isometries at the levels $l\,j!$.

% before: cv-toroidal-abelian-dvr-spread
The conditions at geometric points pass along morphisms of $\mathcal O_E$-semi-abelian schemes. Over a
connected reduced base, they spread from a dense set of points to all geometric points of characteristic
zero: the Lie algebra condition through its criterion by traces, and the torsion condition through
isometries at each level.

% before: cv-toroidal-abelian-dvr
The covers used for the extension of the level structure are spectra of discrete valuation rings, whose
generic point lies in the dense open part where a full-level object is given. Over such a ring the
polarization and the conditions at geometric points extend from the generic fibre.

% before: cv-toroidal-abelian-level
With these pieces, a full level structure over a dense open part of a normal base extends to a
full-level object on an abelian scheme over the whole base. The torsion sections extend first, and the
remaining conditions are then checked on covers by the dense open part and by discrete valuation rings.

% before: cv-toroidal-abelian-valuative
The map from the abelian locus to the full-level family is shown to be proper by lifting along discrete
valuation rings. When the domain is quasi-projective, it suffices to lift the morphisms from fraction
fields of complete discrete valuation rings whose image is a generic point of the domain.

% before: cv-toroidal-abelian-proper
Over a discrete valuation ring, a semi-abelian scheme whose generic fibre is that of an abelian scheme is
itself an abelian scheme, because the isomorphism of generic fibres extends to a surjective homomorphism.
So the limits of points of the abelian locus that exist in the projective compactification stay in the
abelian locus.

% before: cv-toroidal-abelian-fibres
The ample line bundle of a theta presentation has transition functions built from the changes of basis
of theta sections. When the semi-abelian scheme is pulled back from the target, these changes of basis
are trivial along each fibre, so the ample line bundle is trivial on every fibre, and the fibres are
finite.

% before: cv-toroidal-abelian-finite
To apply the finiteness criterion to the map from the abelian locus, the ample theta presentation is
pulled back to the preimage of an affine open subset of the target. The inclusion of that preimage is an
affine morphism, so the pulled-back line bundle is still ample.

% before: cv-toroidal-abelian-birational
A finite map that is an isomorphism over a dense open set is an isomorphism near every point with an
integrally closed local ring, since the ring upstairs is then contained in a localization of the ring
downstairs. This gives the open set near the generic fibre over which the abelian locus and the
full-level family agree.

% before: cv-toroidal-abelian-locus
Over the abelian locus of the compactification, the semi-abelian scheme is an abelian scheme, and its
full level structure extends from the dense open part of the theta family. So the abelian locus maps
back to the theta family, finitely and surjectively, and isomorphically near the generic fibre.

% before: cv-toroidal-moduli-object
The moduli scheme at a level represents objects with a level structure at the integer of the level and a
full level structure. At each geometric point both are read off from one compatible system of
isomorphisms of the lattice with the torsion points.

% before: cv-toroidal-moduli-lift
Changing the compatible system of torsion isomorphisms requires similitudes of the lattice modulo all
integers at once. A similitude modulo one integer lifts to such a family, by Hensel's lemma at the primes
dividing the integer and the Chinese remainder theorem at the others.

% before: cv-toroidal-moduli-similitude
An endomorphism of the lattice that is injective modulo $\mathfrak M$ is, modulo $\mathfrak M$, a similitude
of a family that is the identity modulo a given integer prime to $M$. Twisting a compatible system by the
family keeps the full level structure and moves the level structure at $\mathfrak M$.

% before: cv-toroidal-moduli-tate
A level structure of type $\mathfrak M$ at a geometric point fits into a compatible system with a given
full level structure, after a twist by a family of similitudes. This is what turns a full-level object
with such a level structure into an object at the level.

% before: cv-toroidal-moduli-objects
Objects at the level and full-level objects over $K_l$ differ only by the level structure of type
$\mathfrak M$. At a geometric point it can always be added, and full-level objects pull back along
morphisms.

% before: cv-toroidal-moduli-torsion
Level structures of type $\mathfrak M$ are torsion points of order $M$, which form a finite \'etale scheme
when $M$ is invertible. The conditions on them therefore hold on open and closed subsets.

% before: cv-toroidal-moduli-closed
The level structure of the moduli scheme, pushed to a full-level object, is a point of the finite
\'etale scheme of level structures. The closure of the image carries an object at the level, so it maps
back to the moduli scheme, and the moduli scheme becomes a closed subscheme of a finite scheme.

% before: cv-toroidal-moduli-finite
The moduli scheme at a level of full level $N$ carries a full-level object, so it maps to the theta
family. The map is finite onto the generic fibre, because the remaining level structure, at the
integer of the level, is a point of a finite \'etale scheme over the theta family.

% before: cv-toroidal-level-change
A change of the full level structure by an endomorphism of the lattice, together with a compatible
automorphism of the moduli scheme, induces an automorphism of the theta family. It comes from the
classifying property of the full-level family, applied to the changed structure and to its inverse
modulo $N$.

% before: cv-toroidal-extend-lift
A finite morphism onto a normal scheme that has a section over a flat, dense part has a section
everywhere. This gives the morphism of compactifications from a finite cover of the first one.

% before: cv-toroidal-extend-trivial
A line bundle on a smooth family with integral fibres over a normal base is trivial once it is trivial
over a flat, dense part and along a section. The section rules out the fibres over points of codimension
one as divisors.

% before: cv-toroidal-extend-homomorphisms-torsion
The next five lemmas prepare the extension of homomorphisms of semi-abelian schemes. Over an algebraically
closed field the torsion of a semi-abelian variety is finite in each order. This bounds the points of the closure
of a graph over a closed point.

% before: cv-toroidal-extend-homomorphisms-points
A flat morphism that is a homeomorphism onto its image and surjective on stalks behaves like an open immersion
for factoring morphisms. Points outside its image are reached by discrete valuation rings whose generic points
come from it, which brings the case of a discrete valuation ring to bear.

% before: cv-toroidal-extend-homomorphisms-closure
A morphism defined on a dense part of a scheme is extended through the closure of its graph. These are the
properties of that closure used below: its points over a scheme with a dense part, and its relation to the part
it comes from.

% before: cv-toroidal-extend-homomorphisms-isomorphism
The projection from the closure of the graph is shown to be an isomorphism by a form of Zariski's main theorem.
The flat, dense part supplies the generic point and the birationality that the theorem needs.

% before: cv-toroidal-extend-homomorphisms-quasi-finite
The projection from the closure of the graph has finite fibres because torsion points of a given order are finite
in number. Over a scheme of finite type over the integers, closed points have finite residue fields, so each point
of the group has finite order there.

% before: cv-toroidal-extend-homomorphisms
Homomorphisms of semi-abelian schemes over a normal base are determined by, and extend from, their
restrictions to a flat, dense part. For a discrete valuation ring this is
Theorem~\ref{cv-model-toroidal-extension}.

% before: cv-toroidal-extend-identify
Two $\mathcal O_E$-semi-abelian schemes over a normal base that agree with the same proper one over a
flat, dense part are isomorphic. The isomorphism comes from the extension of homomorphisms, and the
line bundles agree because their quotient is trivial along the zero section and over the dense part.

% before: cv-toroidal-extend-bundles
An ample theta presentation determines its ample line bundle through the theta sections. Over a normal
base, two presentations of the same $\mathcal O_E$-semi-abelian scheme with the same integers give
isomorphic line bundles, because both families of sections are bases over every complete discrete
valuation ring through a generic point.

% before: cv-toroidal-extend-finite
On the normalization of the product of two compactifications, the pullbacks of the two ample line
bundles agree. So the second one is trivial on the fibres of the first projection, and an ample line
bundle can be trivial only on a finite proper scheme.

% before: cv-toroidal-extend
Two compactifications of the same abelian scheme with ample theta presentations of the same level and
integers are related by a morphism that extends the identity of the common open part. Applied to the
theta envelope and its twist by an automorphism of the theta family, it extends the automorphism to
the compactification.

% before: cv-model-toroidal-top
The statements above give the compactification at full level a power of $n$. The moduli scheme at each
level of this full level maps finitely onto its abelian part, and changes of the full level structure
extend to it.

% before: cv-model-toroidal-normalization
A compactification at one level gives one at a deeper level by normalization in the moduli scheme at
the deeper level. The semi-abelian scheme is pulled back, and so are the properties of being
theta-split.

% before: cv-model-toroidal-quotient
To remove the full level structure at $n$, the moduli scheme with that structure is divided by the
group that changes it. This group acts by automorphisms over the moduli scheme without the full level
structure, and each of them changes the full level structure by an endomorphism of the lattice.

% before: cv-model-toroidal-stabilizer
The quotient of a compactification by this group carries a semi-abelian scheme when the stabilizers of
points have invertible order and act trivially on the fibres. Serre's lemma on automorphisms that fix
the torsion points of an order at least $3$ gives the trivial action, and the full level structure then
shows that the stabilizers have order a power of $n$.

% before: cv-model-toroidal-pushforward
The theta sections descend to the quotient through the direct image of the line bundle of the
semi-abelian scheme, which has to be invertible. On the abelian part it is, by the principal
polarization, and the normality of the compactification extends this to all of it.

% before: cv-model-toroidal-descent
When the stabilizers have invertible order and act trivially on the fibres, the quotient of a
compactification by a finite group is a compactification of the quotient. The semi-abelian scheme,
its line bundle and the three properties of being theta-split all descend.

% before: cv-model-toroidal-base
These steps together give, for every level without a full level structure, a compactification over
the integers with $n$ and the level inverted. The model over the local ring at a prime $\ell$ is
obtained from it by normalization.

% before: cv-model-toroidal-degeneration
Over a complete normal local domain, a polarized semi-abelian scheme that is abelian at the generic
point is described by degeneration data. This description, with the action of $\mathcal O_E$ and the
torsion points, is property (2) of the models.

% before: cv-model-toroidal
Models exist at every prime $\ell$ at which the level has a prime-to-$\ell$ part of at least $3$.
Their boundary is described by degenerations of abelian varieties, and this description controls the
Hodge metric near the boundary and the Lie algebra on the special fibre.

% before: cv-model-literature
The models of Theorem~\ref{cv-model-toroidal} are close to the toroidal compactifications of the
literature, which are the usual source of such models. The comparison says which of their properties
those sources give and in which cases.

% before: cv-model-eigenbundles
The determinant line of the two eigenspaces is needed as a line bundle on the whole closure of the
component in the generic fibre of a model. Over a field of characteristic zero that contains every
conjugate of $E$, the action of $\mathcal O_E$ splits the Lie algebra of the semi-abelian scheme into
eigenbundles.

% before: cv-periods-function
At a complex embedding the flat section is first found as a holomorphic function, built from the
periods of the abelian varieties in the complex uniformization. Its absolute value, corrected by the
Hodge Gram determinants of the two bases, is constant, and this gives the archimedean norms.

% before: cv-boundary-situation
The holomorphic function of Lemma~\ref{cv-periods-function} becomes algebraic once it is compared
with algebraic sections on a compactification. The comparison takes place on the closure of the
component in the generic fibre of a model, at one complex embedding.

% before: cv-period-matrix
Near a boundary point, the periods of the abelian varieties form a matrix-valued function on a
product of half-planes and discs, which covers the complement of the boundary in a polydisc. Going
once around the $k$-th component of the boundary adds an integral symmetric matrix to it.
Subtracting the corresponding multiples of the logarithms of the coordinates leaves a function that
these translations do not change.

% before: cv-period-riemann-slices
The bounds on the imaginary part of a period matrix come from extending functions of its entries
across the coordinate hyperplanes of the polydisc. The extension theorem used for this asks for a
limit along each one-variable slice through a point with a single zero coordinate, and builds the
extension from Cauchy's formula on a torus.

% before: cv-period-punctured-order
On such a slice, the functions to be extended are exponentials of holomorphic functions which,
multiplied by a power of the slice coordinate, are bounded. A bounded function on a punctured disc
has a removable singularity, and a function with a logarithm on the punctured disc winds zero times
around the origin. Together these give the sign of the exponent and the limit of the exponential.

% before: cv-period-log-descent
A function on the product of half-planes and discs that the translations by $2\pi i$ in the first
$a$ coordinates do not change is a function on the punctured polydisc. It is holomorphic there,
since near each point a holomorphic branch of the logarithm inverts the exponential.

% before: cv-period-exp-monomial
For an integral vector $v$, the exponential of $2\pi i$ times the value at $v$ of the quadratic form
of a period matrix has absolute value at most one, because its imaginary part is positive definite.
The three preceding lemmas extend this exponential across the boundary, up to a monomial in the
boundary coordinates whose exponents are the values at $v$ of the quadratic forms of the matrices
$B_k$. The same bound forces these exponents to be nonnegative.

% before: cv-period-reduced-extends
The quadratic forms of the reduced matrix at integral vectors are logarithms of the extended
exponentials, and these logarithms exist on the polydisc because it is convex. Polarization
recovers the entries of the reduced matrix, which therefore extends holomorphically across the
boundary. Since rational vectors are dense, the nonnegative exponents make each $B_k$ positive
semidefinite.

% before: cv-period-kernel-positive
On the common kernel of the matrices $B_k$, the reduced matrix and the period matrix have the same
quadratic form, so the imaginary part of the extended reduced matrix is positive there away from the
boundary. The open mapping theorem, applied along a line through a boundary point, carries this
positivity to the boundary.

% before: cv-period-matrix-bounds
The imaginary part of a period matrix is that of the extended reduced matrix plus the matrices $B_k$
weighted by $|\log|z_k||/2\pi$. Near the origin these weights are large, so the $B_k$ give
positivity away from their common kernel, and the previous lemma gives it on that kernel. The
weights grow at most like $\lambda_a$, which gives the upper bound.

% before: cv-boundary-hodge-growth
Near a boundary point of a resolution of the compactification, the Hodge metric degenerates. Its Gram
determinants in algebraic bases grow at most like a power of the logarithm of the distance to the
boundary, a growth that a removable singularity allows.

% before: cv-boundary-unit
Read in algebraic bases of the semi-abelian scheme, the function built from the periods has its
absolute value controlled by the two Gram determinants. Their logarithmic bounds let it extend
across the boundary without zeros.

% before: cv-boundary-gaga
A function that is locally a holomorphic unit times a local generator of a line bundle defines a
holomorphic section without zeros on a proper smooth variety. By GAGA it is algebraic, and since the
field of definition is algebraically closed in the function field, it is defined over that field up
to a constant.

% before: cv-boundary-section
The three preceding lemmas combine on a resolution of the closure of the component. They give a
generator of the determinant line over the whole closure whose values at complex points are a
constant multiple of the function built from the periods.

% before: cv-function-generator
A generator of the determinant line gives a frame function, determined up to a constant of the
coefficient field. Its comparison with the function built from the periods holds at every complex
embedding at which the uniformization is defined, which gives the archimedean norms, and frame
functions built on models at different primes agree up to such a constant.

% before: cv-norms-zariski
At a finite place the norm of the section is constant once the generator has a constant order along
the special fibre. On a normal proper scheme with connected special fibre, a line bundle of degree
zero on every curve of the special fibre has this property.

% before: cv-hodge-lattice-germ
Alpöge builds the lattice of Lemma~\ref{cv-norms-hodge-lattice} from power series. Formal coordinates
of a semi-abelian scheme along its unit section give a formal group law with an action of
$\mathcal O_E$, and trivializations of the biextension defined by its cubical line bundle give power
series satisfying the identities of a symmetric biextension of the formal group by the formal
multiplicative group, in the sense of Mumford \cite{Mum69}. The lattice is read off from the linear
and bilinear coefficients of these series.

% before: cv-hodge-lattice-transport
Two choices of coordinates and trivializations give isomorphic formal germs, and the lattices they
define are compared through the transport maps. The terms of degree one and of bidegree $(1,1)$ in
the identities of a formal germ give the linear algebra that this comparison needs. In
characteristic $\ell$ the conjugate maps relate the Hodge module to the Frobenius twists of the Lie
algebra and of its dual.

% before: cv-germs-chart
The formal germs come from the semi-abelian scheme through coordinates and trivializations. Coordinates
along the unit section turn the group law and the action of $\mathcal O_E$ into power series. Near the
unit section the line bundle is trivial, and in a trivialization of $\Lambda$ the cubical structure, its
symmetry and the Rosati condition become the series $\Phi_1$, $\Phi_2$, $S$ and $P_a$.

% before: cv-germs-formal
Coordinates along the unit section exist wherever the tangent vectors at the unit have a basis. They
are universal for points that are the unit modulo a nilpotent ideal, and this universality makes the
power series and the changes of coordinates unique.

% before: cv-germs-local
Formal charts exist locally on the base. The identities of a formal germ are those of the cubical
structure written in coordinates, together with identities for the generators $w_a$, which hold because
every function on a fibre power of the group comes from the base.

% before: cv-germs-compare
Two formal charts at one point differ by a change of coordinates and a change of trivialization, which
together form an isomorphism of formal germs. The transport maps of these comparisons compose, so the
lattices defined from charts at different points can be compared over the generic point.

% before: cv-germs-unique-lambda
Only the changes of coordinates $\phi$ and $\psi$ are stated to be unique in a comparison. The change of
trivialization $\lambda$ is determined as well, by the same ratio of generators that gives the composition
rule for comparisons.

% before: cv-germs-principal
At a point of residue characteristic $\ell$ the fibre of the semi-abelian scheme is in general not an
abelian variety. Over an algebraically closed field it is an extension of an abelian variety by a torus,
and its line bundle comes from a principal polarization of the abelian variety. The arguments in
characteristic $\ell$ below use this property of the fibres and a version of it for tangent vectors.

% before: cv-germs-principal-base-change
Having a principal fibre depends only on the fibre as a group scheme with a line bundle. So it is kept
under extensions of the field and under pullback along morphisms of semi-abelian schemes.

% before: cv-germs-dvr
A point of the special fibre of the model is reached from the generic fibre through a complete discrete
valuation ring that dominates its local ring. The residue field is enlarged at the same time, so that the
torus over the closed point becomes split.

% before: cv-germs-raynaud
Over such a ring the pullback of the semi-abelian scheme comes from Mumford's construction. Its closed
fibre is then that of the Raynaud extension, an extension of an abelian variety by a split torus, and the
polarization of the abelian part is principal.

% before: cv-germs-principal-fibres
The two properties of the model in Theorem~\ref{cv-model-toroidal} combine with these results at every
point of positive residue characteristic. The first gives a field over which the torus splits, and the
second gives Mumford's construction over a complete discrete valuation ring with that residue field.

% before: cv-germs-torus-lift
A tangent vector of a split torus lifts to a point over the truncated power series ring of order $p$. If
the lift is killed by $p$, the tangent vector vanishes, since the $p$-th power of an element of that ring
only depends on its terms of order at most one.

% before: cv-germs-frobenius-series
In characteristic $p$, multiplication by $p$ in a formal group is a power series in the $p$-th powers of
the variables. Its linear part in these powers is the matrix $W$ in the first component of the conjugate
map.

% before: cv-germs-frobenius-points
Through a formal chart, multiplication by $p$ on points over truncated power series rings is computed by
the matrix $W$. Hence $p$ times such a point depends only on its first-order part, and a tangent vector
gives a point killed by $p$ when $W$ annihilates the $p$-th powers of its coordinates.

% before: cv-germs-descent
For an extension of an abelian variety by a torus, a generator of a translate of a pulled-back line bundle
along a tangent vector can be made invariant under the torus. The obstruction is a cocycle with values in
a square-zero ideal, which is a coboundary, and an invariant generator descends to the abelian variety.

% before: cv-germs-tangent
At a principal fibre, a tangent vector along which the translate of the line bundle is trivial, and which
lifts to a point killed by $p$, vanishes. Its image in the abelian variety vanishes because $K$ is trivial
there, and the torus part vanishes by the lifting lemma for tori.

% before: cv-germs-conj-injective
Over an algebraically closed field every vector is a vector of $p$-th powers, so a vector in the kernel of
the conjugate map gives a point of the formal chart killed by $p$. If the line bundle is trivial along its
first-order part, the tangent property forces the vector to vanish.

% before: cv-germs-kernel-bound
The freeness of the Hodge module at a principal point comes from bounds on the kernels of its
endomorphisms. Linear algebra compares these kernels with those for the abelian quotient and for the
lattice of characters of the torus, after a change of the bilinear coefficients.

% before: cv-germs-kernel-degree
For the abelian quotient of a principal fibre, the kernels of the Hodge endomorphisms are bounded
through the kernels of the elements of $\mathcal O_E$ acting on the abelian variety. These are finite
group schemes whose orders are powers of the norm, with an exponent fixed by the dimension.

% before: cv-germs-kernel-pairing
The kernel of a Hodge endomorphism is bounded by two spaces of tangent vectors: the kernel of the
conjugate action, and the vectors in the kernel of the action that pair to zero with it under the
bilinear coefficients. These two spaces are what the Hopf algebra of the kernel of the element sees.

% before: cv-germs-lattice-kernel
Part of the tangent space at a principal point comes from a lattice with an action of $\mathcal O_E$.
Over a field of characteristic $p$, the kernel of an element of $\mathcal O_E$ on such a lattice has
dimension at most the rank times the $p$-adic valuation of the norm. This is the bound that the
freeness criterion below asks of the Hodge module.

% before: cv-germs-free-local
The kernel bounds give freeness one maximal ideal of $k\otimes\mathcal O_E$ at a time. An element $t$
with distinct values at the characters above one prime of $\mathcal O_E$ separates generalized
eigenspaces, and an element $\pi$ with $\pi t$ nilpotent counts generators on each of them through its
kernel.

% before: cv-germs-free-primes
The elements $t$ and $\pi$ of the local statement come from a prime of $\mathcal O_E$ above $p$. The
element $\pi$ lies in the prime, with norm of $p$-adic valuation the residue degree, and $t$ lifts a
generator of the multiplicative group of the residue field and multiplies the prime into a nilpotent
ideal modulo $p$.

% before: cv-germs-free-criterion
With these elements the local statement holds at every maximal ideal of $k\otimes\mathcal O_E$, and the
Chinese remainder theorem and Nakayama's lemma glue the local generators. The result is a criterion for
freeness over $k\otimes\mathcal O_E$ in terms of characteristic polynomials and dimensions of kernels,
which the Hodge module at a principal point satisfies.

% before: cv-germs-hopf
The kernel bound needs a power of $p$ that divides the order of the kernel of $a$ on the abelian
quotient. That kernel is a finite commutative group scheme, so its coordinate ring is a finite
commutative and cocommutative Hopf algebra. Its primitive elements and its derivations at the counit
count the factors of $p$.

% before: cv-germs-hopf-dual
Cartier duality exchanges the two kinds of elements: the derivations at the counit of a Hopf algebra are
the primitive elements of its dual. The points of the dual are the grouplike elements, which are units.

% before: cv-germs-hopf-frobenius
In characteristic $p$ the $p$-th power of a primitive element is primitive. On a space stable under
$p$-th powers, an induction on the dimension of the ideal an element generates finds an element whose
$p$-th power is $0$ or itself.

% before: cv-germs-hopf-unit
Points of the dual Hopf algebra are tested by evaluating its elements. A subspace that no point
annihilates contains a unit, since a vector space over an infinite field is not a finite union of proper
subspaces.

% before: cv-germs-hopf-powers
The first $p$ powers of a nonzero primitive element are independent. Convolution powers of a linear form
that separates the element from $1$ pick out each power in turn.

% before: cv-germs-hopf-derivation
Multiplication by a primitive element dualizes to a derivation of the dual Hopf algebra, whose $p$-th power
is $0$ or the derivation itself. In both cases the algebra is free of rank $p$ over the kernel of the
derivation, and a derivation at a point is determined by its values on that kernel and at one element.

% before: cv-germs-hopf-quotient
A primitive element whose $p$-th power is $0$ or itself cuts the Hopf algebra down by a factor of $p$.
The quotient is again a Hopf algebra, and the only primitive elements lost in it are the multiples of
that element.

% before: cv-germs-hopf-count
Passing to these quotients one primitive element at a time, and to the Cartier dual for the derivations,
gives the power of $p$. The two spaces must be orthogonal: for $k[f]/(f^p)$ with $f$ primitive, the
primitive element $f$ and the derivation $f\mapsto1$ span spaces of dimension one, and the dimension of the
algebra is $p$.

% before: cv-germs-weil-flat
The descent in the pairing below needs the endomorphism $a$ of an abelian variety to be flat when its
kernel is finite. A finite endomorphism of an integral group scheme fixes the generic point, so it is flat
on a nonempty open set, and translations carry this flatness to every closed point.

% before: cv-germs-weil-isogeny
A finite kernel makes the fibres of the endomorphism finite, as each is a translate of the kernel. On a
proper group scheme the endomorphism is then finite, and on a smooth connected one the preceding lemma
applies.

% before: cv-germs-weil
The primitive elements and the derivations of the kernel come from the formal chart of the abelian
quotient. The generator $w_a$ of the chart restricts to a pairing between the kernels of $a$ and of
$\bar a$, which is bilinear by rigidity and nondegenerate by descent, as $K(\mathcal L)$ is trivial. On
first-order points of the chart it is the series $P_a$.

% before: cv-germs-kernel-hopf
The pairing of a first-order point of the kernel of $\bar a$ with the universal point of the kernel of $a$
gives primitive elements of its Hopf algebra, and first-order points of the kernel of $a$ give
derivations. Their values on each other are the bilinear terms $\beta_a$ of the formal germ.

% before: cv-germs-abelian-extension
A principal fibre is an extension of an abelian variety by a split torus, and the lemmas below use only
this extension. The line bundles $\Lambda_N$ measure how far the pullback of a line bundle is from being
additive in the point.

% before: cv-germs-abelian-units
Line bundles on the fibre are compared with line bundles on its abelian quotient through the action of the
torus, and the comparison produces invertible functions on products with the torus. On a reduced connected
base such a function is a character of the torus times a function on the base.

% before: cv-germs-abelian-seesaw
A line bundle on a product with the abelian quotient is trivial as soon as its pullback to the fibre is
trivial and one of its slices is. The pullback carries a generator; the torus acts on it by a character,
which a unit on the fibre absorbs, and the invariant generator descends.

% before: cv-germs-abelian-descent
The action of $\mathcal O_E$ and the polarization live on the fibre, and the quotient map is faithfully flat
with the torus as kernel. Homomorphisms that kill the torus factor through the quotient, and every
homomorphism to an abelian variety kills it, as a torus has no nonconstant maps to an abelian variety.

% before: cv-germs-abelian-rosati
The line bundle $M$ of a principal fibre comes with no action of $\mathcal O_E$ attached, but the Rosati
condition descends to it. Two applications of the preceding seesaw lemma reduce the condition on the
quotient to the condition on the fibre.

% before: cv-germs-abelian-normal
From $M$ one builds a line bundle with the same Mumford bundle $\Lambda$ that is symmetric, cubical and
normalized by $w_0$. As $w_0+\bar w_0=1$, the Rosati condition makes $\Lambda$ of the new bundle equal to
that of $M$, by the theorem of the cube.

% before: cv-germs-abelian-ktriv
The triviality of $K$ depends only on the Mumford bundle $\Lambda$. So it passes from $M$ to the normalized
line bundle.

% before: cv-germs-abelian-tangent
The tangent space of the fibre is an extension of that of the abelian quotient by that of the torus. The
quotient map is smooth, as its base change along itself is a product with the torus, so tangent vectors
lift.

% before: cv-germs-abelian-characters
The action of $\mathcal O_E$ on the fibre preserves the torus. Endomorphisms of a split torus are
endomorphisms of its character group, so $\mathcal O_E$ acts on the characters, and dually on a lattice
whose tensor product with $k$ is the tangent space of the torus.

% before: cv-germs-fibre-taylor
The points of a formal chart identify the completed local ring of the fibre at the unit with a ring of
power series. A function on the fibre is therefore determined by its Taylor series along the chart.

% before: cv-germs-abelian-bilinear
The bilinear terms of the formal germs are compared through products and substitutions of power series
that are $1$ on both axes. Only the terms of bidegree $(1,1)$ contribute, which gives the formulas below.

% before: cv-germs-abelian-quotient
At a principal fibre the semi-abelian scheme is an extension of an abelian variety by a split torus. The
abelian quotient inherits the action of $\mathcal O_E$ and a principal line bundle, the torus contributes
a lattice to the tangent space, and the bilinear terms of the two formal germs agree up to a change of
coordinates.

% before: cv-germs-cocycle-powers
The proof of Lemma~\ref{cv-germs-cocycle} compares the $p$-th power of a line bundle $\Lambda(y_1,y_2)$
with $\Lambda(py_1,y_2)$ along a formal chart. There the comparison is the product $\Pi_p$ of the series
$\Phi_1([i]x,x,y)$ for $i<p$. Its behaviour under addition in each argument follows from the identities of a
formal germ by induction.

% before: cv-germs-cocycle-series
Multiplication by $p$ on a formal germ in characteristic $p$ is a series in the $p$-th powers of the
variables. The product $\Pi_p$ has the same property in its first argument, which defines the series
$\varphi$. The bilinear terms of $\varphi$ are the matrix $D$ of the conjugate map.

% before: cv-germs-cocycle-kernel
For $c$ in the kernel of the conjugate map, differentiating $\varphi$ in the direction $c$ along $u=0$ gives a
series $\delta_c$ in one family of variables. The identities of the formal germ make $\delta_c$ additive up to
a series in $p$-th powers, and its linear terms are $D^{\mathsf t}c=0$. So all its partial derivatives vanish.

% before: cv-germs-cocycle-nfold
The cubical structure identifies $\Lambda(ny_1,y_2)$ with the $n$-th power of $\Lambda(y_1,y_2)$, through a
product of values of $\theta$. In a formal chart this product is $\Pi_n$, up to the generators $e_m$.

% before: cv-germs-cocycle-generator
For a point $y$ with values in $k_p$ that is killed by $p$, the line bundle $\Lambda(py,h)$ is trivial. The
generator $\kappa_p$ then gives a generator of the $p$-th power of $\mathcal L_y$, whose expansion along a
formal chart is computed from $\varphi$ at the point $y_w$, up to a $p$-th power.

% before: cv-germs-cocycle-first-order
A generator of the $p$-th power of a line bundle on $Y_p$ need not come from a generator of the line bundle
on $Y_1$. Over the open sets where $\mathcal L$ is trivial, the coefficients of $s^p$ of such a generator
differ by $p$-th powers on overlaps. On a normal fibre, if one of them is a $p$-th power in the function
field, so are all of them, and their $p$-th roots give the generator on $Y_1$.

% before: cv-germs-cocycle
Lemma~\ref{cv-germs-conj-injective} assumes that the line bundle $\mathcal L_{y_{w,1}}$ has a generator
whenever $\alpha_B(w_1^p,\dots,w_n^p)=0$. This hypothesis holds for every formal chart, with no condition
on the fibre.

% before: cv-hodge-lattice-model-germs
The model of Theorem~\ref{cv-model-toroidal} carries formal germs at all points of the closure of the
component, each compared with one at the generic point. Their Hodge modules give the stalks of the
lattice. The freeness of these modules over $\mathcal O_E$ and the injectivity of the conjugate maps
in characteristic $\ell$ are what the degree computation uses.

% before: cv-hodge-lattice-curve-degree
On a curve in characteristic $\ell$ the Hodge module is filtered in two ways: by the dual of the Lie
algebra, and by the Frobenius pullback of the Lie algebra. For a piece cut out by an idempotent of
$k\otimes\mathcal O_E$, with $k$ a finite field, the two filtrations show that the Frobenius of $k$
multiplies its degree by $\ell$. As that Frobenius has finite order, the degree vanishes.

% before: cv-norms-hodge-lattice
When $\ell$ divides the discriminant of $E$, the Lie algebra need not split into eigenbundles over the
special fibre. The determinant line is then replaced by the determinant of an eigenspace of a lattice
that extends the Lie algebra by its dual, and this determinant has degree zero on the curves in the
special fibre.

% before: cv-norms-padic
The finite norms of a flat section are read at integral points of the component, in integral bases of
the two eigenspaces. The constant order of Lemma~\ref{cv-norms-zariski}, applied to the lattice of
Lemma~\ref{cv-norms-hodge-lattice}, gives the same norm at all of them.

% before: cv-flat-degree
To bound slopes of automorphic forms of a given weight one shifts the weight by sections
whose norms are constant on each component. The degree of such a flat section is a
period of the CM point it passes through, and Theorem~\ref{cv-chowla-selberg} computes it.
Resolution of singularities and GAGA are needed to produce the section and its norms at
all places.

% before: cv-where-chowla-selberg-enters
The two forms of the Chowla--Selberg formula have different consumers. The difference
matters for reading the overview, which attributes all of these inputs to one upper bound.

% before: cv-functorial-resolution
The arithmetic intersection theory on toroidal compactifications and the construction of
flat sections both blow up singular schemes over number fields and their completions.
The constructions must be compatible with smooth morphisms and with extension of the base
field, which is the functorial form of Hironaka's theorem.

% before: cv-moriwaki-bigness
Arithmetic volumes count small sections, and a lower bound for them is what produces
lower bounds for slopes of automorphic forms. For nef line bundles the arithmetic
self-intersection number is such a lower bound, and it can be computed.

% before: cv-reconstruction-functionals
The arithmetic intersection numbers of a tower of flag bundles are linear forms on polynomials in
the tautological classes, and the proof of the next theorem compares them with linear forms that
live on the flag variety alone. Those forms, and the weights at which they are evaluated, are
defined first.

% before: cv-reconstruction-flag-degree
Every comparison below ends in an evaluation of the degree on the flag variety and of its two
derivatives $\mathrm{Sh}_\tau$ and $\mathrm{Tot}_\tau$ at powers of linear forms and at multiples of
power sums. Each such evaluation is a closed formula.

% before: cv-reconstruction-polarization
A linear form on polynomials of fixed degree is determined by its values on powers of linear forms.
Evaluating it on these powers gives a polynomial in the coefficients, and the coefficients of that
polynomial bound the values on monomials.

% before: cv-wall-lipschitz
The bounds on walls perturb a weight slightly and must control how the intersection numbers move.
When the values on monomials are bounded, the dependence on the weight is Lipschitz.

% before: cv-wall-kleiman
The degree of the perturbed tautological bundle on the generic fibre must be bounded above and below
by degrees of nef bundles. Intersection numbers of nef bundles are nonnegative.

% before: cv-wall-successive-minima
On a wall the arithmetic self-intersection number of the perturbed bundle is squeezed between
bounds for the heights of points. The two inequalities between heights of points and the height of
the variety are those of Zhang's theorem on successive minima.

% before: cv-wall-perturbed-weights
A weight on a wall is moved off the wall by a distance comparable to $1/B$, so that the perturbed
bundle becomes nef on the fibres at level $B$. The moved weight must keep every other property that
the height bounds require.

% before: cv-wall-rational-weights
The height bounds apply to bundles with rational coefficients only. Since the intersection numbers
depend continuously on the weight, a bound at rational weights on slightly thicker walls gives the
bound at every weight.

% before: cv-reconstruction-antisymmetric
Up to a multiple of the flag height, the arithmetic intersection numbers change sign under
exchanges of entries inside one block, and they are small near the walls between the blocks. A
polynomial of this kind is close to a multiple of the Vandermonde product by a block invariant
linear form.

% before: cv-reconstruction-coefficients
A multiple of the Vandermonde product by a block invariant linear form evaluates, on powers of
linear forms, to combinations of $\mathrm{Sh}_\tau$ and $\mathrm{Tot}_\tau$. This translates the
previous lemma into estimates for the intersection numbers of monomials.

% before: cv-reconstruction-flag-height
The flag height enters the comparison in three ways: it vanishes at weights with two equal
neighbouring entries, it is nonnegative at dominant weights, and it does not see the first power
sums. It is also invariant under relabelling the CM types, which allows averaging over them.

% before: cv-transport-power-sums
The relation between two levels of the tower holds for products of power sums with other
polynomials. Every polynomial of degree one more than the dimension of the flag variety is a sum
of such products.

% before: cv-transport-local-terms
The local term at a prime in the relation between two levels splits into a part proportional to
$\mathrm{Sh}$, a part proportional to $\mathrm{Tot}$ and a remainder of size $\log p/p^2$. The first two
parts change the coefficients of the comparison, and the remainders converge.

% before: cv-transport-test-polynomial
The lower bound at the base of the tower is a bound for the intersection number of one polynomial.
It is built from the Plücker class of the Grassmannian of the type and the classes of the fibres
over it, so that its flag height is a height of the Grassmannian.

% before: cv-transport-test-facts
The values of $\mathrm{Sh}$ and $\mathrm{Tot}$ at the test polynomial decide which of the two
coefficients of the comparison the lower bound at the base controls. They depend on which block of
the type is smaller.

% before: cv-transport-remainders
The remainders from the primes above a level are small compared with the main terms once the level
is large. The bounds below are the forms used for the weights of the theorem and for the test
polynomial.

% before: cv-transport-grassmannian
The flag height of the test polynomial is the height of a product of Grassmannians in their
Plücker embeddings. Points spanned by matrices of roots of unity have small height and are Zariski
dense, which bounds it from above.

% before: cv-condition-flag-model
The models of the flag bundle in the theorem below are arithmetic varieties in products of
projective spaces carrying a family of metrized line bundles, one for each tautological class, whose
metrics depend on a regularization parameter. Their arithmetic intersection numbers are limits as
the parameter tends to $0$, and their geometric degrees are integrals over the complex points.

% before: cv-proportionality-analytic-system
The complex points of a model are compared with a fixed compact complex manifold, the toroidal
compactification of the flag bundle over one component, on which the tautological bundles carry
invariant metrics with logarithmic singularities along the boundary. The regularized metrics, and
the bounds they satisfy, are what make the arithmetic intersection numbers converge.

% before: cv-period-monge-ampere
The Chern numbers of an analytic system are integrals of Monge--Amp\`ere measures of potentials
that are singular along the boundary. Off the boundary these measures have densities computed at
each point from the second derivatives of the potentials in a chart, and the integral over the
manifold is assembled from the charts with a partition of unity.

% before: cv-period-flag-domain
The potentials on the flag domain are logarithms of Gram determinants for an indefinite Hermitian
form. They are invariant under the unitary group of the form acting on the left, and they change by
pluriharmonic terms under the upper triangular group acting on the right. This moves the
computation at any point of the flag domain to a neighbourhood of the identity.

% before: cv-period-big-cell
Near an invertible upper triangular matrix, every matrix is the exponential of a strictly lower
triangular matrix times an upper triangular one, with both factors depending holomorphically on
it. The lower triangular factor gives coordinates transverse to the orbits of the upper triangular
group.

% before: cv-period-cell-coordinates
Near a point, a holomorphic map to the flag domain can be replaced by the exponential of a strictly
lower triangular matrix, at the cost of pluriharmonic terms in the potentials. The derivative of
this coordinate vanishes on the directions in which the map moves along an orbit of the upper
triangular group, and only on those.

% before: cv-period-cell-levi
At the identity the Levi forms of the potentials are diagonal in the entries below the diagonal.
After a relabelling that puts the positive part of the form first, an entry contributes with weight
one half to the potentials whose index lies between the labels of its row and its column.

% before: cv-period-chart
The uniformization of an analytic system gives, near each point off the boundary, holomorphic maps
to the flag domain that carry the potentials of the tautological bundles. The pointwise computation
uses one such neighbourhood at a time, and the positivity of the volume uses a point where the maps
have no direction along the upper triangular orbits.

% before: cv-period-levi-diagonal
Pulling the computation at the identity back along the cell coordinates gives the Levi forms of all
the potentials at a point in one set of linear coordinates. At a flag immersion these coordinates
are independent, and this is what makes the Monge--Amp\`ere density nonzero there.

% before: cv-period-density-diagonal
When a Levi form is a weighted sum of squares of linear forms, the determinant of its matrix is the
product of the weights times the squared absolute value of the determinant of the linear forms. This
gives the Monge--Amp\`ere density without computing the complex Hessian entry by entry.

% before: cv-period-pointwise
For a combination of the potentials with real coefficients, the weights of the diagonal Levi form
are sums of consecutive coefficients, one for each pair of indices. Their product becomes the
product of differences of weights in the theorem below, and the remaining factor does not depend on
the coefficients.

% before: cv-proportionality-mumford
When the manifold is locally uniformized by the flag domain, its Chern numbers are those of the
compact dual flag variety up to one positive factor, the volume. This is the source of the
proportionality in the theorem below.

% before: cv-proportionality-literature
The identity just proved has the shape of a classical theorem on locally symmetric varieties. The
comparison says what that theorem covers, since the analytic systems of this chapter are given by
charts and growth bounds.

% before: cv-proportionality-realization
A realization transports geometric degrees from the model to the manifold. Since the manifold does
not depend on the level, neither do the geometric degrees.

% before: cv-tail-shells
The rate of convergence of the arithmetic intersection numbers comes from bounds on the
Monge--Amp\`ere measures of the regularized metrics near the boundary, where the singular potentials
grow like $\log r$. A bound on the mass of the region where $r\ge R$ gives, by summing over dyadic
shells, a bound on the integral of $\log r$ over that region.

% before: cv-tail-cutoff
A potential bounded below only by a multiple of $-\log r$ is cut off at a constant level by a smooth
maximum, which keeps the curvature semipositive. The cut-off potential is bounded, and where the
original one lies above the level the two differ by a constant, so their Monge--Amp\`ere measures
agree there.

% before: cv-tail-situation
The tail bounds concern one compact complex manifold with a boundary divisor of simple normal
crossings, cut out by a section of norm at most $1$. The function $r$ tends to infinity at the
boundary like $-2\log\|\sigma\|$, and an analytic system carries the same data.

% before: cv-tail-poincare-lelong
Integrals of $\log\|\sigma\|$ against products of first Chern forms are moved off the boundary by the
Poincar\'e--Lelong formula. It turns an integral against $dd^cv$ into integrals of $v$ over the
components of the boundary divisor and against the first Chern form of the bundle of $\sigma$.

% before: cv-tail-energy
For potentials bounded between $-M-1$ and $0$, the integral of $r$ against their Monge--Amp\`ere
measure grows at most linearly in $M$. Together with the cutoff at a level of size $\log R$, this
bounds the mass of the region where $r\ge R$.

% before: cv-tail-bound
For potentials bounded below by $-C_1\log r$, the cutoff at the level $C_1\log R+2$ leaves the
measure unchanged where $r<R$, and the energy bound controls the mass where $r\ge R$. The resulting
constants depend on the potentials only through $C_1$.

% before: cv-tail-uniform
On an analytic system, the regularized metrics at small parameters are compared with the metric at
one fixed parameter, whose first Chern form plus the K\"ahler form is positive definite. The
differences of the potentials lie between a multiple of $-\log r$ and $0$, so the tail bound holds
uniformly in the parameters.

% before: cv-tail-rate
The change of an arithmetic intersection number between two parameters is an integral of the
difference of two regularized potentials against a product of first Chern forms. Near the boundary
the difference is of size $\log r$, which the uniform tail bound controls; away from it the
difference is small by the bounds of an analytic system.

% before: cv-condition-convergence
The arithmetic intersection numbers are defined as limits, and the theorem needs the rate of
convergence at the parameter $e^{-1/s^2}$ used for the perturbed weights. The rate comes from
replacing the metrics one factor at a time and bounding each change by an integral over the
manifold.

% before: cv-condition-comparison
Two models realized on the same manifold whose bundles differ only by vertical divisors at a few
primes have nearly the same arithmetic intersection numbers. The difference is bounded by the
logarithms of those primes, weighted by the multiplicities; this gives the bound
$C\log B$ at level $B$.

% before: cv-nef-positivity
The arithmetic intersection numbers of the next theorem need the tautological bundles of dominant
weight to be nef on the generic fibre of each model. This comes from the analytic system the model
is realized on: on a curve, the degree of a realized bundle is the integral of a first Chern form
bounded below by $-\varepsilon(t)\omega$ for every $t$, and letting $t$ tend to $0$ makes it
nonnegative.

% before: cv-origin-hodge-component
The models in the tower of the next theorem are tied to the moduli scheme through covers of its
flag bundle, on which the Lie algebra of the semi-abelian scheme carries complete flags. At the
primes of a finite set the de Rham lattice is replaced by the Hodge--Tate lattice of the torsion, and
the graded pieces of the flags are the line bundles whose intersection numbers the theorem
computes.

% before: cv-origin-tower
The volume theorem is stated for a tower of flag models together with the data that tie each model
to the moduli scheme. For each component these data give one analytic system on which the models of
all levels are realized, Hodge components at the primes of $S_B$ whose root lines carry the
tautological bundles, and a reference model with an ample bundle. A good origin also satisfies the
smoothness, flatness and Hodge--Tate conditions that the proof uses.

% before: cv-origin-realization
The geometric degrees and the arithmetic intersection numbers of the models are computed on the
analytic systems. For this each model must be realized on the analytic system of its component,
with a multiplicity and an integer $k_A$ that do not depend on the level.

% before: cv-origin-effective
The perturbed weights $\lambda-s\nu$ of the volume theorem need a nonzero integral section of a
multiple of the class $\ell(\nu,-1)$ on each model. The section is a product of the generators of the
root lines with an element of the function field of the component, and the effectivity condition of
the origin makes it regular on the generic fibre.

% before: cv-volume-global-model
The volume theorem and the construction of towers start from a model of the moduli scheme over the
ring of integers of the coefficient field. At each prime this model is compared with a system of local
models at all levels prime to that prime, which carry the degenerations, the Hodge--Tate data and the
comparison with Siegel moduli schemes used by the origin of the tower.

% before: cv-origin-trivial-flag
At a prime of $S_B$ the Hodge--Tate data of the origin give frames of the flag of root lines on each
chart that are invertible modulo a power of $p$ and change by triangular matrices between charts. The
definition below records this as a condition on the multidegrees of the root lines alone. It is what the
nefness of the tautological bundles on the fibres at these primes is derived from.

% before: cv-origin-trivial-flag-divisor
The Hodge--Tate data give a trivial flag modulo $p^{R_p-2}$, while the degrees on curves in the fibre at
$p$ need one modulo $p$. Passing to a divisor of the modulus keeps the same frames.

% before: cv-origin-trivial-flag-section
A curve in the fibre of a flag model at $p$ is a curve over $\mathbf F_p$, and a family of chart
elements that glue modulo $p$ defines a section of a line bundle on it. If the section does not vanish
at the generic point of the curve, the line bundle has nonnegative degree there.

% before: cv-origin-trivial-flag-degrees
Minors of the frames on the first $i$ columns glue modulo $p$ to sections of the dual of the sum of the
first $i$ root lines, and the inverse of the determinant to a section of the sum of all of them. One
minor and the inverse determinant are nonzero on a given curve, which bounds the partial sums of the
degrees and fixes their total.

% before: cv-origin-trivial-flag-nef
A dominant weight pairs with the degrees of the root lines through decreasing coefficients. Abel
summation turns the bounds on the partial sums into the nefness of the tautological bundle of that
weight on the fibre at $p$.

% before: cv-volume-theorem
The slopes of automorphic forms of regular weight on a unitary Shimura variety are
controlled by arithmetic intersection numbers of tautological bundles on a flag bundle over
its toroidal compactification. The statement below computes these numbers up to a small
error and bounds the mean slope from below in terms of the root discriminant of the host
field and $L'/L(1,\chi)$.

% before: cv-origin-components
The components of a tower built on a model are its irreducible components, each with its generic
point. A tower has a nonempty finite set of components, so the model must have finitely many
irreducible components and at least one.

% before: cv-origin-reference
The data of a good origin at one component start from a Hodge component at the empty set whose flag has
degree $1$. The construction also needs a boundary ideal with normal crossings in smooth charts, and a block of
coordinates that embeds the generic fibre. Together these form a reference geometry.

% before: cv-origin-reference-char
The coordinates of a reference geometry take values in a field of characteristic zero, the function field of a
component of the model. The model is flat over the ring of integers of the coefficient field, and this gives the
characteristic.

% before: cv-origin-reference-lattices
The constructions of a reference geometry compare the Lie algebras of semi-abelian schemes over several models of one
component. The Lie algebra at the generic point is a vector space over the function field, and the Lie algebras at the
other points are modules over their local rings inside it. Lattice families, their pullbacks and their flags describe these
modules.

% before: cv-origin-reference-pullback
Each model built below carries the pullback of the semi-abelian scheme of $\mathcal X$. Pullbacks keep the action of
$\mathcal O_E$ and the cubical line bundle, and they compose.

% before: cv-origin-reference-lie
Tangent vectors are morphisms from the spectrum of the dual numbers, and they form the Lie algebra. Homomorphisms and the
action of $\mathcal O_E$ act on them by composition.

% before: cv-origin-reference-lie-bundle-base
The Lie algebra of a group scheme at a point with values in a ring is the module of tangent vectors there. For a smooth group
scheme it is the dual of the invariant differentials along the unit section, a finite projective module, and it commutes with
base change.

% before: cv-origin-reference-lie-pull
A morphism of semi-abelian schemes over a dominant morphism of models gives a map of Lie algebras at the generic points,
semilinear over the extension of function fields. For a smooth group scheme the Lie algebra commutes with base change, so this
map carries the Lie algebras at the points of one model to those of the other.

% before: cv-origin-reference-lie-same
When the two models have the same function field, the map of Lie algebras is an isomorphism. The Lie algebra of the pullback
is then the Lie algebra of the original scheme, with the modules at the points generated by those below.

% before: cv-origin-reference-lie-bundle
The Lie algebras of a smooth group scheme at the points of its base are free modules that vary like a vector bundle. The
same holds for their duals, the invariant differentials.

% before: cv-origin-reference-integral
A boundary ideal with normal crossings comes with charts that are smooth over $\mathbf Q$. The local rings of the model at the
points over no prime are localizations of these smooth algebras, so they are normal.

% before: cv-origin-reference-eigen-split
The action of the ring of integers of a number field splits a vector space into eigenspaces by idempotents whose
denominators divide the discriminant. Where the discriminant is a unit, these idempotents preserve a stable lattice family
and split it into vector bundles.

% before: cv-origin-reference-eigen
The labels give the eigenspaces of the Lie algebra of an $\mathcal O_E$-semi-abelian scheme. On a model whose local rings over
no prime are integrally closed, the values of the labels are regular there, and the Lie algebras at these points split
along the eigenspaces.

% before: cv-origin-reference-normalization
Some of the models below are normalizations over the points of interest. Over a point whose local ring is already integrally
closed, such a model has one point, with the same local ring.

% before: cv-origin-reference-refinement-plucker
A lattice family becomes free where one maximal minor of its generators divides all the others. Cramer's rule gives the
freeness from such a minor, and over a local ring a free family of generators produces one.

% before: cv-origin-reference-refinement-twist
On a scheme with coordinates, a locally finitely generated lattice family is generated on each chart by one finite set of
vectors. After a common twist by a power of the coordinate of the chart, a single finite family serves all the charts at once.

% before: cv-origin-reference-refinement-flatten
The maximal minors of the twisted generators form a new family of coordinates. Where the lattice family is free its point lies in
a chart of this family, and on any model over the given one the pullback is a vector bundle at the points in such a chart.

% before: cv-origin-reference-refinement-join
Joining further families of coordinates to those of a model gives a model on which every point lies in a chart of each family.
Over a point that already lies in such charts, the new model has exactly one point, with the same local ring.

% before: cv-origin-reference-refinement
The projections of the Lie algebras onto the eigenspaces are vector bundles only over the points over no prime. A refinement
of the model makes them vector bundles everywhere and gives charts for given families of coordinates. It is a normalization
over the points where these projections are already free.

% before: cv-origin-reference-base-over
The models of the generic fibre and of the integral model come with coordinates over $F_c$ and over $\mathcal O_c$. Their charts
are described by the algebra over that ring generated by the ratios, in place of the ring generated over $\mathbf Z$.

% before: cv-origin-reference-base-integral
Coordinates over a ring that is finitely generated as an abelian group can be replaced by coordinates over $\mathbf Z$. It is
enough to multiply them by generators of the ring, which are integral over $\mathbf Z$.

% before: cv-origin-reference-base-rings
A model over $\mathcal O_c$ and a model over $F_c$ with the same coordinates have the same local rings away from the primes.
Clearing the denominators of $F_c$ moves a point of one chart ring to a point of the other with the same units.

% before: cv-origin-reference-base-projective
A projective scheme over a noetherian ring has coordinates at a point whose local ring is the field. They are the homogeneous
coordinates of the point, and the charts are the standard affine pieces of projective space, cut down to the component of the
point.

% before: cv-origin-reference-base-product
The resolution of the generic fibre gets its coordinates from a projective embedding, and those of the integral model have to be
added to them. Multiplying the two families gives coordinates again, as soon as every point lies in a chart of the second family.

% before: cv-origin-reference-base-patch
The boundary of the resolution is described locally by subrings of the function field. Each is a smooth algebra over $\mathbf Q$
with coordinates, and its localizations are the local rings of the points, with the boundary cut out by some of the coordinates.

% before: cv-origin-reference-base-patches
The charts of the resolution with normal crossings give such patches. A chart is cut down by one function until its ring embeds in
the function field, and then it is a patch through every one of its points.

% before: cv-origin-reference-base-situation
The boundary ideal of the model is built chart by chart from the normal crossings patches of the resolution. It uses the chart
algebras over $\mathbf Q$, their localizations at one element, and the forms whose values are not units at the boundary points.

% before: cv-origin-reference-base-localization
At a point there are a patch and a chart algebra, two finitely generated algebras with the same localization there. As for two
affine neighbourhoods of a point, inverting one element in each makes them equal.

% before: cv-origin-reference-base-shrink
Inverting one element of a normal crossings patch gives a patch again, with the same coordinates. Its primes are those of the patch
that avoid the element, with the same localizations.

% before: cv-origin-reference-base-chart
Away from the primes, the primes of a chart algebra over $\mathbf Q$ are the points of the chart. The localization at the prime of
a point is its local ring.

% before: cv-origin-reference-base-cover
Comparing a patch at a point with the chart algebra there gives a patch that is the chart algebra with one element inverted. So the
patches can be taken inside the charts of the coordinates.

% before: cv-origin-reference-base-points
The primes of such a localized chart algebra are again the points of the chart at which the inverted element is a unit. At these
points the patch decides membership in the boundary through its first coordinates.

% before: cv-origin-reference-base-separate
Off the boundary, a patch gives an element of the chart algebra that is a unit at a given point and at no boundary point of the
chart. It is the product of the first coordinates of the patch, with its denominator cleared.

% before: cv-origin-reference-base-form
Clearing denominators and multiplying by the coordinate of the chart turns such an element into a homogeneous form. The extra factor
keeps the form a non-unit at the boundary points of every other chart.

% before: cv-origin-reference-base-zeros
A complex point gives a homomorphism from the local ring of its centre to $\mathbf C$, which sends the non-units to zero. So a form
vanishes at the coordinates of the point when its value at the centre is not a unit.

% before: cv-origin-reference-base-model
Condition (3) of a reference geometry asks for a cover of each chart by patches, with an ideal whose values cut out their
boundaries. Patches at all the points of a chart give such a cover, as no prime of the chart algebra avoids them all.

% before: cv-origin-reference-base-ideal
The boundary forms generate the boundary ideal. They separate each point off the boundary from the boundary, which gives both the
normal crossings and the complex points of the boundary.

% before: cv-origin-reference-base
The construction starts from a resolution of the generic fibre of $\mathcal X$ on which the complement of the abelian locus has
normal crossings. A single block of coordinates, with the function field of the component, covers it.

% before: cv-origin-reference-rank
The moduli problem fixes the signature of the action of $\mathcal O_E$ on the Lie algebra at each point of the moduli scheme.
The eigenspaces of the Lie algebra at the generic point of a model have the dimensions $p$ and $q$ that the signature
prescribes.

% before: cv-origin-reference-coords
The first step towards a reference geometry is a model of one component of $\mathcal X$ with coordinates. On it
the eigenspaces of the Lie algebra of the semi-abelian scheme extend to vector bundles, and the locus where the
scheme is not abelian is cut out with normal crossings.

% before: cv-origin-reference-flag-bundle-twist
The coordinates of the flag bundle are built from linear forms that generate the dual of a vector bundle on every chart. A frame on
each piece of a finite cover of a chart gives such forms after a twist by a power of the coordinate of the chart, with one power
for all the charts.

% before: cv-origin-reference-flag-bundle-inverse
The matrix of the twisted forms on a frame has a left inverse. At a point the inverse has regular entries, and at the generic
point it is a left inverse over the field.

% before: cv-origin-reference-flag-bundle-coords
The minors of the generic flag, written in the twisted forms, give the coordinates of the flag bundle. Their ratios generate the
field of rational functions in the entries of the unipotent matrices.

% before: cv-origin-reference-flag-bundle-valuative
The model built from these coordinates is a Proj of a graded ring generated by finitely many elements of degree one. It is therefore
proper over $\mathbf Z$, and it is integral.

% before: cv-origin-reference-flag-bundle-onto
A morphism from such a model to a scheme with coordinates is surjective. A valuation ring dominating the local ring of a point of
the target gives, by the valuative criterion of properness, a point of the model over it.

% before: cv-origin-reference-flag-bundle-minors
In a chart of the flag bundle one minor of each flag is a unit times the others. Cramer's rule turns this into integral vectors in
the flag that are dual to integral vectors of the frame.

% before: cv-origin-reference-flag-bundle-adapted
Over a local ring, such dual systems give an invertible matrix whose columns span the flag step by step. Its columns are chosen one
at a time, with residues independent of the earlier ones.

% before: cv-origin-reference-flag-bundle-frame
The smoothness of the flag bundle is checked in a chart, where the generic basis is written in a local frame. In a chart of the base,
after inverting one function, the matrix of the twisted forms factors through matrices over the ring of the chart.

% before: cv-origin-reference-flag-bundle-normalized
Over a local ring, the minors of a matrix with a left inverse control those of its second factor. This gives an order of the rows in
which every leading minor divides all the minors of its size.

% before: cv-origin-reference-flag-bundle-cell
With the rows in that order, an LU decomposition expresses the ratios of the minors through the entries of a unipotent matrix. These
entries are algebraically independent, so the ring of the chart of the flag bundle is a localization of a polynomial ring.

% before: cv-origin-reference-flag-bundle
On the bundle of complete flags of a family of vector bundles the flags become tautological. Its generic point is given by a
unipotent matrix for each bundle, whose entries below the diagonal are algebraically independent over the function field of
the base.

% before: cv-origin-reference-flag-linear
Generic flags and flag bundles are moved along base changes and isomorphisms of vector spaces. The degree of a generic flag
does not change, and in degree $1$ the entries of the unipotent matrices generate the field.

% before: cv-origin-reference-flag-coords
On the flag bundle, the coordinates of the base are joined to those of the fibre. The charts of the joined point cover the
flag bundle, and the coordinates of a complex point restrict to coordinates of its image.

% before: cv-origin-reference-flag-snc-chart
The smooth charts of a boundary ideal are localizations of the rings of ratios of a chart. A smooth morphism between two schemes
with coordinates maps a chart of the source into a chart of the target, so it makes the ring of ratios of the source chart a
smooth algebra over that of the target chart.

% before: cv-origin-reference-flag-snc
A smooth morphism whose generic fibre adds algebraically independent coordinates raises the dimension of the smooth charts by
their number. The normal crossings of the boundary ideal pass to the charts of the joined point.

% before: cv-origin-reference-flag-constants
The flags of the eigenspace bundles are given on the pullbacks of the lattices of the base. At the points over no prime these
pullbacks are the eigenspace parts of the Lie algebras of the new semi-abelian scheme, which therefore carry the flags.

% before: cv-origin-reference-flag
The flag of a Hodge component lives on the bundle of flags of the eigenspaces of the Lie algebra. Passing to this
bundle raises the dimension from $|\Psi|pq$ to $|\Psi|N(N-1)/2$, and the generic bases of the flag are its
coordinates.

% before: cv-origin-reference-descent-quotients
A flag of a vector bundle is a flag of subbundles when every quotient by a step of the flag is again a vector bundle. Then
each step is locally a direct summand.

% before: cv-origin-reference-descent-model
A Hodge component needs the eigenspaces of its Lie algebra to carry flags at every point. A refinement applied to the
quotients by the steps of the flags makes all of them vector bundles, and the flags then become flags of subbundles.

% before: cv-origin-reference-descent-roots-piece
Line bundles are generated by global sections after a twist, and the proof works on the affine charts of a scheme with
coordinates. Inside a chart, the points at which a ratio polynomial is a unit form a standard open set, and the elements
regular on it are the ratio polynomials divided by a power of that polynomial.

% before: cv-origin-reference-descent-roots-twist
A line bundle in the function field is generated near each point by one element. After multiplication by a power of the
coordinate of the chart, finitely many elements of the function field generate it at every point of every chart. These
elements become a new block of coordinates.

% before: cv-origin-reference-descent-roots
The root lines of a flag are line bundles, built from the graded pieces of the flags and their duals. Further blocks of
coordinates make each of them generated on every chart by a monomial in the coordinates times a fixed element of the
function field.

% before: cv-origin-reference-descent
The scheme built on the bundle of flags still has to be made into a Hodge component, with a flag and root lines
that have generators on every chart. Over the generic fibre the new scheme keeps the local rings of the old one,
so the boundary and the coordinates of complex points can be transferred.

% before: cv-origin-snc-charts
The lemmas that follow argue chart by chart. On the affine open subscheme of a chart, the functions regular at a point
are the fractions of the ring of ratios with denominator a unit there, and the points correspond to the primes of that
ring.

% before: cv-origin-segre-coordinates
Adjoining blocks to a point changes the charts of a scheme, and schemes are compared through the elements of the
field that are regular at their points. A finite family of elements gives charts, and the scheme has the family as
coordinates when its charts are affine open subschemes whose rings are the rings of ratios. The covering of
Definition~\ref{cv-origin-hodge-component}(1) is the case of the Segre products of a point.

% before: cv-origin-segre-regular
Regularity at a point is read off the local ring, and the local rings sit inside the field when they embed in it.
Regular elements then stay regular on an open neighbourhood, the sections over an open set are the elements regular
on it, and a morphism compatible with the points of the field carries regular elements and units back.

% before: cv-origin-segre-model
A family with a chart whose ratios give the whole field by fractions has a scheme with these coordinates. It is the
Proj of the graded ring of powers of the span of the family, whose standard charts are the spectra of the rings of
ratios.

% before: cv-origin-segre-map
Morphisms into a scheme with coordinates are built chart by chart. If every point of the domain lies in a chart of
these coordinates, the ring of ratios of each chart is regular on an open subset of the domain, which gives a
morphism into the affine chart, and these morphisms glue.

% before: cv-origin-segre-unique
The morphism is unique when its domain also has coordinates. Two morphisms compatible with the points of the field
send each point to points that specialize each other, and they agree on sections because the local rings of the
domain embed in the field.

% before: cv-origin-segre-join
Adjoining a second family multiplies the coordinates. A scheme with the products as coordinates exists, because a
chart of the first family at the image of the field already gives the field by fractions, and it maps to the
original scheme.

% before: cv-origin-reference-segre
The block that embeds the generic fibre is the Segre block, whose coordinates are the products of one coordinate
from each block. Adjoining it must keep the covering of the scheme by affine charts with rings of ratios.

% before: cv-origin-reference-snc
The boundary ideal is constructed for the coordinates of the bundle of flags, and the reference geometry has more
blocks. The normal crossings charts transfer to the larger point when the two schemes have the same local rings
over the generic fibre.

% before: cv-origin-reference-points
The boundary condition of a reference geometry is stated through the coordinates of complex points. Restricting
the coordinates to the blocks of a smaller point gives coordinates of the image point, when the local rings over
the generic fibre agree.

% before: cv-origin-reference-embeds
A block whose charts are affine with the rings of its ratios embeds the generic fibre, since local rings are then
localizations of these rings. This applies to the Segre block of any Hodge component.

% before: cv-origin-reference-exists
Every component of the model carries a reference geometry. The construction chains the coordinates of the
component, the bundle of flags, the Hodge component over it and the Segre block.

% before: cv-origin-analytic
The analytic side of a component is an analytic system realized on the complex points of the reference geometry by
a bijective map of degree $1$. Its metrized line bundles are those of the Hodge metrics of the root lines, up to
the singular weights of the system.

% before: cv-origin-hodge-norm-gram
The Hodge metrics are defined through the Hodge inner products of tangent vectors at the complex points of a Hodge
component. These come from a uniformization of the fibre by a complex vector space, in which the line bundle has a
factor of automorphy of Appell--Humbert type. A family in the generic Lie algebra then has a Gram matrix at every
complex point at which its members have values.

% before: cv-origin-hodge-norm-eigen
At a point over no prime the Lie algebra splits into its eigenspaces for $\mathcal O_E$. The difference of two labels at
an element of $\mathcal O_E$ is an algebraic integer whose inverse is regular away from the primes, so the eigenspace
components of a vector of the Lie algebra at the point stay in it.

% before: cv-origin-hodge-norm-values
A vector of the Lie algebra at a point has a value at every complex point with that centre. As the Lie algebra of a
smooth group scheme commutes with base change, these values span the tangent space of the fibre, and a linear form
that is regular on the Lie algebra at the point gives a linear form on that tangent space.

% before: cv-origin-hodge-norm-form-ample
The fibres of the semi-abelian schemes will be shown to carry ample line bundles in a concrete form: finitely many sections of
a power of the line bundle whose nonvanishing loci are affine and cover, with ratios that generate their rings. This form
is what produces theta functions on a uniformization.

% before: cv-origin-hodge-norm-form-descent-ample-affine
The descent of ampleness through a field extension needs every set where a section of a power of an ample line bundle does
not vanish to be affine. On a scheme that is universally closed over an affine scheme, the ratios of finitely many covering
sections to a given section generate the unit ideal: a common zero would survive as a point of a closed image under a
projection.

% before: cv-origin-hodge-norm-form-descent-ample-embedding
Relative ampleness by sections asks for sections whose ratios generate the rings of their nonvanishing loci. A power of an
ample line bundle on a scheme proper over a ring has such sections. Sections of a high power that extend generators of the
rings of affine nonvanishing loci give an immersion into projective space, and properness makes it closed.

% before: cv-origin-hodge-norm-form-descent-ample
Over a ring, relative ampleness by sections of a line bundle on a proper scheme can be tested pointwise: every point should
lie in an affine set where some section of a power does not vanish. This pointwise form passes to and from field extensions
through the sections, which commute with flat base change.

% before: cv-origin-hodge-norm-form-descent
A complex point of the model in the image of the moduli scheme lifts to the moduli scheme only after a field extension. The
fibre there is a base change of the universal abelian scheme, which is proper with an ample line bundle, and these
properties must then descend to the original field.

% before: cv-origin-hodge-norm-form-fibre
At an abelian point of a Hodge component the fibre is a fibre of the model's semi-abelian scheme over its generic fibre,
since the ring of integers of the coefficient field embeds in $\mathbf C$. The abelian part of the generic fibre is the moduli
scheme, so the fibre is proper with an ample line bundle.

% before: cv-origin-hodge-norm-form-torus-exp-log
A holomorphic local law that is commutative and associative becomes addition after a change of coordinates. The change is
built from the doubling map: halving is a contraction near the origin, and its iterates multiplied by powers of $2$ converge
to a logarithm of the law.

% before: cv-origin-hodge-norm-form-torus-exp
A chart at the origin carries the group law of the complex points to a holomorphic local law, and a logarithm of that law
makes it additive. The inverse of the logarithm is then a homomorphism near the origin, which halving extends to the whole
tangent space, and compactness and divisibility of the group make it a uniformization.

% before: cv-origin-hodge-norm-form-torus-points
The complex points of a scheme carry the coarsest topology in which the regular functions are continuous. For a group scheme
the group law is continuous, as regular functions on a product of affine charts are generated by those of the factors. An
ample line bundle in the concrete form of sections confines every point to one of finitely many compact pieces.

% before: cv-origin-hodge-norm-form-torus-chart-etale
Near a point, an algebra that is \'etale over a polynomial ring is described by as many equations as extra variables, with
an invertible Jacobian. The holomorphic inverse function theorem then solves these equations near the point, as
holomorphic functions of the coordinates of the polynomial ring.

% before: cv-origin-hodge-norm-form-torus-chart
A smooth scheme is \'etale over affine space near each point. Solving the \'etale equations holomorphically gives a chart
of the complex points whose coordinates are regular functions, and these coordinates detect tangent vectors, as an \'etale
algebra has no nonzero derivations over the polynomial ring.

% before: cv-origin-hodge-norm-form-torus-divisible
The uniformization is surjective only if the group of complex points is divisible. Multiplication by an integer is finite
and flat, so it is open and closed and hence surjective, and its finite fibres have complex points.

% before: cv-origin-hodge-norm-form-torus
A complex abelian variety is a quotient of its tangent space by a lattice. Coordinates of a holomorphic chart at the origin
whose derivative along the uniformization is the identity will identify tangent vectors with vectors of $\mathbf C^g$.

% before: cv-origin-hodge-norm-form-tangent
A tangent vector at the origin is a ring homomorphism to the dual numbers, and so is the jet of a vector of $\mathbf C^g$
through the uniformization. Over an affine chart with the coordinates of the uniformization these two kinds of
homomorphism are the same.

% before: cv-origin-hodge-norm-form-appell-datum
Pulled back along a uniformization, a line bundle with a cubical structure becomes a cocycle of holomorphic functions on the
uniformizing space together with a holomorphic trivialization of its cube on triples of points. Both are invariant under the
lattice. The Appell--Humbert form of the line bundle is derived from these analytic data alone.

% before: cv-origin-hodge-norm-form-appell-pair
The cubical structure lives on the triple product of the group scheme, so its values along the uniformization are holomorphic
only if a triple of holomorphic families of points is a holomorphic family in the product. Regular functions on an affine
chart of a product are sums of products of functions from the factors.

% before: cv-origin-hodge-norm-form-appell-analytic
The transition functions and the cubical structure of the line bundle, read through the uniformization, form a periodic cube
datum. They inherit holomorphy from the uniformization and periodicity from its kernel.

% before: cv-origin-hodge-norm-form-appell-faces
A holomorphic function on the uniformizing space that is invariant under a spanning lattice is constant. Ratios of the
trivializations of the cube that do not depend on the charts are such functions, and they give two identities near the
origin: the cube is trivial on the faces, and it satisfies a doubling relation.

% before: cv-origin-hodge-norm-form-appell-local
Near the origin the two identities determine the trivialization of the cube up to an exponential of a third difference. The
logarithm vanishes to order three, a doubling equation for the third difference can be solved by a convergent series, and a
function of order three that satisfies the homogeneous equation vanishes.

% before: cv-origin-hodge-norm-form-appell-global
The local description of the cube produces trivializing functions of the line bundle near the origin. The cube identity then
extends them from a ball to the ball of twice the radius, by expressing their values at a point through their values at
half of it.

% before: cv-origin-hodge-norm-form-appell-factor
Global trivializing functions compatible with the cube have a factor of automorphy under the lattice. The cube identity at translated
triples makes this factor an exponential of a linear form, up to a constant.

% before: cv-origin-hodge-norm-form-appell-normal
A factor of automorphy that is a constant times the exponential of a linear form can be brought to Appell--Humbert form by
multiplying the trivializing functions by the exponential of a quadratic polynomial. In coordinates, the linear forms split
into a symmetric part and a Hermitian part; the Hermitian part gives the Hermitian matrix, and the rest of the constant
gives a unitary character.

% before: cv-origin-hodge-norm-form-appell
The Hodge inner products are read off the factor of automorphy of the line bundle along the uniformization. The cubical
structure of the line bundle brings this factor to Appell--Humbert form, with a Hermitian matrix and a unitary character.

% before: cv-origin-hodge-norm-form-theta
Sections of a power of an ample line bundle become theta functions on the uniformization, with the factor of automorphy raised
to that power. Their ratios generate the rings of affine charts, so no nonzero direction leaves all of them invariant.

% before: cv-origin-hodge-norm-form-positive
The Hermitian matrix of the factor of automorphy must be positive definite for the Hodge inner products to be a metric. A
nonzero theta function grows at most like the exponential of the Hermitian form, so the form cannot be negative anywhere,
and a degenerate direction would leave every theta function invariant.

% before: cv-origin-hodge-norm-form-lie
The tangent space at a complex point of the base is identified with the tangent space at the origin of the fibre. Through the
uniformization, differentials then give a linear isomorphism with $\mathbf C^g$, as they add and scale with tangent vectors.

% before: cv-origin-hodge-norm-form-lift-open
The lift of a holomorphic homomorphism is built first near the origin, where the uniformization has a local inverse. That
inverse can be composed with the homomorphism only if points near the origin go to images of points near the origin. A
Baire category argument gives this from functions that separate the points of an affine chart.

% before: cv-origin-hodge-norm-form-lift-real
A local lift of a homomorphism is additive near the origin up to elements of the kernel of the uniformization. As that kernel
is discrete, the local lift is additive near the origin, and halving extends it to a real-linear lift.

% before: cv-origin-hodge-norm-form-lift
Two uniformizations of the same fibre must be compared. A holomorphic homomorphism from a complex vector space to the
complex points of an abelian variety lifts to a linear map to its uniformizing space.

% before: cv-origin-hodge-norm-form-pullback-trivial
Trivializing functions of two isomorphic line bundles differ by a holomorphic function without zeros on the uniformizing
space. Its factor of
automorphy has the Appell--Humbert shape, and the function and its inverse grow at most like the exponential of the
Hermitian form. This forces the form and the character to be trivial.

% before: cv-origin-hodge-norm-form-pullback
The Appell--Humbert data of a line bundle are determined by the line bundle. When a homomorphism pulls one line bundle back to
another, its linear lift carries one Hermitian form to the other.

% before: cv-origin-hodge-norm-form-compare
The Hodge inner product is defined through some uniformization, so it must not depend on the choice. Two uniformizations of the
fibre differ by a linear lift of the isomorphism of the fibres, which matches their differentials and their Hermitian forms.

% before: cv-origin-hodge-norm-form
At a point over which the semi-abelian scheme is abelian, the fibre is a complex abelian variety with an ample line
bundle. Its Hodge inner products are then the values of one positive definite Hermitian form on the tangent space.

% before: cv-origin-hodge-norm-adapted
Two bases of a lattice adapted to the same flag differ by a triangular matrix over the local ring. The leading minors
of their Gram matrices then differ by the squared absolute values of its diagonal entries, so the Hodge norm of a
graded piece is fixed up to the value of a unit.

% before: cv-origin-analytic-hodge-points
The Hodge metrics are functions of the coordinates, so a complex point must be determined by its coordinates, and
scaling the coordinates in a block must keep the point. Complex conjugation of a point conjugates its coordinates and the
values of functions and tangent vectors there.

% before: cv-origin-analytic-hodge-conjugate
The Hodge metrics are to be invariant under conjugation. Conjugating a uniformization of the fibre at a point gives one
of the fibre at the conjugate point, with the conjugate Hermitian matrix, so Hodge inner products are conjugated and
Hodge norms are kept.

% before: cv-origin-analytic-hodge
The metrized line bundles of the analytic data start from Hodge metrics on the root lines. These are positive
homogeneous functions on all of $\hat C$ that give the Hodge norm of the generator at the abelian points.

% before: cv-origin-analytic-local-tame
The Hodge potentials will be assembled from logarithms of holomorphic functions and from potentials of eigenbundles, by
sums and telescoping. Tameness near the boundary, and plurisubharmonicity at the origin of a chart, survive these
operations, and the logarithm of the absolute value of a holomorphic function without zeros has both.

% before: cv-origin-analytic-local-roots
In a normal crossings chart each root line is a graded piece of the flag of an eigenbundle, or its dual. Its Hodge metric
is therefore a holomorphic factor times the square root of a ratio of two consecutive leading minors of the Gram matrix of a
frame, and the logarithms of these minors are the potentials that carry the estimates.

% before: cv-origin-analytic-local-central
The last Hodge potential involves all root lines of one embedding, and the potentials of the eigenbundles then cancel in
pairs. What is left is the determinant of the Hodge bundles, which is flat, so near an interior point its logarithm is
pluriharmonic.

% before: cv-origin-analytic-local
An analytic system needs bounds on its weights near the boundary and plurisubharmonicity in the interior. In each
normal crossings chart these come from the Hodge potentials of the root lines, and the potential of the last root
line is pluriharmonic.

% before: cv-origin-analytic-system-pullback
The analytic system is built on the complex points of the kernel of the coordinates, and its line bundles are
pulled back from metrized line bundles on the product of projective spaces. The Fubini--Study bundles give the smooth
metrics, and the positivity of a pullback is tested by Levi forms in charts.

% before: cv-origin-analytic-system-arith
The construction needs the kernel of the coordinates to be an arithmetic subvariety of dimension $D+1$. Condition
(3) of a reference geometry gives a smooth chart algebra with $D$ independent differentials, and so the
transcendence degree of the function field.

% before: cv-origin-analytic-system-boundary
The boundary of the analytic system is the zero set of the boundary ideal on the complex points. The construction of
the boundary bundle uses that this set is closed and stable under the torus and under conjugation, which follows
from the description of the abelian locus in condition (4).

% before: cv-origin-analytic-system-ideal
The boundary bundle is built from forms that vanish on the boundary, and the boundary ideal need not contain all of
them. Replacing it by the ideal of the polynomials that vanish on the torus orbits of the boundary points keeps
condition (3) and the zero set.

% before: cv-origin-analytic-system-crossings-dehomogenize
The components of the boundary are minimal primes of a multihomogeneous ideal, while the charts of the complex points
are affine. Putting one variable of each block equal to $1$ matches the multihomogeneous primes that avoid these
variables with primes of the affine chart, and keeps track of which forms lie in them.

% before: cv-origin-analytic-system-crossings-regular
The charts of the complex points are built from regular local rings of the affine charts. Without the smoothness of
the chart algebras, condition (3) still gives, near each point, a localization of the chart algebra whose module of
differentials is spanned by $e$ elements, and this bounds the tangent spaces.

% before: cv-origin-analytic-system-crossings-params
The elements $z_1,\dots,z_e$ of condition (3) are rational functions on the arithmetic variety. Shifted by their values
at a complex point, they become local parameters there, and the boundary passes through the point along the zero sets
of those among the first $n$ that vanish at it.

% before: cv-origin-analytic-system-crossings-local
In a regular local ring the elements of a regular system of parameters are pairwise distinct primes. So an ideal whose
primes are detected by some of these parameters has the product of those parameters as radical, up to a factor that
does not vanish at the point.

% before: cv-origin-analytic-system-crossings-components
The local parameters describe the boundary near one point, while the sets $W_l$ of the normal crossings divisor have to
be global. They are the zero sets of the minimal primes over the boundary ideal, and near each point each of them is the
zero set of one parameter.

% before: cv-origin-analytic-system-crossings-mult
A multihomogeneous form whose zeros on the complex points lie in the boundary is, near each point, a unit times a
monomial in the boundary parameters. The exponent at each parameter is the order of the form along the corresponding
minimal prime, so it is the same at every point of that component.

% before: cv-origin-analytic-system-crossings
Condition (3) makes the chart algebras smooth over $\mathbf Q$ and puts the boundary in normal crossings. On the
complex points this gives charts in which the boundary is a union of coordinate hyperplanes, and in which every form
that vanishes on the boundary is divisible by the product of these coordinates.

% before: cv-origin-analytic-system-section
From these local forms a line bundle is glued, with a section whose zero set is the boundary. Its metric is built
from finitely many forms that vanish on the boundary, together with their conjugates, so that the norm of the
section is invariant under the torus and under conjugation.

% before: cv-origin-analytic-system-snc
The three previous lemmas combine to give the boundary bundle on the manifold of complex points. Its section has
norm at most $1$ and a divisor with simple normal crossings, as an analytic system requires.

% before: cv-origin-analytic-system-positive
The metrics of the analytic system are pullbacks of Fubini--Study bundles. Their functions are local potentials of
metrics because the lifts agree up to the torus, and the K\"ahler form and the bundle of the block $k_0$ are positive
because the lifts are immersions.

% before: cv-origin-analytic-system-charts
The local bounds of Lemma~\ref{cv-origin-analytic-local} are stated on polydiscs in which the boundary is a union of
coordinate hyperplanes. Normal crossings charts of the section give such polydiscs, on which the function $r$ of the
section is bounded below by the local logarithmic function up to a constant.

% before: cv-origin-analytic-system-weight
The singular weights of an analytic system and their regularizations satisfy the bounds of
Definition~\ref{cv-proportionality-analytic-system}(1) one bundle at a time. Near the boundary the weights are
controlled by tame functions, whose derivatives may grow like the exponential of the local logarithmic function.

% before: cv-origin-analytic-system-enlarge
The weights of the different bundles come with different constants, while an analytic system has one pair of
constants. The bounds survive a larger constant and a smaller range of $t$, and a smooth weight is a regularized
weight with constant regularization.

% before: cv-origin-analytic-system-weights-family
The regularized weights are built in two steps: the singular weight off the boundary, and then smooth functions that
approximate it. The conditions on the second step alone are those of an analytic system that involve the parameter $t$.

% before: cv-origin-analytic-system-weights-conjugation
The bundles of the realization have to be invariant under conjugation, and the weights are built on the manifold
$\tilde X$. Complex conjugation of the complex points descends to $\tilde X$, where it is an antiholomorphic involution,
and it changes the Fubini--Study potentials only by pluriharmonic functions.

% before: cv-origin-analytic-system-weights-tame
The local bounds on the Hodge potentials are stated in normal crossings charts, one point at a time. Compactness turns
them into a bound by a multiple of $\log r$ on all of $\tilde X$, and gives bounds on derivatives in charts around every
point.

% before: cv-origin-analytic-system-weights-line
A smooth weight on $\tilde X$ that is invariant under the involution has to be realized by a metric on the multiprojective
space, invariant under conjugation. The extension of a smooth function to a metric of degree zero does this once the
metric is averaged with its conjugate.

% before: cv-origin-analytic-system-weights-singular-disc
Near the boundary the weight is subharmonic on discs that meet the boundary in finitely many points, where it may grow,
but more slowly than any multiple of the logarithm of the distance. Such points do not spoil the mean value inequality
on the whole disc, by the maximum principle applied after subtracting the Dirichlet solution and adding small multiples
of these logarithms.

% before: cv-origin-analytic-system-weights-singular-extension
In a polydisc the boundary is a union of coordinate hyperplanes. Moving one coordinate at a time to a circle where the
function is at least as large reaches a compact set away from the boundary, on which the function is bounded.

% before: cv-origin-analytic-system-weights-singular
A regularized weight is bounded above, while the tame bounds give only a bound by a multiple of $\log r$. Near the
boundary, plurisubharmonicity improves this growth to a bound from above.

% before: cv-origin-analytic-system-weights-flat-extend
A function whose Levi form vanishes is pluriharmonic, and a bounded pluriharmonic function extends across a set
contained in finitely many complex hyperplanes. In a chart where the boundary is a union of coordinate hyperplanes, this
extends a bounded flat weight smoothly across the boundary.

% before: cv-origin-analytic-system-weights-flat
When the weight is pluriharmonic off the boundary, no regularization is needed: it extends smoothly across the
boundary. The constant family then has the pluriharmonic local potentials that the last bundle of each flag needs.

% before: cv-origin-analytic-system-weights-regularize-atlas
The regularization works in finitely many charts of a fixed size, in each of which the boundary is a union of
coordinate hyperplanes. In such a chart the function $r$ differs by a bounded amount from an explicit sum of logarithms
of the coordinates.

% before: cv-origin-analytic-system-weights-regularize-gradient
The bounds on the derivative of the weight come in arbitrary charts around each point. Changing to the charts of the
atlas multiplies them by the norm of a holomorphic derivative, which is bounded near each point, and compactness makes
one constant work for all the charts.

% before: cv-origin-analytic-system-weights-regularize-global-notation
The regularization is built in polydiscs of the charts of the atlas, off the coordinate hyperplanes that make up the
boundary. Its tools are suprema over small polydiscs, which bound a plurisubharmonic function from above, and averages
against a smooth kernel, which make it smooth.

% before: cv-origin-analytic-system-weights-regularize-global-numbers
The loss of positivity $\varepsilon_C(t)$ decreases with $t$, tends to $0$, and is larger than $t$ for small $t$. These
elementary facts are used to compare the constants of the construction.

% before: cv-origin-analytic-system-weights-regularize-global-transfer
In a chart of the atlas the function $r$ and the explicit sum $r_{T,\mu}$ differ by a bounded amount. Bounds stated with
one of them then hold with the other, at the cost of a larger constant.

% before: cv-origin-analytic-system-weights-regularize-global-refine
The local potentials of the metric are defined only near each point, while the construction in a chart uses one
potential on the whole chart. Shrinking and recentring the charts gives an atlas with this property, with the same shape
of the boundary and the same kind of bounds on the weight.

% before: cv-origin-analytic-system-weights-regularize-global-extend
A plurisubharmonic function off the coordinate hyperplanes that is bounded above becomes plurisubharmonic across them
after it is truncated from below. Its suprema over polydiscs are then convex in the logarithm of the radius, and they
agree with the suprema of the original function off the hyperplanes.

% before: cv-origin-analytic-system-weights-regularize-global-oscillation
The supremum of the weight over a small polydisc changes little when the radius is multiplied by a bounded factor. The
logarithmic bound near the boundary and the convexity in the logarithm of the radius bound the change by
$\varepsilon_C(t)$.

% before: cv-origin-analytic-system-weights-regularize-global-mollifier
Averaging a plurisubharmonic function against a smooth kernel of radius $t$ gives a smooth plurisubharmonic function that
lies between the function and its supremum over the polydisc of radius $t$. The mean value inequality bounds it from
below by that supremum, up to the oscillation over a slightly larger polydisc.

% before: cv-origin-analytic-system-weights-regularize-global-mollify
In one chart, the weight truncated from below at $-C\log\log(1/t)$ and averaged at scale $t$ is smooth and
plurisubharmonic. By the oscillation bound it stays within $\varepsilon_C(t)$ of the supremum of the weight over the
polydisc of radius $t$.

% before: cv-origin-analytic-system-weights-regularize-global-near
Where $r$ is at most $\log(1/t)$ up to a constant, the regularized weight should exceed the weight by at most $t^{1/3}$.
The derivative of the weight grows at most like a power of $r$ times $e^{r/2}$, so by the mean value theorem its
supremum over the polydisc of radius $t$ is close to its value at the centre.

% before: cv-origin-analytic-system-weights-regularize-global-chart
Together, the averaging and the bounds near the boundary give the regularization in one chart. Besides the bounds of a
regularized weight, it stays within $\varepsilon_C(t)$ of the supremum over the polydisc of radius $t$, a quantity that
can be compared between charts.

% before: cv-origin-analytic-system-weights-regularize-global-compare
The suprema over polydiscs of radius $t$ in two charts differ by at most $\varepsilon_C(t)$, once the difference of the
local potentials is taken into account. A change of charts maps a polydisc of radius $t$ into one of radius $C_1t$, and
the oscillation bound absorbs the factor $C_1$.

% before: cv-origin-analytic-system-weights-regularize-global-max
Local functions that are close on the overlaps of charts are glued with Demailly's regularized maximum. It is smooth and
convex, it keeps lower bounds on Levi forms, and it ignores arguments that are smaller than another one by more than its
parameter.

% before: cv-origin-analytic-system-weights-regularize-global-max-limit
The glued functions must converge with two derivatives off the boundary. The second derivative of the regularized
maximum with parameter $\eta$ is of size $1/\eta$ and meets only the differences of the first derivatives of its
arguments, so convergence holds when the parameter is large compared with the square of these differences.

% before: cv-origin-analytic-system-weights-regularize-global-glue
The local regularizations differ by at most $\delta(t)$ on the overlaps, and the regularized maximum of copies lowered
near the edges of their charts glues them. The lowering costs a multiple of $\delta(t)$ in the positivity of the
curvature.

% before: cv-origin-analytic-system-weights-regularize-global-bounds
The glued function inherits the bounds of the local regularizations up to multiples of $\varepsilon_{C}(t)$. Where $r$
exceeds $\log(1/t)$, the term $\varepsilon_C(t)r$ absorbs the difference from the weight, which is at least $-C_5\log r$.

% before: cv-origin-analytic-system-weights-regularize-global
With the atlas and one constant for the bounds on the weight and its derivative, the regularizing functions are built
chart by chart and then glued. The loss of positivity in this construction is the $\varepsilon(t)$ of an analytic system.

% before: cv-origin-analytic-system-weights-regularize
In general the singular weight is only plurisubharmonic off the boundary and has to be approximated by smooth functions
whose curvature is almost semipositive. The derivative bounds near the boundary control the loss of positivity in this
approximation.

% before: cv-origin-analytic-system-weights-symmetric
The regularization need not commute with the involution, and the bundles of the realization must be invariant under
conjugation. Averaging the regularization with its composite with the involution restores the invariance at the cost of
larger constants.

% before: cv-origin-analytic-system-weights
The weights of the bundles of the root lines are the differences between the Hodge potentials and the Fubini--Study
potentials. Tame bounds near the boundary and plurisubharmonicity in the interior give regularized weights, and an
average over conjugation keeps the bundles of the realization invariant under conjugation.

% before: cv-origin-analytic-system
From the local bounds in every normal crossings chart, the analytic system is built on the manifold of complex
points itself. The realization is then the identity map, which is bijective of degree $1$ with a holomorphic
inverse.

% before: cv-origin-analytic-exists
The two previous lemmas combine to give analytic data on every reference geometry of a global model. The
holomorphic inverse of $\pi$ is used for the immersion condition (3d) of a good origin.

% before: cv-origin-weight-birational
The sections of the effective weight are elements of the function field of the arithmetic variety of the
coordinates. Its dimension comes from the transcendence degree of the flag over the function field of the
component.

% before: cv-origin-boundary-points
Whether a form vanishes on the boundary is tested at complex points with prescribed values on a chart. A homomorphism
from the ring of ratios of a chart to a field gives a point with values in that field, and a point of the zero set of
the kernel gives a point with those coordinates.

% before: cv-origin-weight-boundary
The section of the effective weight must vanish where the semi-abelian scheme is not abelian. A multihomogeneous
form that vanishes on this boundary supplies the vanishing.

% before: cv-origin-weight-nef-degree
The positivity of the degree polynomial passes to the combination of the bundles of a realization through its
degree. The realization turns the integral over the complex points into an integral over its manifold, and the
map to the analytic system multiplies that integral by its degree.

% before: cv-origin-weight-nef-generic-fibre
An integral of first Chern forms over the complex points is first compared with the degrees of the components of
the fibre over a prime. An arithmetic intersection system gives the comparison: its recursion, applied to the
trivial bundle with a constant metric and to the constant sections $p$ and $1$, computes one arithmetic
intersection number in two ways.

% before: cv-origin-weight-nef-generic-conservation
The degrees of the components of the fibre over a prime, counted with multiplicities, add up to the intersection
number on the generic fibre. This is the conservation of number for the flat family that the arithmetic
subvariety forms over the integers localized at that prime.

% before: cv-origin-intersection-resolution
Integrals over the complex points of an arithmetic subvariety are compared with intersection numbers on a
resolution of these points. The resolution needs to be an isomorphism only off the zero set of an ideal, and it
turns the forms that vanish only there into divisors with normal crossings.

% before: cv-origin-intersection-resolution-exists-arith
The resolution is built for a point with one more block of coordinates. Adding coordinates that are algebraic over the
old ratios does not change the dimension of the kernel.

% before: cv-origin-intersection-resolution-exists-join
Resolution of singularities applies to the point with an extra block of coordinates. It gives the extra block, an ideal
whose zero set contains the bad locus, and a description of the extra coordinates near the good points.

% before: cv-origin-intersection-resolution-exists-embedded-regular
The complex points are a manifold where the local rings of the affine charts are regular of the right dimension. The
condition on the boundary ideal covers each chart by smooth localizations, and one of them contains any given complex
point, which bounds its cotangent space.

% before: cv-origin-intersection-resolution-exists-embedded-snc
The boundary ideal must pull back to a divisor with normal crossings in holomorphic coordinates. Algebraically, this
means a regular system of parameters at each point in which the boundary is a union of coordinate hyperplanes, and in
which a form vanishing only on the boundary is a unit times a monomial.

% before: cv-origin-intersection-resolution-exists-embedded
Where the ring of ratios is covered by smooth charts in which the ideal cuts out coordinate hyperplanes, these charts
glue to a compact complex manifold. It maps onto the complex points, and the ideal pulls back to a divisor with normal
crossings.

% before: cv-origin-intersection-resolution-exists-descent
Forgetting the extra block of coordinates turns a resolution of the joined point into one of the original point. Away
from the zeros of a suitable element the forgotten coordinates are functions of the old ones.

% before: cv-origin-intersection-resolution-exists
Every arithmetic subvariety has such resolutions, off the zero set of a smaller ideal. They come from Hironaka's
resolution of singularities in the form of Lemma~\ref{cv-resolution-coordinates}.

% before: cv-origin-intersection-poincare-lelong
The recursion of an arithmetic intersection system subtracts an integral of the logarithm of the norm of a section
over the complex points. When the section or a metric changes, the Poincar\'e--Lelong formula on these points
accounts for the change of that integral.

% before: cv-origin-intersection-poincare-lelong-all-twist
The Poincar\'e--Lelong formula is available for metrized line bundles invariant under conjugation. A product with a
conic $w_0^2+pw_1^2=0$ has two complex copies of each point, exchanged by conjugation, and a metric that is the given
one on one copy and its conjugate on the other is invariant.

% before: cv-origin-intersection-poincare-lelong-all-prime
The product with the conic stays integral only for suitable primes $p$. A domain finitely generated over $\mathbf Z$
of characteristic zero contains a square root of $-p$ in its fraction field for only finitely many $p$.

% before: cv-origin-intersection-poincare-lelong-all-metric
The twisted metric interpolates between the given one and its conjugate. The weight $\chi_p$ is $1$ at one point of the
conic and $0$ at the other, and conjugation exchanges them.

% before: cv-origin-intersection-poincare-lelong-all-components-prime
Adjoining a square root of $-p$ to the fraction field of $R/W$ gives a field exactly when $W$ is twistable at $p$. The
quotient by the twist then embeds in a polynomial ring over that field, and it is integral over a polynomial ring in
one variable over $R/W$.

% before: cv-origin-intersection-poincare-lelong-all-components
The twist must again be an arithmetic subvariety of the same dimension, and the divisor of a section must twist
component by component. These are the algebraic inputs of the deduction.

% before: cv-origin-intersection-poincare-lelong-all-double-projection
The integrals over the complex points are computed in affine coordinates. A linear projection of the Segre coordinates
that misses the complex points presents them as a finite branched cover of $\mathbf C^e$, and the Monge--Amp\`ere density
is summed over the local holomorphic sheets of this cover.

% before: cv-origin-intersection-poincare-lelong-all-double-good-normalization
A good projection must be linear in the Segre coordinates, so the normalizing elements have to be linear forms. Over an
infinite field a linear change of variables brings any algebraic relation among the forms into monic shape in one of
them, and that form can then be dropped.

% before: cv-origin-intersection-poincare-lelong-all-double-good
Noether normalization of the algebra generated by the Segre monomials gives at most $e+1$ linear forms over which it is
integral, and these forms define a projection that misses the complex points. Its fibres are finite because each Segre
monomial satisfies, on the complex points, an equation whose other terms have lower degree.

% before: cv-origin-intersection-poincare-lelong-all-double-sheets
The sheet density uses a chosen sheet through each point of a fibre. Two sheets through one point differ by a blockwise
rescaling, which changes the potentials by pluriharmonic functions, as long as their Segre images agree, and they agree
near almost every point.

% before: cv-origin-intersection-poincare-lelong-all-double-bridge-ma
The integral through a projection and the integral over a resolution compute the Monge--Amp\`ere density in different
coordinates. Under a holomorphic change of coordinates each complex Hessian is conjugated by the derivative, so the
density changes by the squared modulus of the Jacobian determinant.

% before: cv-origin-intersection-poincare-lelong-all-double-bridge-area
The integral over a compact complex manifold is defined with charts and a partition of unity. A holomorphic map to
$\mathbf C^d$ that is a local biholomorphism off a set of measure zero, with finite fibres, pushes it forward to an
integral over $\mathbf C^d$ of sums over the fibres. On a piece where the map is injective this is the change of
variables formula, and countably many pieces cover the manifold.

% before: cv-origin-intersection-poincare-lelong-all-double-bridge-norm
Points of the fibres of a projection at which a section vanishes must be avoided. A multihomogeneous element outside the
prime satisfies a monic equation over the coordinates of the projection, and its lowest nonzero coefficient gives a
nonzero polynomial on $\mathbf C^{e+1}$ that vanishes on the image of its zero set.

% before: cv-origin-intersection-poincare-lelong-all-double-bridge-map
A good projection, composed with the Segre map and the maps of an analytic resolution, gives a map from the resolution to
$\mathbf C^{e+1}$. Its affine coordinates map the resolution to $\mathbf C^e$, and the rescaled Segre point lies in the
fibre over the image. The points where this fails to be a local biholomorphism, or lie over the boundary, are set aside.

% before: cv-origin-intersection-poincare-lelong-all-double-bridge-fibres
Away from the points set aside, the fibres of the affine projection correspond to the fibres of the good projection.
The correspondence is injective because the resolution meets each orbit at most once off the boundary, and it is onto
for almost every fibre.

% before: cv-origin-intersection-poincare-lelong-all-double-bridge-null
The points set aside, and their images, have measure zero. The ramification would otherwise contain an open set on which
the affine projection has an infinite fibre, which the finiteness of the fibres of the good projection excludes.

% before: cv-origin-intersection-poincare-lelong-all-double-bridge-sheet
At a point of the resolution where the affine projection is a local biholomorphism, its local inverse composed with the
resolution, rescaled in one block, is a sheet of the good projection. Its density is the Monge--Amp\`ere density on the
resolution divided by the squared Jacobian.

% before: cv-origin-intersection-poincare-lelong-all-double-bridge
An analytic resolution turns the integral through a projection into an integral over a compact complex manifold. The
manifold does not depend on the projection.

% before: cv-origin-intersection-poincare-lelong-all-double-independent
Every arithmetic subvariety has an analytic resolution off some ideal. So the integral over its complex points is finite
and does not depend on the good projection that computes it.

% before: cv-origin-intersection-poincare-lelong-all-double-fibre
A complex point of the twist has its last block on the conic $w_0^2+pw_1^2=0$, so it is a multiple of $w_+$ or of $w_-$.
The fibres of a projection for the twist therefore split into two fibres of projections for $X$.

% before: cv-origin-intersection-poincare-lelong-all-double-density
Along each half of the split fibre, the sheets of the twist are the sheets of $X$ with the last block held at $w_\pm$. The
twisted metrics restrict there to the given metrics or to their conjugates.

% before: cv-origin-intersection-poincare-lelong-all-double-conjugate
Complex conjugation maps the complex points of $X$ to themselves, as $X$ is defined over $\mathbf Z$. It carries the
integral for the conjugate metrics to the integral for the given ones.

% before: cv-origin-intersection-poincare-lelong-all-double
The integrals over the twist are the analytic input. Each counts both copies of the complex points, and conjugation
identifies the two contributions.

% before: cv-origin-intersection-poincare-lelong-all
An arithmetic intersection system is defined on all smooth metrized line bundles, which need not be invariant under
conjugation. A twist of the arithmetic subvariety by a conic without real points reduces the formula for them to
the invariant case.

% before: cv-origin-intersection-family-numbers
The numbers are defined through the recursion with chosen sections. Every metrized line bundle is a difference of two
whose multidegrees are nonnegative, and these have monomial sections.

% before: cv-origin-intersection-family-curves-points
On an arithmetic curve the integral over the complex points involves no forms. It is a finite sum over one point in
each orbit, and these points correspond to the embeddings of the function field into $\mathbf C$.

% before: cv-origin-intersection-family-curves-vertical-ring
The vertical terms of the recursion on an arithmetic curve are computed in the local rings of the curve at its closed
points. These rings are built from quotients of multihomogeneous elements of the same multidegree, and the functions
regular at every closed point form an order in the function field.

% before: cv-origin-intersection-family-curves-vertical-local
The local rings of the definition are the homogeneous localizations of the coordinate ring. Their units are read off from
the numerators, and the localization of the coordinate ring at the same prime is a polynomial ring over them, localized.
The second fact lets multiplicities in the coordinate ring be computed in the local ring.

% before: cv-origin-intersection-family-curves-vertical-chain
Several dimension bounds below come from one chain of primes. A homomorphism to a field that kills a multihomogeneous
element outside a prime gives, by switching off the blocks one at a time, a chain of length one more than the number of
blocks above that prime.

% before: cv-origin-intersection-family-curves-vertical-closed
On an arithmetic curve the components of a divisor are closed points lying over a prime number. The dimension formula for
quotients of the polynomial ring over $\mathbf Z$ identifies the relevant minimal primes of a section with the closed points
that contain it.

% before: cv-origin-intersection-family-curves-vertical-component
At a component of the divisor of a section, the local ring is a one-dimensional local domain with a finite residue field
of $p^{\deg(W)}$ elements. The multiplicity of the component is the length of the quotient by the section divided by a
unit, so its vertical term is the logarithm of the order of that quotient.

% before: cv-origin-intersection-family-curves-vertical-order
The functions regular at every closed point form an order $A$ in the function field, whose maximal ideals are the closed
points. The orders of the local quotients by an element of $A$ then multiply to the order of its quotient, which is the
absolute value of its norm.

% before: cv-origin-intersection-family-curves-vertical
The vertical terms of the recursion on a curve are the finite places of its function field. Their sum for the
quotient of two sections is the sum of the logarithms of its absolute values at the embeddings, by the product
formula.

% before: cv-origin-intersection-family-curves
On an arithmetic curve the recursion has no horizontal terms. Its independence of the section is the product formula
for the rational function given by two sections.

% before: cv-origin-intersection-family-product
The right side of the recursion is additive under tensor products, with the product of the sections. This lets the
induction split a line bundle into pieces of nonnegative multidegree.

% before: cv-origin-intersection-family-sections-avoid
The proofs below replace a section by one whose divisor has no component in common with given divisors. Prime
avoidance gives such a section in every large enough multidegree.

% before: cv-origin-intersection-family-sections-degree
The vertical terms of the recursion are intersection numbers on the fibres over the primes. They are multilinear and
symmetric, and they satisfy their own recursion through the divisor of a section.

% before: cv-origin-intersection-family-sections-position
Additivity in the last line bundle in one dimension less gives additivity of the recursion in the bundle before the
last. The integral is additive there, as the Monge--Amp\`ere density is additive in each potential.

% before: cv-origin-intersection-family-sections-cycles-local
Locally at a component of codimension two, the double sum for two sections becomes a sum over the height-one primes of a
two-dimensional local domain. The length of the quotient by the two elements splits into this sum and a kernel that is
the same for both orders, so the sum is symmetric.

% before: cv-origin-intersection-family-sections-cycles
Exchanging two sections in the recursion produces two double sums over components of codimension two. As cycles,
these agree when the two divisors have no common component.

% before: cv-origin-intersection-family-sections-analytic-defect
The analytic part of the exchange of two sections is an identity between integrals that depends on the metrics of the
remaining bundles. Its defect, the difference of the two sides, is a function of these bundles, and the proof studies it
as such.

% before: cv-origin-intersection-family-sections-analytic-additive-slot
An integral over the complex points is an integral against a Monge--Amp\`ere measure on an analytic resolution. That
measure is symmetric in its potentials and additive in each of them, and a tensor product of metrics adds the
potentials.

% before: cv-origin-intersection-family-sections-analytic-additive
Each integral in the defect is a product of first Chern forms integrated against a logarithm. A tensor product in one
place adds the forms, so the defect is additive in each bundle.

% before: cv-origin-intersection-family-sections-analytic-twist
Additivity reduces a smooth bundle to two bundles with plurisubharmonic metrics. A large enough Fubini--Study twist
makes any smooth metric plurisubharmonic, as the Levi form of a torus-invariant function is bounded by that of the
Fubini--Study potential.

% before: cv-origin-intersection-family-sections-analytic-semipositive-regularized
The logarithm of the norm of a section is singular along its zeros. Adding a small constant under the logarithm gives a
smooth metric of multidegree zero that decreases to it.

% before: cv-origin-intersection-family-sections-analytic-semipositive-first
The comparison of the two sides of the defect expands integrals against a tensor product in the first place. The sheet
densities add there, as in the last place.

% before: cv-origin-intersection-family-sections-analytic-semipositive-curvature
The regularized metric need not be semipositive. A fixed Fubini--Study twist makes it so for every value of the
regularizing constant, as its Levi form is bounded below by a multiple, between zero and one, of that of the original
metric.

% before: cv-origin-intersection-family-sections-analytic-semipositive-absolute
Passing to the limit under the integral needs a dominating function. The fibre sums of the absolute values of the sheet
densities provide one.

% before: cv-origin-intersection-family-sections-analytic-semipositive-limit
As the regularizing constant tends to zero, the integrals of the regularized logarithm tend to those of the logarithm of
the norm. This is dominated convergence along the fibres of a good projection.

% before: cv-origin-intersection-family-sections-analytic-semipositive-mono
Against semipositive bundles the integrals are monotone in the integrand. The Monge--Amp\`ere density of positive
semidefinite Hessians is nonnegative.

% before: cv-origin-intersection-family-sections-analytic-semipositive-stokes
For two metrics of multidegree zero, exchanging the integrand and the first form does not change the integral. This is
integration by parts on a resolution of the complex points.

% before: cv-origin-intersection-family-sections-analytic-semipositive
For bundles whose metrics are plurisubharmonic everywhere, the defect has a sign. The proof regularizes the logarithms of
the sections and applies the Poincar\'e--Lelong formula and Stokes's theorem.

% before: cv-origin-intersection-family-sections-analytic
After the cycles, what remains of the exchange of two sections is an identity between integrals over the complex
points. It is the analytic content of the symmetry, and it rests on the Poincar\'e--Lelong formula.

% before: cv-origin-intersection-family-sections-swap
With the cycles and the integrals compared, the recursion can exchange its last line bundle with the bundle whose
section it uses. This needs the recursion in one dimension less, to expand the terms of the components.

% before: cv-origin-intersection-family-sections
In higher dimension the independence of the section uses the recursion in one dimension less. Exchanging two line
bundles in the recursion is where the Poincar\'e--Lelong formula enters.

% before: cv-origin-intersection-family-recursion
With the previous lemmas the defined numbers satisfy the recursion for every section. They are also additive in the
last line bundle, by induction on the dimension.

% before: cv-origin-intersection-family-symmetry
Symmetry is the remaining condition on an arithmetic intersection system. It follows from the recursion with every
section by exchanging adjacent line bundles.

% before: cv-origin-intersection-family
An arithmetic intersection system is built by induction on the dimension through its recursion. The
Poincar\'e--Lelong formula shows that the numbers obtained do not depend on the sections used and are symmetric.

% before: cv-origin-intersection-full
Theorem~\ref{cv-volume-theorem} asks for two further properties of an arithmetic intersection system, for the
projections of a product of projective spaces onto some of its factors. Intersection numbers of pulled-back bundles
are multiplied by the degree of the projection, and they vanish when the image has too small a dimension.

% before: cv-origin-intersection-projection-fibre
The projection formula is a statement about the generic fibre of the projection. It is a projective variety over the
function field of the image, of dimension the difference of the dimensions.

% before: cv-origin-intersection-projection-generic-rational
A line bundle of any multidegree has a rational section, the quotient of two monomials. The recursion holds with such
a section, with the components of both divisors and the integral of the logarithm of the quotient.

% before: cv-origin-intersection-projection-generic-cycles-fibres
Each component of the divisor of a pulled-back section lies over a prime of the image, and its generic fibre over that
prime has a degree. These degrees weight the comparison of the components with those on the image.

% before: cv-origin-intersection-projection-generic-cycles-horizontal
A horizontal component lies over a horizontal prime of the image of at most its dimension. If the dimension drops, the
pulled-back intersection number vanishes by contraction; otherwise the prime is a component on the image, and the
projection formula in lower dimension applies.

% before: cv-origin-intersection-projection-generic-cycles-vertical-degree
The projection formula over a finite field weights the image by the degree of the extension of function fields. For a
vertical component this degree is the degree of its generic fibre over the prime below, computed by the ratio fields as
in characteristic zero.

% before: cv-origin-intersection-projection-generic-cycles-vertical
A vertical component and the prime below it are varieties over the same finite field. Their degrees are compared by the
projection formula over that field.

% before: cv-origin-intersection-projection-generic-cycles-pushforward-situation
The comparison of multiplicities over a component $W$ on the image takes place in two rings. One is the local ring of
the image at $W$; the other is the ring of functions on the source that are regular at every prime over $W$.

% before: cv-origin-intersection-projection-generic-cycles-pushforward-fibres
The primes over $W$ are the components of the pulled-back divisor that lie over $W$. None contains another, since the
fibre over the generic point of the image has dimension zero.

% before: cv-origin-intersection-projection-generic-cycles-pushforward-local
On a variety of any dimension, the multiplicity of a component of a divisor is the length of a quotient of its local
ring. The argument given for curves applies without change.

% before: cv-origin-intersection-projection-generic-cycles-pushforward-order
Over a one-dimensional local domain, the length of the quotient of a torsion-free module by an element is the rank of
the module times the order of the element. This is the form of the projection formula used for the multiplicities.

% before: cv-origin-intersection-projection-generic-cycles-pushforward-integral-centre
Integrality over the local ring is tested on valuation rings. A valuation ring of the function field of a multiprojective
variety has a centre on it, found by dividing each block of coordinates by one coordinate such that all the quotients lie
in the valuation ring.

% before: cv-origin-intersection-projection-generic-cycles-pushforward-integral-centre-fibre
A valuation ring that contains the local ring at $W$ has its centre over $W$ or at the generic point of $X'$. The
principal ideal theorem leaves no prime strictly between $Q$ and $W$, and the elimination relation of the generic fibre
excludes any other centre over the generic point of $Q$.

% before: cv-origin-intersection-projection-generic-cycles-pushforward-integral
The ring of functions regular over $W$ has to be finite over the local ring at $W$. The first step is that it is
integral, which is the valuative form of the properness of the projection.

% before: cv-origin-intersection-projection-generic-cycles-pushforward-closure
The integral elements form a finite module because the local ring is a localization of a ring of finite type over the
integers. Such a ring has finite integral closure in every finite extension of its fraction field.

% before: cv-origin-intersection-projection-generic-cycles-pushforward-finite
Integrality and the finiteness of the integral closure together give a finite module. It is torsion-free, as both rings
lie in fields.

% before: cv-origin-intersection-projection-generic-cycles-pushforward-rank
The rank of this module is the degree of the extension of function fields. By the count of the generic fibre, that is
the degree of the generic fibre.

% before: cv-origin-intersection-projection-generic-cycles-pushforward-max
The maximal ideals of the semilocal ring correspond to the primes over $W$. Localizing at them recovers the local rings
of the source at these primes.

% before: cv-origin-intersection-projection-generic-cycles-pushforward-residue-fibre
The residue degrees will be degrees of generic fibres, so these fibres must be finite. A component above $W$ has the
dimension of $W$, so its generic fibre over $W$ has dimension zero.

% before: cv-origin-intersection-projection-generic-cycles-pushforward-residue
Each residue field of the semilocal ring is a finite extension of the residue field at $W$. Its degree is the degree of
the generic fibre of the corresponding component over $W$.

% before: cv-origin-intersection-projection-generic-cycles-pushforward-length
The length of the quotient by a local equation of the divisor splits over the maximal ideals. Each one contributes its
multiplicity times its residue degree.

% before: cv-origin-intersection-projection-generic-cycles-pushforward
Over a component of the divisor on the image, the components above it are the maximal ideals of a semilocal ring that
is finite and free over the local ring below. Its length modulo the section counts them with their multiplicities and
residue degrees.

% before: cv-origin-intersection-projection-generic-cycles
The components of the divisor of a pulled-back section lie over those of the section on the image. Counted with their
multiplicities and residue degrees, they give the degree of the generic fibre times the terms on the image.

% before: cv-origin-intersection-projection-generic-integral-section-map
An integral over the source is an integral over a resolution of its complex points. To compare it with one over the
image, the resolution is mapped to the base of a good projection for the image.

% before: cv-origin-intersection-projection-generic-integral-section-null-orbits
Away from a hypersurface of $X'$, the points over a given point of $Q$ fall into finitely many torus orbits. The ratios of
the coordinates of $X'$ satisfy equations over $Q$ whose leading coefficients define the hypersurface.

% before: cv-origin-intersection-projection-generic-integral-section-null-isolated
Without the hypersurface, the points of $X'$ over a point of $Q$ may form a positive-dimensional family. Its isolated
points still fall into finitely many orbits, and these are the points that the resolution reaches off its excluded set.

% before: cv-origin-intersection-projection-generic-integral-section-null-zero
A section of a line bundle on $X'$ that does not vanish identically on $X'$ vanishes only on a negligible set of the
resolution. Otherwise the bridge map of $X'$ would carry an open set into the negligible set of values over its zeros.

% before: cv-origin-intersection-projection-generic-integral-section-null-flat
The map to the base of $Q$ cannot have identically vanishing Jacobian on an open set. Otherwise a fibre would contain
infinitely many points of the resolution, while the finiteness of orbits allows only finitely many.

% before: cv-origin-intersection-projection-generic-integral-section-null
The change of variables along this map needs a negligible set of bad points and finite fibres. These are the
properties shown for the affine projection of a resolution in the bridge lemma.

% before: cv-origin-intersection-projection-generic-integral-section-count-orbits
The points of $X'$ over a point of $Q$ off a hypersurface fall into as many torus orbits as the degree of a chart field of
the generic fibre. Read through the good projection, the hypersurface meets the fibres over a negligible set of the base,
so the count holds over almost every point of the base.

% before: cv-origin-intersection-projection-generic-integral-section-count-relation
The ratios of the coordinates of $X'$ in a chart are algebraic over the function field of $Q$, and so is the value there
of a multihomogeneous element. Clearing denominators in its minimal polynomial gives a relation over $R$ whose constant
term is not in $Q$.

% before: cv-origin-intersection-projection-generic-integral-section-count-elimination
The resolution of $X'$ may fail to be injective on the zero set of $J'$, and this set lies over a proper closed subset of
$Q$. An equation for that subset comes from eliminating the ratios of the coordinates of $X'$, which are algebraic over
the function field of $Q$.

% before: cv-origin-intersection-projection-generic-integral-section-count-negligible
Two sets of points of the base are negligible: the values over which a point of $X'$ in the zero set of $J'$ lies, and
the values of the map at its excluded points. Off them, every point of the fibre of $Q$ has the expected preimages in the
resolution.

% before: cv-origin-intersection-projection-generic-integral-section-count
Over a general point of the base, each point of the fibre of the image has as many points of the resolution above it
as the degree of the generic fibre. The specialization of the generic fibre at a general point keeps its degree.

% before: cv-origin-intersection-projection-generic-integral-section-density-sheet
At a point of the resolution off the excluded set, the map to the base has a holomorphic local inverse. Composed with the
resolution and rescaled in one block, it is a sheet of $Q$, and the chain rule for the Monge--Amp\`ere density turns the
density on the resolution into the density along this sheet.

% before: cv-origin-intersection-projection-generic-integral-section-density
At each of these points the density of the pulled-back integrand is the sheet density of the image. The local inverse
of the map is a sheet of the image.

% before: cv-origin-intersection-projection-generic-integral-section
The integral of the logarithm of a quotient of two sections splits into one integral for each section. For a single
section, the integrand and the forms on the source are pulled back from the image.

% before: cv-origin-intersection-projection-generic-integral
The integrand and the forms on the source are pulled back from the image. Over almost every point of the image the
number of points above it is the degree of the generic fibre.

% before: cv-origin-intersection-projection-generic
When the dimensions agree, the generic fibre is a finite set of points. The intersection numbers of pulled-back bundles
are its degree times those on the image.

% before: cv-origin-intersection-projection-count-chart
The coordinate ring of a multihomogeneous prime in a chart is generated by the ratios of the coordinates to the chosen
ones. Its dimension over the base field is the degree of the field these ratios generate.

% before: cv-origin-intersection-projection-count-scaling
The chosen coordinates of each block carry what the ratios forget. For a multihomogeneous kernel they are
algebraically independent over the field of the ratios.

% before: cv-origin-intersection-projection-count-disjoint
Adjoining algebraically independent elements to two fields does not change the degree of one over the other. This
removes the coordinates of the image blocks from the comparison of degrees.

% before: cv-origin-intersection-projection-count
The degree of the generic fibre is a degree of fields. In a chart its coordinate ring is the extension of the fields of
ratios of the two points.

% before: cv-origin-intersection-projection
The projection formula holds for every arithmetic intersection system. The generic fibre of the projection is a
finite set of points, whose number, counted with multiplicities, is the degree of the extension of function fields.

% before: cv-origin-intersection-contraction-image
The contraction property compares an arithmetic subvariety with its image under a projection that forgets blocks of
coordinates. The arguments below need the image to be an arithmetic subvariety of no larger dimension, and its
dimension to drop when a function pulled back from it vanishes.

% before: cv-origin-intersection-contraction-ratios
The integrals in the recursion vanish because the image has small dimension. A variety of dimension $d$ has a function
field of transcendence degree $d-1$, so any $d$ ratios of coordinates satisfy a polynomial relation on its complex
points.

% before: cv-origin-intersection-contraction-rank
Along a holomorphic map into the complex points of a variety of small dimension, such a relation forces some
directions to move the image only along the orbits of the torus. Their number is bounded below in terms of the
dimension.

% before: cv-origin-intersection-contraction-euler
A metric on $\mathcal O(a)$ changes along an orbit of the torus by the factor $\prod_k|c_k|^{a_k}$. So the logarithm of
a pulled-back metric is pluriharmonic in the directions of the previous lemma, and its complex Hessian kills them.

% before: cv-origin-intersection-contraction-discriminant
The Monge--Amp\`ere density is an alternating sum of determinants of sums of complex Hessians. It vanishes when enough
of the Hessians kill a large enough subspace.

% before: cv-origin-intersection-contraction-forms
The three previous lemmas combine into the vanishing of the integrals in the recursion. First Chern forms pulled back
from the image vanish in a product with more factors than the complex dimension of the image.

% before: cv-origin-intersection-contraction-field
The vertical terms of the recursion are intersection numbers on varieties over finite fields. Over a field the same
vanishing holds, by induction through the divisors of sections pulled back from the image.

% before: cv-origin-intersection-contraction-vertical
On a vertical component the image lies in a special fibre of the image of the arithmetic subvariety, of smaller
dimension. So its degrees vanish when enough of the bundles are pulled back.

% before: cv-origin-intersection-contraction-pulled
The case in which every line bundle is pulled back comes first. A section pulled back from the image keeps the
horizontal components of its divisor in the same situation, with images of smaller dimension, so the recursion gives
an induction on the dimension.

% before: cv-origin-intersection-contraction
The contraction property also holds for every arithmetic intersection system. In its recursion, the terms that
remain involve line bundles pulled back from an image whose dimension is too small to carry them.

% before: cv-origin-intersection-exists
Together these give a full arithmetic intersection system. Theorem~\ref{cv-volume-theorem} takes one as given, and
its recursion compares integrals over the complex points with intersection numbers on the fibres.

% before: cv-origin-weight-nef-generic
An integral of first Chern forms over the complex points is an algebraic intersection number on the generic
fibre. There the results on nef line bundles over a field apply.

% before: cv-origin-weight-nef-kodaira
A nef line bundle with positive top self-intersection on a projective variety is big. A multiple of it twisted
down by the ample bundle $\mathcal O(1,\dots,1)$ then has a nonzero section.

% before: cv-origin-weight-nef-sections
The sections needed are integral: on every chart they lie in the ring generated over $\mathbf Z$ by the ratios.
A rational section becomes one after its denominators on the finitely many charts are cleared.

% before: cv-origin-weight-nef-shift
The boundary form lowers the multidegree of the sections by its own multidegree. Raising a section to a power
and multiplying it by coordinates absorbs any such twist.

% before: cv-origin-weight-nef
Integral sections of large multiples of a combination of the bundles of a realization exist when the degree
polynomial of the analytic system is positive there. The combination may be twisted down by any multihomogeneous
form, and this is where the boundary form enters.

% before: cv-origin-weight-gap-uniformized
The positivity of the degree polynomial comes from a proportionality of Hirzebruch--Mumford type, which needs the
analytic system to be uniformized by a flag domain. The period charts of the uniformization come from the periods of the abelian
varieties over the complex points of the reference geometry.

% before: cv-origin-weight-gap
The weight $\nu$ enters through the gaps between its consecutive entries. The degree polynomial of the analytic
system is positive at these gaps when $\nu$ decreases strictly.

% before: cv-origin-weight-data
An integral section of the right multidegree and the boundary form combine into the effective section of the
origin. The computation is local on a chart, using the generators of the root lines.

% before: cv-origin-weight
The effective section of a good origin exists on every reference geometry with analytic data. It is the product of
the section given by positivity, the boundary form and the inverse generators of the root lines.

% before: cv-origin-reference-closure
The elements $y^c_i$ of the data at a component must generate its field. The block $k_0$ embeds the generic fibre,
in particular at the generic point.

% before: cv-origin-orbit-forms
The coordinates of the block $k_0$ determine those of the other blocks through homogeneous forms. At a point of the
zero set of the kernel, the embedding of the generic fibre by this block gives such forms that do not all vanish there.

% before: cv-origin-reference-orbit
The map $\mathrm{pt}_c$ must be injective up to scaling. Points of the arithmetic variety with proportional
coordinates in the block $k_0$ are equal.

% before: cv-origin-ample-families
The forms relating the other blocks to the block $k_0$ were found point by point. For ampleness finitely many families
of them must serve every point of the generic fibre at once, and valuation rings give such families.

% before: cv-origin-reference-ample
Condition (3e) of a good origin asks that the bundle $\bar A$ be ample on the generic fibre. Its multidegree is
concentrated on the block $k_0$, which embeds the generic fibre.

% before: cv-origin-reference-constants
The field embedding of the data at a component must carry algebraic numbers into the coefficient field. A flag of
degree $1$ adds only independent transcendental elements to the function field of the component.

% before: cv-origin-reference-immersion
Condition (3d) asks that the coordinates of the block $k_0$ give an immersion of the analytic system into
projective space. On the manifold of complex points this follows from the immersion of the lifts.

% before: cv-origin-analytic-inverse
The immersion of Lemma~\ref{cv-origin-reference-immersion} is on the manifold of complex points, while (3d) is a
statement on the analytic system. A holomorphic inverse of $\pi$ transports it.

% before: cv-origin-extension-tate
A Tate system at a geometric point must survive the passage to a larger algebraically closed field. The torsion points
do not grow, as the scheme of points of order dividing $m$ has finitely many points over the smaller field and is of
finite type there.

% before: cv-origin-extension-level
The level cover starts from a full level structure over a finite extension of the residue field. The torsion points of
a Tate system over an algebraic closure are defined over a finite extension, as there are finitely many of them modulo
$N$ and each is a point of a scheme of finite type.

% before: cv-origin-extension-charts
A scheme in projective space whose local rings embed in the function field has coordinates given by the point of the
function field. Its charts are the inverse images of the standard affine opens.

% before: cv-origin-extension-normal
The level cover needs a finite normal scheme over the model with a point in the finite extension. It is the
normalization, in that extension, of the closure of the point in projective space; its local rings are integrally
closed, and its projective embedding gives coordinates.

% before: cv-origin-extension-cover
The points of finite order of $\mathcal A$ at a generic point are made rational by a level cover. Over a finite
extension of the residue field the abelian variety receives a full level structure, and a finite normal scheme over
$\mathcal X$ carries the pulled back semi-abelian scheme with this structure.

% before: cv-origin-change-extension
At the primes of $S'$ the Hodge component must carry the Hodge--Tate periods of its points of order $p^4$. These
points and the corresponding roots of unity are defined over a finite extension of the function field.

% before: cv-origin-coords-join-map
After a finite extension of the field, the new scheme carries a point with values in the larger field, and the map to
the old scheme must be compatible with the embedding of fields. The map is built chart by chart, as for a single field,
with the embedding applied to the ratios of the old coordinates.

% before: cv-origin-coords-join
A finite extension of the field is reached by adjoining a block of coordinates whose ratios generate it. The scheme
with the joined coordinates exists because the old coordinates already give the smaller field by fractions, and it maps
to the old scheme compatibly with the embedding of fields.

% before: cv-origin-coords-normal
The new coordinates must have normal local rings. On each chart the integral closure of the ring of ratios is a finite
module, and a common choice of generators for all charts gives a family whose charts have these integral closures as
rings of ratios.

% before: cv-origin-coords-crossings
Condition (3) of a reference geometry asks, on each chart, for smooth localizations of the ring of ratios with
coordinates in which the boundary is cut out by some of the coordinates. The same condition for a single element is the
form in which the resolution of singularities delivers it.

% before: cv-origin-coords-snc-chart
The resolution gives a smooth scheme covered by charts in which the boundary is a union of coordinate hyperplanes. On an
affine open containing the generic point this becomes a condition on the ring of values, after a cover by open sets that
are standard in both the affine open and the charts.

% before: cv-origin-coords-snc
The generic fibre of the new coordinates must be smooth, with a boundary ideal of simple normal crossings type. Both
come from the resolution of singularities in characteristic zero, applied to the closure of the point in projective
space.

% before: cv-origin-change-coords
Over the finite extension the component needs new coordinates with normal local rings. Their generic fibre is
resolved, as the analytic resolution of the next construction requires.

% before: cv-origin-change-factor
The construction of a Hodge component at $S'$ starts from a point of the moduli scheme. The generic points of the
model lie in the image of the moduli scheme.

% before: cv-origin-change-projective-periods
A Hodge component at $S'$ is cut out at the primes of $S'$ by Hodge--Tate periods of its semi-abelian scheme. They make
sense for a commutative group scheme with a line bundle over a scheme with a generic point, once its points of order
dividing $n$ are defined over the function field and carry a Weil pairing there.

% before: cv-origin-change-projective-weil-pullback-characters
A generic Weil pairing pulled back along an embedding of fields must stay perfect, so every character of the points of
order dividing $n$ with values in the larger field has to come from the smaller one. Counting characters shows this for
every finite group whose characters with values in the smaller field separate its points.

% before: cv-origin-change-projective-weil-pullback
The semi-abelian schemes of the construction are pulled back along morphisms whose generic points embed the function
fields. Full torsion and generic Weil pairings pass to such pullbacks, with the pairing composed with the embedding, since
the points of the pullback with values in a field are those of the original scheme.

% before: cv-origin-change-projective-weil-pairing
A generic Weil pairing exists wherever the points of order dividing $n$ are those of an object at level $l$. The Tate
systems of the object at the geometric points describe these points with their commutator pairings, which gives the
pairing over the function field once the torsion and the roots of unity are defined there.

% before: cv-origin-change-projective-weil
The semi-abelian schemes of the construction are pulled back from the model, and their generic points lie over the
moduli scheme, whose abelian scheme is an object at level $l$. So the model has a generic Weil pairing at these points,
and it pulls back to the semi-abelian schemes.

% before: cv-origin-change-projective-local-torsion
Near a point, Hodge--Tate periods come from a flat and quasi-finite group scheme over an affine chart whose points over
the generic point are the points of order dividing $n$. A function on it with the values of a character at these points
has a derivative modulo $n$ along the Lie algebra, and this derivative is the Hodge--Tate form of the character.

% before: cv-origin-change-projective-local-chart
The construction takes place over an affine neighbourhood of the point. One can be chosen whose ring is a noetherian
integrally closed domain with fraction field $\Omega$, as the local rings at the specializations of the generic point are
integrally closed in $\Omega$.

% before: cv-origin-change-projective-local-inputs-multiplication
The torsion schemes of the construction are kernels of multiplication by $N$ on semi-abelian schemes. Such a kernel is flat
and quasi-finite over the base when multiplication by $N$ is, and this is checked on the fibres, where an extension of an
abelian variety by a torus has only finitely many points of order dividing $N$.

% before: cv-origin-change-projective-local-inputs-torsion
Over a chart whose generic point carries all the points of order dividing $N$, the kernel of multiplication by $N$ is a
torsion scheme. Its generic fibre is \'etale, as $N$ is invertible in characteristic zero, and its points over the generic
point are the given points of order dividing $N$.

% before: cv-origin-change-projective-local-inputs-zero
A torsion scheme is flat over a domain, and its generic fibre is reduced with the points $\gamma_t$ as its only points. So a
function on an open subset of it is determined by its values at the $\gamma_t$, and functions on it can be compared there.

% before: cv-origin-change-projective-local-inputs-multiplicative
A Hodge--Tate form is read off from a function on a torsion scheme with the values of a character at the points
$\gamma_t$. Such a function is multiplicative on all points, as the values at pairs of the $\gamma_t$ determine a function
on the product; this makes its derivative along the Lie algebra additive.

% before: cv-origin-change-projective-local-inputs-hopf
Restricted to the points $\gamma_t$, the functions on a torsion scheme of type $C$ give an algebra of functions on $C$
with values in $R$. A function with the values of a character exists when the character lies in this algebra, and the
properties of a Hopf order are what put the characters there.

% before: cv-origin-change-projective-local-inputs-hopf-characters
The dual of a Hopf order is a finitely generated ring under convolution, and pairing it with a character is a ring
homomorphism. So the pairings are integral over $R$, which, for $R$ integrally closed, puts the character into the order.

% before: cv-origin-change-projective-local-inputs-hopf-order-subgroup
Over a Henselian local ring, the part of a torsion scheme that is finite over the base is an open subgroup, and its ring of
functions is a finite algebra. The proof that the functions on a torsion scheme form a Hopf order compares the functions on
such an open subgroup with those on the whole torsion scheme.

% before: cv-origin-change-projective-local-inputs-hopf-order-glue
A function on an open subset of a torsion scheme is determined by its values at the points of the generic fibre, so
functions with prescribed values on the sets of an open cover glue. On a quasi-compact open set every vector of values is
attained after multiplication by a nonzero element of the ring, because distinct points give linearly independent
homomorphisms.

% before: cv-origin-change-projective-local-inputs-hopf-order-induced
A Hopf order on a subgroup induces one on the whole group: a function belongs to it when each of its translates restricts
to an element of the order. The functions on a finite open subgroup will give a Hopf order on the subgroup of labels whose
points it contains.

% before: cv-origin-change-projective-local-inputs-hopf-order-finite
On a finite open subgroup, the ring of functions is a finite projective algebra over the base, and the values at the
generic points embed it into the functions on a subgroup of labels. The group law gives the comultiplication and the
inverse, so this ring is a Hopf order there, and the order it induces is the one defined by the open subgroup.

% before: cv-origin-change-projective-local-inputs-hopf-order-closed-fibre
Over a local ring, an open set that contains the closed fibre receives every morphism from a scheme finite over the base,
as finite morphisms are closed. Over a Henselian local ring, the closed fibre has an open neighbourhood finite over the
base, by the structure of quasi-finite algebras over Henselian rings.

% before: cv-origin-change-projective-local-inputs-hopf-order-base-change
The comparison of functions passes through base changes of the torsion scheme to Henselian rings. The generic fibre of a
torsion scheme is a disjoint union of copies of the spectrum of the fraction field, so base change keeps the conditions on a
torsion scheme, and along a flat base change a function with prescribed values descends.

% before: cv-origin-change-projective-local-inputs-hopf-order-translate
Two finite open subgroups, one inside the other, give the same functions. A point of the larger one is moved into the
smaller one by translation by a section over the base, which extends a point of the generic fibre because the base is
integrally closed.

% before: cv-origin-change-projective-local-inputs-hopf-order-henselization
To compare functions near a point of the torsion scheme, the base is replaced by a Henselian ring over the prime under
that point. The henselization of the local ring at the prime serves, as it is flat and stays an integrally closed domain.

% before: cv-origin-change-projective-local-inputs-hopf-order-all
Over an integrally closed domain, the functions on a finite open subgroup with the prescribed values extend to the whole
torsion scheme. Near a point, after base change to a Henselian ring, the inverse image of the subgroup lies in the finite
part of the base change, which has the same functions, and the extension descends.

% before: cv-origin-change-projective-local-inputs-hopf-order-flat
The functions on a torsion scheme commute with flat base change of the ring. This follows from the flat base change of
global sections, as the functions in question form the kernel of a map between modules built from those sections.

% before: cv-origin-change-projective-local-inputs-hopf-order-local
Being a Hopf order can be checked after flat base change to local rings over the maximal ideals, once the order has full
rank. Projectivity comes from flatness by faithfully flat descent, and the comultiplication and the inverse descend
because membership in a submodule does.

% before: cv-origin-change-projective-local-inputs-hopf-order
The functions on a torsion scheme over a noetherian integrally closed domain form a Hopf order. The proof passes to the
fraction field and to flat Henselian integrally closed local extensions of the ring, over which the torsion scheme is base
changed.

% before: cv-origin-change-projective-local-inputs-character
A Hodge--Tate form of a character needs a function on the torsion scheme with the values of the character at the points
$\gamma_t$. It exists because the character lies in the Hopf order of the functions on the torsion scheme.

% before: cv-origin-change-projective-local-inputs-tangent
A tangent vector of the group scheme over the chart, reduced modulo $N$, is a point of order dividing $N$, as $N$ times it
vanishes modulo $N$. So it lifts, in one way only, to the torsion scheme, where a function can be evaluated on it.

% before: cv-origin-change-projective-local-inputs-forms
A Hodge--Tate form of a character is the derivative modulo $N$ of a function with its values, taken along the lifts of
tangent vectors to the torsion scheme. The function is multiplicative, so this derivative is linear, and it lifts to a
linear form, as the Lie algebra of a smooth group scheme is projective.

% before: cv-origin-change-projective-local-inputs-chart
On an affine chart, a Hodge--Tate form of the character given by the generic Weil pairing yields a Hodge--Tate period at
every point of the chart. The function on the torsion scheme extends to a neighbourhood of the unit section, as the torsion
scheme is closed in the group scheme over the chart, and its derivative computes the form.

% before: cv-origin-change-projective-local-inputs-unique
The Hodge--Tate periods of a point are determined up to $n$ times the dual lattice. Two functions with the values of the
same Weil pairing agree on the torsion scheme over the local ring, so their derivatives along the Lie algebra agree modulo
$n$.

% before: cv-origin-change-projective-local
The modules spanned by the Hodge--Tate periods are to be made into vector bundles, which needs them to be locally finitely
generated. This follows from a local statement: near every point the periods of all points of order dividing $n$ can be
chosen at once, and those of one point differ by $n$ times forms in the dual of the Lie algebra.

% before: cv-origin-change-projective-endomorphism
An endomorphism of the group scheme moves the points of order dividing $n$, and through the generic Weil pairing it moves
their Hodge--Tate periods by its differential. This makes the modules of periods stable under the action of
$\mathcal O_E$, which acts by endomorphisms.

% before: cv-origin-change-projective-coherent
A modification of the model turns locally finitely generated lattice families into vector bundles. The modules spanned by
the Hodge--Tate periods are such families, and they lie between the dual Lie algebras and $n$ times them, so they agree
with the dual Lie algebras away from the primes dividing $n$.

% before: cv-origin-change-projective-pullback-transport
Over the pullback the Hodge--Tate periods are taken for a second generic Weil pairing, which need not be the pullback of
the first. Two such pairings on groups in bijection still match point by point: each point on one side has a partner on
the other side with the same pairings, up to the embedding of fields.

% before: cv-origin-change-projective-pullback-derivative
Condition (b) of a Hodge--Tate period reads off the coefficient of $\varepsilon$ of a function along tangent vectors. This
coefficient is linear on the Lie algebra, since the group law on points of first order is addition.

% before: cv-origin-change-projective-pullback-extend
To check condition (b) at a point it suffices to check it on generators of the Lie algebra at that point. A vector that
extends to the point is a combination of the generators with coefficients in the local ring, and the derivative of the
function follows the same combination.

% before: cv-origin-change-projective-pullback-lie
The Lie algebras of a smooth commutative group scheme and of its pullback are compared at every point as in
Lemma~\ref{cv-origin-reference-lie-pull}. The same bases serve both, so the dual lattices pull back as the lattices do.

% before: cv-origin-change-projective-pullback-situation
The comparison of Hodge--Tate periods along a pullback takes two steps with the same data. The pullback carries its own
line bundle and generic Weil pairing, and the maps $j$ and $j^\vee$ transport vectors and forms.

% before: cv-origin-change-projective-pullback-torsion
Condition (a) of a Hodge--Tate period concerns the values of a function at the points of order dividing $n$. For points
matched through the two generic Weil pairings it pulls back, as both commutator pairings are read off from the generic
pairings.

% before: cv-origin-change-projective-pullback-periods
With both conditions, a Hodge--Tate period pulls back to one of the matched point. Condition (b) is checked on the bases
of Lemma~\ref{cv-origin-change-projective-pullback-lie}, along which the function and the forms change by the local
homomorphism of the morphism.

% before: cv-origin-change-projective-pullback
After the modification the semi-abelian scheme is replaced by its pullback, whose Hodge--Tate periods must be the
pullbacks of the old ones. Then the vector bundles that the modification produces are those of the new semi-abelian
scheme.

% before: cv-origin-change-projective-bundles-weil
At the points over no prime the modules $\mathcal V_s$ are compared with the Hodge--Tate periods, which needs a generic
Weil pairing. It exists once the torsion of the semi-abelian scheme itself is defined over the function field, since its
pullback to the generic point is an object at level $l$.

% before: cv-origin-change-projective-bundles-predual
At the primes of $S$ the modules $\mathcal V_s$ consist of the vectors of an eigenspace at which the Hodge--Tate periods
take regular values. They form a vector bundle when the restricted periods do, through dual bases.

% before: cv-origin-change-projective-bundles-saturate
For the labels $\tau$ the modules are further saturated by the different of $E$. A nonzero ideal of a Dedekind domain
becomes principal locally, so saturating a vector bundle multiplies its local bases by one scalar.

% before: cv-origin-change-projective-bundles-generic
At the points over no prime the integers $p^{R_p}$ are units, and the Hodge--Tate periods span the dual of the Lie algebra.
There the modules they cut out, saturated or not, are the eigenspaces of the Lie algebra.

% before: cv-origin-change-projective-bundles-glue
The modules $\mathcal V_s$ are glued from a vector bundle on the points over no prime of $S$ and vector bundles near each
prime of $S$, which agree at the points over no prime. On a noetherian scheme two bases that agree at the points over no
prime agree away from finitely many primes, and this allows the gluing.

% before: cv-origin-change-projective-bundles
The modules $\mathcal V_s$ of a Hodge component are cut out by the Hodge--Tate periods at the primes of $S'$ and by the
Lie algebra elsewhere. They form vector bundles once the restrictions of the periods to the eigenspaces of
$\mathcal O_E$ do.

% before: cv-origin-change-projective-frames
The modification must not change the model over the generic fibre. There the Hodge--Tate periods are the dual of the
Lie algebra, and its restrictions to the eigenspaces of $\mathcal O_E$ are already free, so the modification is an
isomorphism there.

% before: cv-origin-change-projective
A Hodge component at $S'$ needs the vector bundles of Hodge--Tate periods at the primes of $S'$. They are built
over the normal model of the finite extension, without changing it over the generic fibre.

% before: cv-origin-change-comparison
The flag of the reference geometry must be carried to the new Hodge component. The Lie algebras of the two
components are compared along the morphism between them, and a generic flag of degree $1$ becomes one of the
degree of the extension.

% before: cv-origin-change-bases
The transported bases are generic, but the new component may not yet have root lines with generators on every
chart. A further Hodge component over it supplies them, and it agrees with the old one where the bases are
already adapted.

% before: cv-origin-change-smooth
The new coordinates must keep a smooth generic fibre, for the analytic resolution of the next construction.
Smoothness passes along maps that are isomorphisms over the generic fibre.

% before: cv-origin-change
The data at a finite set of primes start from a change of the reference geometry. The Hodge component moves to a
finite cover over which the points of order dividing $p^4$ are defined, for the primes $p$ of $S'$.

% before: cv-origin-model-resolution-coords
The complex points of the new coordinates are made into a manifold one affine chart at a time. Near a point with
regular local ring, a system of parameters is a system of holomorphic coordinates on the zeros of the ideal.

% before: cv-origin-model-resolution-regular
The affine charts have regular local rings at their complex points because the generic fibre is smooth. The
dimension comes from the arithmetic variety, and smoothness bounds the tangent spaces.

% before: cv-origin-model-resolution
On the new coordinates the complex points must form a compact manifold with local lifts, as for the reference
geometry. Smoothness of the generic fibre gives such a manifold with no exceptional locus.

% before: cv-origin-lift-situation
The degree of the map between the two manifolds is computed from the coordinates. Its fibres are the points of one
arithmetic variety over a point of another, whose blocks are among those of the first after renaming, and they are
counted up to the action of the torus on each block.

% before: cv-origin-lift-fibre
The count over a general point is the degree of a field extension. A generic fibre puts the coordinates of the two
varieties in a pair of fields, and the ratios in the blocks that are not renamed generate the extension whose degree
is the count.

% before: cv-origin-lift-scaling
The degree of the generic fibre is compared with the degree of the function fields by scaling each block by one of
its coordinates. For a variety cut out by multihomogeneous equations, these scaling coordinates satisfy no relation
over the field generated by the ratios.

% before: cv-origin-lift-degree
For the two points of Lemma~\ref{cv-origin-model-lift}(2), the generic fibre is built from the cone over the larger
point. Its degree is the degree of the extension of function fields, because the scaling coordinates do not change
degrees.

% before: cv-origin-lift-finite
The count needs a finite extension. When the two varieties have the same dimension, each ratio of the new blocks is
algebraic over the smaller function field, by a comparison of transcendence degrees.

% before: cv-origin-lift-spec
In an affine chart the count is the specialization of a finite extension at a point. A primitive element turns the
points of the fibre into the roots of one separable polynomial, once finitely many denominators are inverted.

% before: cv-origin-lift-chart
Normalizing each orbit by the chosen coordinates of the new blocks turns the orbits over a point into the points of an
affine fibre. The general count then holds in one chart, away from the zeros of one form.

% before: cv-origin-lift-orbits
The charts together cover every orbit over a general point, and they all have the same chart field. So the count in
one chart is the count of all orbits, away from the zeros of the product of the forms of the charts.

% before: cv-origin-lift-dense
The exceptional set of the map must have empty interior. This rests on the density of the complement of a hypersurface
in the complex points of a variety defined over the rationals.

% before: cv-origin-lift-nonvanishing
An arithmetic variety is defined over the integers, and its complex points lie in the product of punctured spaces.
The density passes to these points because the ideal stays prime over the rationals and the product is open.

% before: cv-origin-lift-interior
On a manifold that resolves an arithmetic variety, the zeros of a form that does not vanish on the variety have empty
interior. A point of the interior would be a limit of points where the form does not vanish, and the
injectivity of the lifts would put one of these points in the interior.

% before: cv-origin-model-lift
The new manifold maps to the manifold of the reference geometry, compatibly with the coordinates. The map has the
degree of the extension of function fields, and so has its composite with $\pi$.

% before: cv-origin-model-pullback
Regular functions, units and values at complex points must be compared along the morphism of Hodge components.
These comparisons follow from the local homomorphisms of local rings along it.

% before: cv-origin-model-hodge
The Hodge metrics on the new root lines are obtained from the old ones. The Lie algebras, the graded pieces of
the flags and their Hodge norms correspond along the morphism of Hodge components.

% before: cv-origin-model-metric-norm
At a point over which the semi-abelian scheme is abelian, one generator of a graded piece fixes the Hodge norms of
all its elements. Another element has the norm of the generator times the absolute value of its ratio to it. So
the norms on the two Hodge components can be compared through a single generator pulled back from the old one.

% before: cv-origin-model-metric-chart
A Hodge metric is defined by an identity in every chart of the coordinates in which the point has no zero
coordinate. The identity in one such chart implies it in the others, since the ratios of the chosen coordinates are
regular at the point with the ratios of the complex coordinates as values.

% before: cv-origin-model-metric
The new root lines have different generators, so the pulled-back Hodge metrics need a correction. With it, the
Hodge norms of the generators agree.

% before: cv-origin-correction-extend
The correction is built from a bundle of multidegree zero whose logarithm on the manifold of the analytic data is a
given smooth function. The blocks of a point give rank-one projections that do not change along orbits. Through them
the manifold embeds in a vector space, where the function extends by local inverses and a partition of unity.

% before: cv-origin-model-correction
The correction is a smooth metrized line bundle given by the ratio of the old and new generators. Its potentials
are pluriharmonic, so it does not change the first Chern forms of the realization.

% before: cv-origin-model-section
The effective section of the reference geometry must remain effective on the new Hodge component. Generators of
the root lines and units correspond along the morphism, which gives the transfer.

% before: cv-origin-model
The flag model at $S'$ is the pullback of the reference flag model to the new coordinates, corrected on the root
lines. Its realization is the new manifold, mapped to the analytic system through the reference one.

% before: cv-origin-data
The data of a good origin split into data at one component, which do not depend on the level, and data
at one finite set of primes. The first come with a reference model of degree $1$ realized bijectively on
the analytic system, and the second are built from them for any finite set of primes that contains the
primes dividing the discriminant.

% before: cv-volume-good-tower
The volume theorem assumes a tower with a good origin. Such a tower exists on every global model
whose components are defined over the coefficient field, so the volume theorem applies to the moduli
schemes of the deep covers.

% before: cv-point
The covolume bound is uniform over the points of the families of definite-side data built
in Chapter~\ref{chapter-assembly}. Its statement and the steps of its proof all refer to
the same lattices, weights and scale at a point.

% before: cv-covolume-bound
The comparison inequality needs, for each of the $m$ indefinite groups of the chain, an
upper bound for the covolume of the lattice of lifted forms. The main term below is the
contribution of the period $\Omega$; the error grows like $e^{x^{\varepsilon}}$, which is
the scale $b$ of the windows.

% before: cv-separated
The level group is a product of general linear groups over the finite fields at the level places,
and the lifted forms are cut by an idempotent of its group algebra. Green's parametrization labels
each irreducible character by a polynomial, and the idempotent keeps the characters whose labels
are squarefree, have well separated roots and are far from the tame values.

% before: cv-bad-part
The part of the lifted forms removed by the separated idempotent must have small rank. When the
representation is a sum of regular representations of large subgroups, that rank is at most the
share of the labels that are not separated, measured by the Gelfand--Graev idempotents.

% before: cv-squarefree-mass
The weights that Green's parametrization gives to squarefree labels must account for almost all of
the total weight. Counting squarefree polynomials gives this up to an error of order $1/g$.

% before: cv-local-estimates
Each condition in the definition of a separated tuple fails for a small proportion of labels,
counted with their weights. The estimates below bound these proportions one condition at a time,
for a single polynomial.

% before: cv-property-w
The share of the labels that are not separated is a probability over random tuples of
polynomials, with the weights that Green's parametrization gives to squarefree labels. Each
condition of separation fails with small probability, and the bound below adds these
probabilities.

% before: cv-local-share
At the points of the families the level primes are very large compared with the ranks, and the
auxiliary primes of the host are large too. The bound for the share then decays faster than any
power of $x$.

% before: cv-lifting-map
The comparison with automorphic forms in the next proposition passes through a map from cusp forms
on the Shimura variety of $U(W_j)$ to lifted forms, linear only over the coefficient field of the
Shimura variety. The lifting of forms identifies the cusp forms with a rational structure on the
lifted forms and matches the norms at the finite places, and the properties of the map follow from
these two facts alone.

% before: cv-lifting-rational
The lifted forms of a given type are cut out of all lifted forms by eigenvalue equations whose
coefficients lie in the smaller field. Two facts about extending scalars from a number field $K$ to
$L$ control such subspaces: their solutions are spanned by solutions over $K$, and the
$\mathcal O_L$-span of an $\mathcal O_K$-module meets a subspace defined over $K$ in the span of the
corresponding $\mathcal O_K$-module.

% before: cv-lifting-type
When the type is defined over the smaller field, the lattice of integral lifted forms of that type
is generated by integral rational forms of the type. These are the lifts of integral cusp forms of
the type, which is how the lattice of lifted forms of the prescribed type is compared with integral
cusp forms in the next proposition.

% before: cv-realization-linear
The bad-set condition compares ranks of lattices of lifted forms of a given type, while
Lemma~\ref{cv-bad-part} counts dimensions in a complex representation. Extending scalars from the
coefficient field preserves the dimension of a subspace cut out by a character of a subgroup and by
the kernel of an element of the group algebra. The second part builds the complex space from forms
on a cover, which is possible when the forms invariant under the kernel of the covering group come
from below.

% before: cv-realization-split
On the definite group the Shimura variety has dimension zero, and its forms are a module over a
finite product of separable field extensions of the coefficient field. Over $\mathbf C$ this module
splits along the complex points, and a group acting freely on the points permutes the pieces
freely. The proof of Proposition~\ref{cv-point-conditions} deduces from this the hypothesis of
Lemma~\ref{cv-bad-part} with $c=1$ for the definite group.

% before: cv-regular-group-determinant
On the indefinite groups the embedding into regular representations comes from multiplying by one
form whose translates are independent in a strong sense: its group determinant is not a zero
divisor. Multiplying by the translates embeds the forms of one weight into regular representations
and bounds an eigenspace of a character from below by the forms of a smaller weight. The two bounds
are compared through the ratio of the dimensions of the spaces of forms of two weights, which gives
the factor $3/2$ in Proposition~\ref{cv-point-conditions}.

% before: cv-cusp-setting
The bound for the maximal slope in Proposition~\ref{cv-point-conditions}(4) needs two facts on the
integral cusp forms on the Shimura variety of $U(W_j)$: a lower bound for their Petersson norms and an
upper bound for the dimension of the space they span. Both come from comparing the supremum of a
cusp form with its Petersson norm. The setting fixes the readings of forms in the uniformizations,
these two norms, and the explicit constants that enter the bounds.

% before: cv-decay-majorant
The local bound for a cusp form is proved in the charts of one component. It uses three objects attached
to a frame: a positive definite form that measures lattice vectors at the frame, the unipotent
transvections along an isotropic vector, and the factor relating the value of a form to its pointwise
norm.

% before: cv-decay-flow
Moving a frame along the transvections of an isotropic vector multiplies its leading minors by explicit
factors. These factors decide when the moved family is still a frame, and how far the pointwise norm of a
form can move from the absolute value of the form.

% before: cv-decay-fourier
Along an isotropic vector, a form becomes a periodic holomorphic function of the transvection parameters,
with a Fourier expansion. A bounded pointwise norm forces the frequencies to be nonnegative, and for a
cusp form the constant term vanishes, so the form decays as the parameter goes to infinity.

% before: cv-decay-schwarz
A periodic holomorphic function on a half-plane that tends to zero at infinity is small at a point when it
is bounded on a line below that point. This is the Schwarz lemma after the change of variable to the
punctured disc.

% before: cv-decay-maximum
At a point where a cusp form attains at least half of its supremum, the lattice $L'$ has no short vectors
for the majorant. A short vector would be isotropic, and the decay along it together with the Schwarz lemma
would make the form at that point smaller than the choice of the point allows.

% before: cv-submean-chart
Near the centre of the chart at an orthonormal frame, the normalizing factor of a form and the density of
the invariant measure are close to $1$. This turns the sub-mean value property of holomorphic functions on
a ball into a bound for the pointwise norm.

% before: cv-returns-majorant
The number of group elements that bring a small ball of the chart back to itself is bounded through the
majorants. These change by at most a factor $2$ across the ball, are carried along by the group, and do
not change under upper triangular matrices.

% before: cv-returns-short-basis
A self-dual lattice with no short vectors for the majorant has a basis of bounded length, by transference
to its dual for the trace form. Self-duality bounds the vectors of that dual below in terms of the
discriminant of $E$.

% before: cv-returns-count
A group element that brings the ball back to itself moves the short basis to lattice vectors of bounded
length, and is determined by their images. Counting the lattice vectors of bounded length therefore
bounds the number of such elements.

% before: cv-returns-unfold
The integral over the ball is compared with the integral over a fundamental set by unfolding. Each point of
the fundamental set is counted at most once for each group element that brings the ball back to itself.

% before: cv-tube-chart
The total measure of the fundamental sets is bounded through covolumes in products of special unitary
groups. For a group of elements of determinant $1$, the measure of a fundamental set in the chart is a fixed
multiple of such a covolume.

% before: cv-tube-det
The groups attached to the components are principal congruence subgroups of level $M\ge3$. Their elements
have determinant $1$, so they lie in the special unitary groups.

% before: cv-tube-index
The stabilizer of a self-dual lattice is covered by a bounded number of cosets of its principal congruence
subgroup. The number is at most the number of possible images of a set of generators modulo $M$.

% before: cv-tube-archimedean
The covolumes are measured with Haar measures on the special unitary groups, while the volume formula of
Theorem~\ref{cv-volume-mass} uses the measure $m$ on $U(P,Q)$. The two are compared through a projection
from $U(P,Q)^\Phi$ to the special unitary group.

% before: cv-tube-bound
Tubes over the components, built from packing sets for the covolumes, are disjoint in the adelic quotient
because the components lie in distinct double cosets. The volume formula and the index bound then bound
the sum of the covolumes by the mass.

% before: cv-cusp-frame-graph
The lower bound for integral forms evaluates them at a frame attached to a basis of a lattice adapted to the
signature of $G_Z$. Such a basis is built from lattice vectors close to large multiples of the coordinate
axes. Vectors near the first $p$ axes span a subspace on which the form of signature $(p,r-p)$ is positive,
and the vectors orthogonal to them have negative length.

% before: cv-cusp-frame-covering
A lower bound for the lengths of the lattice vectors at a frame gives, by transference, a bound for the
covering radius of the lattice and short linearly independent vectors in it. As in
Lemma~\ref{cv-returns-short-basis}, the dual lattice is described through the trace of $G_Z$, and the
discriminant of $E$ carries it back into the lattice.

% before: cv-cusp-frame-basis
Lattice vectors near the first $q$ axes give the positive part of an adapted basis. Projecting short vectors
orthogonally to them, with the Gram determinant clearing denominators, gives the negative part. The sizes of
both parts are bounded by Lemma~\ref{cv-cusp-frame-graph}.

% before: cv-cusp-frame-block
The frame at a CM point is found near an orthonormal frame by tilting the span of its first $q$ columns.
The tilt at each embedding is a complex $(r-q)\times q$ matrix, and its size is measured by the operator
norm.

% before: cv-cusp-frame-lattice
At an orthonormal frame with the same positive span as a frame where the lattice has no short vectors, the
basis of Lemma~\ref{cv-cusp-frame-basis} has a positive span close to the coordinate one. A small tilt of the
frame then has that positive span, and the lengths of the basis vectors are bounded in terms of the lower
bound.

% before: cv-cusp-frame-chart
Small tilts of an orthonormal frame are again frames, and they change the normalizing factors and the
majorants only slightly. This lets a form that is large at the orthonormal frame stay nonzero at the tilted
one.

% before: cv-cusp-frame-nonvanishing
A form whose pointwise norm at a frame is at least half of its supremum has no zeros on a ball around that
frame in the chart of tilts. Along a complex line through a zero, the Schwarz lemma would make the value at
the centre smaller than it is.

% before: cv-cusp-frame
A frame where a cusp form nearly attains its supremum, and where the lattice has no short vectors, can be
tilted to a frame whose positive span is spanned by lattice vectors of bounded length. The form does not
vanish there, which is what the evaluation at a CM point in Lemma~\ref{cv-cusp-inputs}(3) needs.

% before: cv-cusp-cm-lines
The lower bound of Lemma~\ref{cv-cusp-inputs}(3) comes from the product formula for one nonzero coefficient
of a polynomial attached to the form at a CM point. The abelian variety at the CM point and the variables of
the polynomial carry lines over a number field, each with an ideal at the finite places and a length at each
infinite place. The degree of such a line adds up these contributions, and it is additive in tensor products.

% before: cv-cusp-cm-cauchy
The sizes of a polynomial at the points of a polydisc bound each of its coefficients. Averaging the values at
the points whose coordinates are roots of unity times the radii isolates one coefficient.

% before: cv-cusp-cm-numbers
A nonzero coefficient that lies in the ideal of a metrized line and has bounded lengths at the infinite places
bounds the degree of that line below, by the product formula. Bounds on the degrees of the lines from which
it is built turn this into a lower bound for the mean of the logarithms of the archimedean bounds. These
degree bounds are the conditions (i) to (iii).

% before: cv-cusp-cm-hadamard
Condition (i) bounds the degree of a determinant line below by the degrees of lines whose generators map into
it. At the finite places a comparison of ideals gives the bound, and at the infinite places Hadamard's
inequality for a Gram matrix does.

% before: cv-cusp-cm-key
Condition (iii) relates the line of theta sections of a principally polarized abelian variety to the line of
the determinant of its Lie algebra. The relation is Moret-Bailly's key formula,
Theorem~\ref{cv-definite-key-formula}. Two metrized lines are tied to its local and archimedean forms by the
conditions below.

% before: cv-cusp-cm-theta-degree
At each prime the key formula compares the order of the number $\kappa$ with the ideals of the two lines, and at
each infinite place it compares the absolute value of $\kappa$ with their lengths. The product formula for
$\kappa$ then gives the degree of the dual of the theta line.

% before: cv-cusp-cm-qr
The archimedean bound at the CM point needs the size of a function with a transformation law under upper
triangular matrices at families that are not unitary. The Gram--Schmidt process writes such a family as a
unitary family times upper triangular matrices, and the transformation law reduces the bound to the unitary
case.

% before: cv-cusp-cm-archimedean
A bound on the families that are orthonormal for given Hermitian forms gives a bound on all families whose
entries are small for the corresponding lengths. A square root of each form makes the orthonormal families
unitary, and Lemma~\ref{cv-cusp-cm-qr} applies after a small shift that makes the matrices invertible.

% before: cv-cusp-cm-integral
Integrality of the form at the finite places becomes integrality of the coefficients of its polynomial. A
polynomial with integral values at all integral matrices of unit determinant has integral coefficients once the
residue field is large, and an unramified extension makes it large. A diagonal change of variables then carries
this over to the ideals of the lines of the variables.

% before: cv-cusp-cm-product
The local conditions put a nonzero coefficient of the polynomial in the ideal of a line built from the theta
line and the lines of the variables, and the archimedean conditions bound its lengths. The key formula and the
Hadamard bound give the degrees of the lines, and Lemma~\ref{cv-cusp-cm-numbers} turns these into a lower
bound for the mean of the logarithms of the archimedean bounds.

% before: cv-cusp-cm-point-polynomial
The polynomial of Proposition~\ref{cv-cusp-cm-product} records the values of a section at all changes of frame of
one framed point over a number field. Over the ring of all changes of frame the value is a polynomial divided by a
power of the determinants, and the transformation law of the weight removes the denominator, as for the deep cover
in Lemma~\ref{cv-volume-forms-frame-polynomial}. A nonzero value at a complex framed point with the same point makes
the polynomial nonzero.

% before: cv-cusp-cm-point-good
The CM point supplies the remaining data of Proposition~\ref{cv-cusp-cm-product}. Its abelian variety is
isogenous, over the ring of integers, to a product of $p$ CM abelian schemes of type $\Phi$ and $q$ of the
complementary type. The images of their eigenvectors are the rows of the matrix $U$, and integral models and Hodge
norms make these eigenvectors into metrized lines.

% before: cv-cusp-cm-point-letters
The isogeny of a good CM point has an inverse up to multiplication by an integer. So a linear relation among the
images of the eigenvectors of the factors gives one among the eigenvectors of a single factor, which form a basis,
and the matrix $U$ is invertible.

% before: cv-cusp-cm-point-model
The conditions at a finite prime are read off from integral models. A model matrix converts coordinates in a basis
over the generic fibre into coordinates in an integral basis at the prime, so it tells which vectors lie in the
Lie algebra of the abelian scheme after inverting the elements prime to it.

% before: cv-cusp-cm-point-local-ideal
The lines attached to the eigenvectors of a good CM point are described prime by prime through model matrices.
Taken over all primes together, these conditions identify the ideal of such a line with the multipliers that put
the eigenvector into the Lie algebra of the abelian scheme over $\mathcal O_K$.

% before: cv-cusp-cm-point-local-row
Each row of $U$ expresses a basis vector of the framed point as the image of an eigenvector of $A^\Phi$ or of $A'$.
Scaled by a local generator of its line, it therefore has integral coordinates in an integral basis of the Lie
algebra of $B$ at the prime. This gives the determinant bound in condition (L), and later the integrality of the
matrices $U'_\tau$.

% before: cv-cusp-cm-point-local-model
The integrality of the polynomial is tested at framed points over complete discrete valuation rings. Such a point is
integral when its abelian variety extends over the ring with bases of the eigenspaces of the Lie algebra and a
generator of the theta sections that restrict to those of the point.

% before: cv-cusp-cm-point-local-eigenbasis
Over a discrete valuation ring the eigenspaces of the Lie algebra of a model are free. A change of frame over the
fraction field therefore makes the bases of a framed point integral, and changes of frame over the ring keep them
integral.

% before: cv-cusp-cm-point-local-generator
The theta sections of a principally polarized model form a free module of rank one. So a scalar change of frame
makes the generator of a framed point integral, whatever the change of its bases.

% before: cv-cusp-cm-point-local-base
The test of integrality runs over all extensions of the ring. Base change carries integral bases and integral
generators to integral ones, since the Lie algebra and the theta sections commute with flat base change.

% before: cv-cusp-cm-point-local-situation
At a prime of $\mathcal O_K$, the abelian scheme $B$ of a good CM point gives a model of the framed point over any
complete discrete valuation ring that contains the local ring of the prime. The lines $\Theta$ and $\mathcal L_j$ of
the good CM point are described through models of the same abelian schemes, so the integral frame of this model can
be compared with them.

% before: cv-cusp-cm-point-local-unit
The scalar that makes the generator integral is determined up to a unit. The key formula compares it with the
local generator of the line $\Theta$, which is adapted to the key datum at the prime.

% before: cv-cusp-cm-point-local-letters
The change of frame to an integral basis is controlled by the lines $\mathcal L_j$. Its inverse becomes integral
after scaling each basis vector by a local generator of its line, because these vectors, so scaled, lie in the Lie
algebra of the model.

% before: cv-cusp-cm-point-local
Condition (L) of Proposition~\ref{cv-cusp-cm-product} has two parts. The bound on the determinant of $U$ comes
from the integral models: each row of $U$ is the image under a homomorphism of abelian schemes of an eigenvector
that the local generator of its line makes integral. The integrality of the polynomial comes from the integrality
of the section at an integral model of the framed point.

% before: cv-cusp-cm-point-metric
Condition (A) of Proposition~\ref{cv-cusp-cm-product} compares the values of a cusp form at the frames of the CM
point over $\mathbf C$ with its supremum norm. The comparison uses the Hodge form of the tangent vectors of a
framed point and the mean of its theta function, both read in a uniformization.

% before: cv-cusp-cm-point-metric-change
A change of frame acts on these data by the rules of linear algebra. The Gram blocks change by congruence with the
matrices of the change, and the constant by the square of the absolute value of the scalar.

% before: cv-cusp-cm-point-hodge
At a complex embedding the framed point of a good CM point is metric, with Gram blocks given by the Hodge matrix of
its key datum. That matrix has the determinant that condition (A) asks for, as the polarization is principal, and the
polarization multipliers bounded by $\ell^2$ bound the lengths of the rows of $U$.

% before: cv-cusp-cm-point-cusp-bound-ortho-hermitian
At a frame whose minors have absolute value $1$, the uniformization gives a Hermitian form whose imaginary part on the
periods is the alternating form of the lattice. At an orthonormal frame the periods are the coordinates of the lattice
vectors in the frame, up to conjugation, so this form is a multiple of the standard one.

% before: cv-cusp-cm-point-cusp-bound-ortho
The cusp form is bounded through its readings, which are its values at the framed points of frames. At an orthonormal
frame the uniformization gives the Hodge form and the theta function explicitly, so the metric data of that framed
point are known.

% before: cv-cusp-cm-point-cusp-bound-unique-key
Two uniformizations of one principally polarized complex abelian variety differ by a linear isomorphism, which carries
one Hermitian form to the other. The mean of the theta function changes with the volume of the fundamental domain, and
the key formula ties that change to the change of the determinant of the tangent vectors.

% before: cv-cusp-cm-point-cusp-bound-unique
The metric data of a framed point are defined through a choice of uniformization. Comparing two framed points through
their metric data needs that the constant and the Gram blocks do not depend on that choice.

% before: cv-cusp-cm-point-cusp-bound-good
At a good embedding a framed point is a triangular change of the framed point of an orthonormal frame. The metric data
fix the size of the diagonal of that change, and the transformation law turns it into a factor in front of the
reading.

% before: cv-cusp-cm-point-cusp-bound-conjugate
The other embeddings are the complex conjugates of the good ones. Conjugating a framed point conjugates its metric
data, so the identity Gram blocks are kept.

% before: cv-cusp-cm-point-cusp-bound
At a framed point whose bases are orthonormal for the Hodge form, the value of a cusp form is bounded by its
supremum norm. The power of the constant accounts for the theta generator, whose norm the weight $n$ involves.

% before: cv-cusp-cm-point-arch
These three facts give condition (A). A change of frame that makes the Gram blocks the identity turns the bound of
Lemma~\ref{cv-cusp-cm-point-cusp-bound} into a bound for the polynomial on the families that condition (A)
quantifies over.

% before: cv-cusp-cm-point-eigen-degree-line
The hypothesis on the degrees $\widehat\deg\mathcal L_j$ in Proposition~\ref{cv-cusp-cm-product} concerns lines of
eigenvectors in the Lie algebra of a CM abelian scheme. A nonzero vector of the Lie algebra of the generic fibre
determines such a line, through its integral multiples and its Hodge norms, and its degree is computed after base
change to larger number fields.

% before: cv-cusp-cm-point-eigen-degree-curve
The CM curve of Lemma~\ref{cv-cm-point}(1) is defined over a number field that need not contain the field of the
abelian scheme. After base change to a larger field it is still a CM curve, and the mean of the logarithms of its
Hodge norms, which the Chowla--Selberg formula computes, does not change.

% before: cv-cusp-cm-point-eigen-degree-base
Under an extension of number fields the ideal of the line of an eigenvector extends and its Hodge norms are restricted,
so its degree is multiplied by the degree of the extension. The normalized degree can therefore be computed over any
larger field, where the abelian scheme decomposes.

% before: cv-cusp-cm-point-eigen-degree-serre
An abelian scheme whose CM type is induced from $\mathbf Q(\sqrt{-d})$ is built from a single CM elliptic curve and a
lattice over $\mathcal O_E$. In Serre's tensor construction its points are the tensor product of the lattice with the
points of the curve, and its polarization is given by an $\mathcal O$-valued form on the lattice.

% before: cv-cusp-cm-point-eigen-degree-tensor-family
A Serre tensor product is recognized through its homomorphisms from the curve. A family of such homomorphisms, indexed
by the lattice and compatible with the actions and the polarizations, is the datum that is constructed over
$\mathbf C$ and then carried to the ring of integers of a number field.

% before: cv-cusp-cm-point-eigen-degree-tensor-complex-positive
On the part of a period where $\sqrt{-d}$ acts by a single complex number, the element $w_0$ acts by a scalar $\omega$.
There the alternating form at $(w_0x,x)$ is the imaginary part of $\omega$ times a positive number, which fixes the
sign of the form $h$.

% before: cv-cusp-cm-point-eigen-degree-tensor-laws
Several statements of this section take standard facts on abelian schemes as hypotheses: the theorems of the cube and
of the square, the commutator pairings of theta groups, base change of polarizations, and the extension of line
bundles over discrete valuation rings. They are collected here, each in the form in which it is assumed, so that each
statement can say which of them it needs.

% before: cv-cusp-cm-point-eigen-degree-tensor-complex-pairing-hermitian
Over $\mathbf C$ the commutator pairings at torsion points are exponentials of the imaginary parts of the
Appell--Humbert forms. So two pairs of homomorphisms have the same pairings exactly when the forms pulled back along their
linear lifts agree; this turns the Rosati condition, and later the identities of a family, into identities of Hermitian
forms.

% before: cv-cusp-cm-point-eigen-degree-tensor-complex-datum-structure
A lattice with an action of $\mathcal O_E$, a perfect adjoint form and a period of a given type is the lattice of a PEL
datum. The period identifies the real span of the lattice with $\mathbf C^g$, whose multiplication by $i$ is the complex
structure, and the type gives its signature.

% before: cv-cusp-cm-point-eigen-degree-tensor-complex-datum
The lattice of a uniformization of the abelian variety, with the action of $\mathcal O_E$ and the imaginary part of the
Appell--Humbert form, is a PEL lattice. The dimensions of the eigenspaces of the Lie algebra diagonalize the action in
the coordinates of a type, which gives a period.

% before: cv-cusp-cm-point-eigen-degree-tensor-complex-homs-line
On the pure part of a period, $w_0$ acts by a scalar, so the multiples of $P(a)$ by $\mathcal O$ lie in the period
lattice. The Hermitian form on the complex line through $P(a)$ is then fixed by the value $\langle w_0a,a\rangle$.

% before: cv-cusp-cm-point-eigen-degree-tensor-complex-homs-curve
For the curve the same computation runs the other way. The lattice $c(\mathbf Z+\mathbf Z\frac{1+i\sqrt d}2)$ is known,
and the unimodularity of the imaginary part of the form fixes the form.

% before: cv-cusp-cm-point-eigen-degree-tensor-complex-homs-points
The constructions below define regular functions and morphisms on a complex abelian variety by their values at the
images of a uniformization. These are all the complex points, and on a reduced scheme of finite type over $\mathbf C$ the
complex points determine open sets, functions and morphisms.

% before: cv-cusp-cm-point-eigen-degree-tensor-complex-homs-sign
Pulling a theta function back along a linear map multiplies its factor of automorphy by a power of the semicharacter of
the second uniformization. For the line bundle of a polarization, which is symmetric, the semicharacter takes the values
$\pm1$, so squares of pulled back theta functions are again theta functions.

% before: cv-cusp-cm-point-eigen-degree-tensor-complex-homs-glue
A morphism to a separated scheme can be built from functions. Suppose that near each complex point the regular functions on
an affine open set of the target pull back to regular functions with the right values. Then the local morphisms they
define agree on overlaps and glue.

% before: cv-cusp-cm-point-eigen-degree-tensor-complex-homs-hom
A linear map between uniformizations that carries one lattice into the other and scales the Hermitian forms is the lift
of a homomorphism. Pulled back along it, theta functions for the second uniformization become theta functions for the
first, so they give sections that define the homomorphism.

% before: cv-cusp-cm-point-eigen-degree-tensor-complex-homs-bundle
The same theta functions compare the pullback of the polarization of the second abelian variety with a power of that
of the first. The comparison holds after raising both line bundles to a power.

% before: cv-cusp-cm-point-eigen-degree-tensor-complex-homs
An element $a$ of the pure part of the lattice gives a complex line $\mathbf CP(a)$ whose lattice contains a copy of that
of the curve. The linear map from the uniformization of the curve onto this line scales the Hermitian forms, and so it
comes from a homomorphism of abelian varieties.

% before: cv-cusp-cm-point-eigen-degree-tensor-complex-pairing-polarize
A Hermitian form on the image of the lattice is determined by its diagonal values once $\sqrt{-d}$ acts by a scalar.
Polarization gives the off-diagonal values from the diagonal ones, and the action of $\sqrt{-d}$ separates the real and
imaginary parts.

% before: cv-cusp-cm-point-eigen-degree-tensor-complex-pairing-lambda-fibre
On the fibre of a product over an $n$-torsion point, the line bundle $\Lambda$ pulled back along multiplication by $n$
has a generator that translation by an $n$-torsion point multiplies by the commutator pairing. This is how the pairings
control the line bundle $\Lambda$ on these fibres.

% before: cv-cusp-cm-point-eigen-degree-tensor-complex-pairing-lambda-quotient
Multiplication by $n$ on an abelian variety over an algebraically closed field, with $n$ invertible, is the quotient by
its kernel. A function on the inverse image of an open set that is invariant under the translations by $n$-torsion points
comes from that open set.

% before: cv-cusp-cm-point-eigen-degree-tensor-complex-pairing-lambda-descent
When functions invariant under a family of maps descend along a morphism, so do generators of line bundles. Two line
bundles whose pullbacks have generators with an invariant ratio are then isomorphic.

% before: cv-cusp-cm-point-eigen-degree-tensor-complex-pairing-lambda-dense
The seesaw argument needs the triviality on a dense set of fibres. The fibres over the torsion points suffice, as these
points are dense.

% before: cv-cusp-cm-point-eigen-degree-tensor-complex-pairing-lambda
An isomorphism of line bundles $\Lambda$ gives equal commutator pairings. The converse holds over an algebraically closed
field of characteristic $0$, through the seesaw principle on the fibres over the torsion points.

% before: cv-cusp-cm-point-eigen-degree-tensor-complex-pairing
The last condition of a family, the isomorphisms of the line bundles $\Lambda$, is checked on Hermitian forms. On the
diagonal they are given by the isomorphisms of powers of line bundles, and polarization recovers the off-diagonal
values.

% before: cv-cusp-cm-point-eigen-degree-tensor-complex
Over $\mathbf C$ both the curve and the abelian variety are quotients of complex vector spaces by lattices. When the
Lie algebra of the abelian variety has the type induced from $\mathbf Q(\sqrt{-d})$, its lattice is a PEL lattice and
each of its elements defines a homomorphism from the curve.

% before: cv-cusp-cm-point-eigen-degree-tensor-identities-section
The isomorphisms in a family of homomorphisms are isomorphisms of line bundles, and they descend from a field to a
Dedekind domain through the triviality of a quotient line bundle. The first step produces a global section that is a
generator after a flat base change, from the flat base change of global sections.

% before: cv-cusp-cm-point-eigen-degree-tensor-identities-field
On a proper geometrically integral scheme over a field, the global functions are constants, also after any extension of
the field. So the section of the previous lemma is a generator after the extension, and hence before it.

% before: cv-cusp-cm-point-eigen-degree-tensor-identities-dedekind
Over a Dedekind domain a line bundle on an abelian scheme can be trivial on the generic fibre without being trivial,
through the pullback of a nontrivial ideal class. Triviality along the unit section excludes this: the section of the
generic fibre is divided by the constant it takes along the unit section, place by place.

% before: cv-cusp-cm-point-eigen-degree-tensor-identities-line
The two previous lemmas combine into a descent statement for isomorphisms of line bundles, applied to the quotient of
the two line bundles. The field may be any field into which the Dedekind domain embeds.

% before: cv-cusp-cm-point-eigen-degree-tensor-identities-lambda
The line bundles $\Lambda$ of the identities of a family are trivial along the unit section, as the homomorphisms
preserve it, and their formation commutes with base change. So an isomorphism between two of them over the field
descends to the Dedekind domain.

% before: cv-cusp-cm-point-eigen-degree-tensor-identities
The identities of a family are equalities of morphisms and isomorphisms of line bundles. Over a Dedekind domain they
can be checked after an injective base change to a field, which both the descent and the extension below use.

% before: cv-cusp-cm-point-eigen-degree-tensor-descent-specialize
A homomorphism over $\mathbf C$ is defined over a finitely generated subalgebra of $\mathbf C$, and it is specialized to
a number field by a homomorphism of that subalgebra. The specialization must fix the elements that are algebraic over the
base field. Those are recognized by having only countably many conjugates.

% before: cv-cusp-cm-point-eigen-degree-tensor-descent-countable
The specialization of a homomorphism is compared with the original through countability. Over $\mathbf C$ a
homomorphism of abelian varieties lifts to a linear map that preserves the lattices of uniformizations. This leaves
countably many homomorphisms, each determined by its action on tangent vectors.

% before: cv-cusp-cm-point-eigen-degree-tensor-descent-points
A family of points parametrized by a finitely generated subalgebra of $\mathbf C$ has coordinates in that subalgebra.
If its specializations take only countably many values, these coordinates are algebraic, so a specialization that
fixes algebraic elements does not change the family.

% before: cv-cusp-cm-point-eigen-degree-tensor-descent-hom
The descent of one homomorphism to a number field combines the three previous statements. Spreading out puts it over a
finitely generated subalgebra of $\mathbf C$, specialization puts it over a number field, and the two complex
homomorphisms agree on tangent vectors because they agree on a family of points with countably many specializations.

% before: cv-cusp-cm-point-eigen-degree-tensor-descent
The family over $\mathbf C$ is indexed by a finitely generated group, so finitely many of its homomorphisms determine
it. Each of these is defined over a number field, and so the whole family is defined over their compositum.

% before: cv-cusp-cm-point-eigen-degree-tensor-extend
Over the ring of integers the homomorphisms of the family extend, since homomorphisms of abelian schemes over a
Dedekind domain extend from the generic fibre. The identities then extend with them.

% before: cv-cusp-cm-point-eigen-degree-tensor-recognition-tensor
To recognize an abelian scheme as a Serre tensor product, one first needs a tensor product to compare it with. Serre's
construction tensors the curve with the lattice; the positivity of $\varepsilon\langle w_0x,x\rangle$ makes its line
bundle relatively ample, and the perfectness of the form makes the polarization principal.

% before: cv-cusp-cm-point-eigen-degree-tensor-recognition-cube
The points of a Serre tensor product are sums of pure tensors, and the line bundles $\Lambda$ are additive in each
variable when the theorem of the cube holds. So an isomorphism of line bundles $\Lambda$ on a tensor product can be
checked on pure tensors.

% before: cv-cusp-cm-point-eigen-degree-tensor-recognition-lift
A family of homomorphisms from the curve defines a homomorphism from a Serre tensor product, through the universal
property of the tensor product of modules on points. The identities of the family become the compatibility of this
homomorphism with the actions and with the polarizations.

% before: cv-cusp-cm-point-eigen-degree-tensor-recognition-onto
Surjectivity of a morphism of abelian schemes over a domain can be read off from one injective complex point of the
base. The complex points are dense in the fibre, and the fibre at the generic point is dense in the abelian scheme.

% before: cv-cusp-cm-point-eigen-degree-tensor-recognition-iso
A homomorphism that pulls one polarization back to another has its kernel inside the kernel of the second
polarization. When that polarization is principal, the homomorphism is a monomorphism, and a proper surjective
monomorphism onto a reduced scheme is an isomorphism.

% before: cv-cusp-cm-point-eigen-degree-tensor-recognition
A family indexed by the whole lattice defines a homomorphism to the abelian scheme from the Serre tensor product. When
the family is onto at a complex point, this homomorphism is surjective, and it respects principal polarizations, so it
is an isomorphism.

% before: cv-cusp-cm-point-eigen-degree-tensor
After a finite extension of the base, every abelian scheme of this type is a Serre tensor product. The lattice comes
from the complex uniformization, and the homomorphisms from the curve descend to a number field and extend over its
ring of integers.

% before: cv-cusp-cm-point-eigen-degree-ideal
The eigenvector lines need only the case of a lattice of rank one over $\mathcal O_E$. Such a lattice is a fractional
ideal of $E$ with a trace form, and its self-duality ties the norms of the ideal and of the coefficient of the form to
the discriminant of $E$ and to $d$.

% before: cv-cusp-cm-point-eigen-degree-dual
The Lie algebra of a Serre tensor product is described through the form $h$. Its perfectness over $\mathcal O$
provides dual families, which recover every element of the lattice from its values under $h$.

% before: cv-cusp-cm-point-eigen-degree-rank-one-norm-weil-translate-group
The Hodge norms of the tensor product are compared with those of the curve through commutator pairings. A point killed
by $n$ lifts to the theta group of $\mathcal L^n$, and the commutator of two lifts is a unit on the base, which defines
the pairing.

% before: cv-cusp-cm-point-eigen-degree-rank-one-norm-weil-translate-pairing
Over a base with one point, the cubical structure gives explicit lifts to the theta group. From them one builds a
generator of the pullback of $t_x^*\mathcal L\otimes\mathcal L^{-1}$ along multiplication by $n$, which translation by a
point killed by $n$ multiplies by the commutator pairing.

% before: cv-cusp-cm-point-eigen-degree-rank-one-norm-weil-translate
Translation by a point killed by $n$ acts on a line bundle built from $\Lambda_{\mathcal L}$ through the commutator
pairing of the theta group of $\mathcal L^n$. This follows from the previous lemma by pulling back along the
homomorphism to $A$.

% before: cv-cusp-cm-point-eigen-degree-rank-one-norm-weil
An isomorphism between two line bundles of the form $\Lambda_{\mathcal L}$ on a product of abelian varieties carries
one generator of the previous lemma to a constant multiple of the other. Comparing them at a point killed by $n$ and at
the origin shows that the commutator pairings agree.

% before: cv-cusp-cm-point-eigen-degree-rank-one-norm-riemann-theta-sign
The sign of the exponent in the commutator pairings depends on the conventions for theta groups and for the factors of
automorphy. A direct computation with explicit generators confirms the sign of the theorem in the conventions used here.

% before: cv-cusp-cm-point-eigen-degree-rank-one-norm-riemann-theta
A principally polarized complex abelian variety is a quotient of $\mathbf C^g$ by a lattice, and the sections of the
powers of its line bundle are theta functions. The Hermitian form of the Appell--Humbert data is integral and
unimodular on the lattice, and it determines the commutator pairings of the theta groups.

% before: cv-cusp-cm-point-eigen-degree-rank-one-norm-riemann
Over $\mathbf C$ the commutator pairings are read off from the Hermitian form of the Appell--Humbert data. This holds
for the theta uniformization of a principally polarized abelian variety, and any other uniformization with
Appell--Humbert data differs from it by a linear isomorphism.

% before: cv-cusp-cm-point-eigen-degree-rank-one-norm
The polarization of a tensor product pulls back along each homomorphism from the curve to that of the curve composed
with an endomorphism. Through the two previous lemmas this relation of polarizations becomes a relation of Hermitian
forms, and so a formula for the squared Hodge norm of the image of a tangent vector.

% before: cv-cusp-cm-point-eigen-degree-rank-one
For a tensor product with a lattice of rank one the eigenvector lines can be computed directly. Their ideals come from
the fractional ideal of the lattice and their Hodge norms from those of the curve, so the degree is the
Chowla--Selberg mean of the curve plus a term from the discriminant of $E$.

% before: cv-cusp-cm-point-eigen-degree
The degrees of the lines of the eigenvectors enter Proposition~\ref{cv-cusp-cm-product} through the hypothesis on
$\widehat\deg\mathcal L_j$. For an abelian scheme whose CM type is induced from $\mathbf Q(\sqrt{-d})$, these lines
all have normalized degree $\frac12\log\Omega_E(\chi)$, which is where the Chowla--Selberg formula enters.

% before: cv-cusp-cm-point-degree
The lines $\mathcal M^\Phi_i$ and $\mathcal M'_i$ of a good CM point are lines of eigenvectors of this kind. The
abelian schemes $A^\Phi$ and $A'$ have the types of $\Phi$ and of its complement, which are induced from the two
embeddings of $\mathbf Q(\sqrt{-d})$.

% before: cv-cusp-cm-point-exists-lattice-pair
The lattice of the component is spanned by a fractional ideal times the self-dual lattice of the model. The trace form of
the model, restricted to that ideal, makes it a lattice of rank one, of either signature according to the sign.

% before: cv-cusp-cm-point-exists-lattice-pure
The vectors of a basis adapted to a frame have pure periods. Those of the first $q$ span the columns of the frame read by
one kind of coordinate, and the others are orthogonal to them.

% before: cv-cusp-cm-point-exists-lattice
The basis of Lemma~\ref{cv-cusp-frame} splits the lattice of the component, up to a finite index, into rank one pieces
whose periods are of pure type. Each piece is a lattice of signature $(1,0)$ or $(0,1)$, and the form of the component
restricts to it as a multiple of its own form, with a multiplier bounded by $\ell^2$.

% before: cv-cusp-cm-point-exists-object
Rank one lattices of pure type are the period lattices of abelian varieties with complex multiplication by
$\mathcal O_E$. These are defined over a number field and have good reduction everywhere after an extension, and the maps
into the lattice of the family uniformization become homomorphisms to an abelian scheme whose generic fibre is the
abelian variety at the period.

% before: cv-cusp-cm-point-exists-good
From such an integral CM object and a key datum on its generic fibre, the remaining data of a good CM point are the key
data of the CM factors, their lines, the letters and the framed point. The bound $\ell^2$ on the multipliers passes to
the polarization condition of the good CM point.

% before: cv-cusp-cm-point-exists
The CM point is built at the frame and the basis of Lemma~\ref{cv-cusp-frame}. The basis splits the lattice into
lines of positive and negative norm, and the CM abelian schemes and the isogeny of a good CM point come from this
splitting.

% before: cv-cusp-cm-point
The data of Proposition~\ref{cv-cusp-cm-product} come from a CM point at the frame of
Lemma~\ref{cv-cusp-frame}. Its abelian variety is isogenous to a product of CM abelian varieties of
dimension~$f$, whose eigenvectors in the Lie algebra span lines of normalized degree
$\frac12\log\Omega_E(\chi)$. The values of the form at the frames of the CM point make up the polynomial.

% before: cv-cusp-inputs
Three inputs enter the two facts. The first bounds the value of a cusp form at a point by its
integral over a fundamental set, and the second bounds the total volume of the fundamental sets by a
mass. The third bounds integral forms below through their values at CM points, which is where
$\log\Omega_E(\chi)$ enters.

% before: cv-cusp-torus-bound
The bounds for cusp forms below carry a hypothesis on the classes of the norm-one torus of the host
field, which their proofs do not use. It bounds the number of these classes modulo the determinants of
a principal congruence subgroup by the class number times a power of the level.

% before: cv-cusp-torus-hilbert90
An element of norm one of the finite adeles of $E'$ is a quotient $x/\bar x$ of a unit, as in Hilbert's
Theorem 90. The two candidates $1+t$ and $d(1-t)$, with $\bar d=-d$, are combined place by place with an
idempotent, so that the result is a unit at every place.

% before: cv-cusp-torus-ideal-classes
Writing an element of norm one as $z/\bar z$ with $z$ integral attaches to it the ideal of $E'$ that $z$
generates. Two elements whose ideals lie in the same class differ by a global element of norm one and an
integral element of norm one, so there are at most $h_{E'}$ classes modulo these.

% before: cv-cusp-torus-congruence
An integral element of norm one has a residue modulo $M$, which takes at most $M^{[E':\mathbf Q]}$ values.
Dividing it by a fixed element of norm one with the same residue leaves an element congruent to $1$ modulo
$M$.

% before: cv-cusp-torus-norms
The determinants in the torus class bound come from unitary maps that multiply one vector by a scalar and
fix its orthogonal complement. Such a vector must have a length that is a unit at every finite place. It
exists because, when $E'$ is unramified over $F'$, the lengths of the vectors of a self-dual lattice
generate the unit ideal.

% before: cv-cusp-torus-determinant
An element of norm one congruent to $1$ modulo $M$ is the determinant of the unitary map that multiplies
by it a vector whose length is a unit, and fixes the orthogonal complement. This map preserves the completed
lattice and is congruent to the identity modulo $M$.

% before: cv-cusp-torus-bound-holds
The torus class bound is a theorem, so the bounds that carry it hold unconditionally. Its proof splits a class
of the torus into a class of ideals, a class modulo the level, and a determinant.

% before: cv-cusp-sup-bound
Summing the local bound over the components turns it into a bound by the Petersson norm, because the
adelic probability measure unfolds into the invariant measures on the fundamental sets. The constant
is the product of the local constant and the mass bound.

% before: cv-cusp-bergman
A bound for the values of every function in a space by its $L^2$ norm, uniform over the points,
bounds the dimension of the space. The argument is Bergman's: an orthonormal basis has pointwise sum of
squares at most the constant, and integrating that sum gives the dimension.

% before: cv-cusp-rank-bound
The dimension of the space of cusp forms of a dominant weight is bounded by reading the forms at one
good embedding. Readings preserve linear independence, and by Lemma~\ref{cv-cusp-sup-bound} the space of
readings satisfies the hypothesis of Lemma~\ref{cv-cusp-bergman}.

% before: cv-cusp-norm-bound
For an integral form, the mean over the embeddings of the logarithms of its supremum norms is bounded
below by Lemma~\ref{cv-cusp-inputs}(3). The supremum norm is at most a constant times the Petersson norm
at every embedding, so the same mean for the Petersson norms is bounded below as well.

% before: cv-realization-level-move
The level group of a point acts on the Shimura varieties of the chain by changing the level structure
through an endomorphism of the lattice. Such a change is an automorphism of the moduli scheme together
with an isomorphism of the universal abelian scheme over it, and it moves framed points. The bad set of
Proposition~\ref{cv-point-conditions} is bounded by comparing the values of forms at framed points and at
their moves.

% before: cv-realization-rigidity
A level automorphism that fixes a point acts with finite order on the period lattice of the abelian
variety at that point. The lemma below shows that an endomorphism of finite order of a lattice is the
identity once it is unipotent modulo a prime of residue characteristic larger than $r+1$. It serves here
as Serre's lemma does for automorphisms congruent to $1$ modulo an integer at least $3$.

% before: cv-realization-fixed
The level group acts freely on the points of the zero-dimensional Shimura variety of the definite group
if no nontrivial element has a level automorphism fixing a point. At a fixed point the automorphism of the
abelian variety has finite order, and the level structure reads its matrix on the period lattice, where
Lemma~\ref{cv-realization-rigidity} applies.

% before: cv-realization-definite
For the definite group of the chain the moduli scheme is the spectrum of a product of fields, and the level
group acts on its functions through the level automorphisms. Multiplication by functions is compatible
with the action on forms, and the action on complex points is free when no nontrivial element of the
group fixes the point of a framed point. These are the hypotheses of Lemma~\ref{cv-realization-split}.

% before: cv-realization-levels
On a deep cover the level automorphisms form a larger group, which maps onto the level group of the point
by reduction modulo the primes at the level places. The Chinese remainder theorem splits this map by a
homomorphism whose image commutes with the kernel, which is the structure Lemma~\ref{cv-bad-part}
requires.

% before: cv-realization-pullback
Pulling forms back from the Shimura variety to a deep cover must be compatible with the level
automorphisms of the two, and the forms on the cover that are invariant under the kernel of the reduction
must be constant on the fibres of the level map. Both follow from the moduli interpretation: a level
automorphism changes the level structure, the level map multiplies it by a power of $15$, and the framed
points over a given framed point form one orbit of the kernel.

% before: cv-point-group-determinant
The regular representations on the deep cover come from forms whose group determinant over a group of
level automorphisms is not a zero divisor. Evaluating forms at the moves of one complex point on each
component gives finitely many linear maps, and the group determinant of the form must be nonzero for each
of them. Over a field of characteristic $0$ a generic linear combination of forms has this property.

% before: cv-point-regular-section
A nonzero value at one complex point on each component is enough for a section with an exponent not to be a
zero divisor. At the generic framed point of a component, scaling the frame separates the parts of a
section by their exponents, through the linear independence of characters.

% before: cv-point-conditions-separation
The auxiliary forms of the bad-set estimate need complex points on each component of the cover at which cusp
forms separate the moves by a unipotent group of level automorphisms. Separation by sections of a large weight
comes from the positivity of the Hodge line bundle on a compactification.

% before: cv-point-conditions-spanning
The ranks in the bad-set estimate are dimensions over the coefficient field only if the integral cusp forms
span all cusp forms. A cusp form becomes integral after multiplication by a nonzero constant.

% before: cv-point-conditions
The covolume bound follows from six properties of the lattices at each point. The first
four concern the lifted forms and the scale; the last two compare the lifted forms with
integral automorphic forms on the Shimura variety of $U(W_j)$ and bound the slope of the
latter.

% before: cv-definite-orbits
For the definite group of the chain the Shimura variety is a finite set of points, and a form is a
family of values, one at each point. The forms of the prescribed type are eigenvectors of a group of
level automorphisms that permutes the points freely, so they are determined by their values at one
point of each orbit. Their slope is then bounded below by the slopes of the lattices at single points.

% before: cv-definite-fibre-lattice
At a single point the integral forms are cut out by local matrices at every place, acting on a tensor
product of dual Weyl modules. The slope of such a lattice is a sum of degrees of determinants, less
half the logarithm of the dimension, which comes from comparing the basis of bideterminants with an
orthonormal basis.

% before: cv-definite-key-formula
At a point the degree of the theta line is compared with the degree of the determinant of the Lie
algebra by Moret-Bailly's key formula. Its archimedean form contributes the constant
$\frac14\log2\pi$ to the bound.

% before: cv-definite-situation
The next lemmas concern the Shimura variety of a definite unitary group. They describe its points and
its forms through frames: bases of the eigenspaces of the Lie algebra and a generator of the theta
line, on which a form is a function with a fixed law under triangular changes of frame.

% before: cv-definite-point
Each component of the definite Shimura variety is a rational point, and its frames at all places
come from one rational framed point changed by local matrices. The degrees of these matrices are tied
to the flat sections of Theorem~\ref{cv-flat-degree} and, through the key formula, to the theta line.

% before: cv-definite-fibre
Lemmas~\ref{cv-definite-fibre-lattice} and~\ref{cv-definite-point} together describe the lattice of
integral forms supported on one component and bound its slope below in terms of the degree of the
flat sections. On the definite group this gives the
lower bound of Proposition~\ref{cv-arithmetic-volume-step}.

% before: cv-isotypic-twisted-trace
On the deep cover of the indefinite groups, the forms of the prescribed type are the eigenvectors of
a character of a finite group of level automorphisms. Multiplying a form by the translates of one
auxiliary form and averaging against the character gives one map between weights for each element
of the group. When the group determinant of the auxiliary form is not a zero divisor, these maps
are jointly injective, and the dimension of the eigenspace times the order of the group is at most
the dimension of the space of the larger weight.

% before: cv-lattices-galois-different
The coefficient field of the deep cover is a Galois extension of that of the Shimura variety. An
integral combination of pulled-back forms need not be a combination of integral pulled-back forms
with integral coefficients, but it becomes one after multiplication by any element of the
different. This is a property of Galois-stable lattices, proved with the trace form.

% before: cv-lattices-different-slope
With this, the slope of the lattice of integral forms on the cover exceeds the slope of the lattice
it comes from by at most twice the logarithm of the root discriminant of the larger field. The
pullback preserves norms, so the only loss is the index of the extended lattice, which the
different controls.

% before: cv-height-components-numerics
The lower bound for the slope on the deep cover combines, component by component, a lower bound of
Hilbert--Samuel type for degrees, comparisons of ranks with geometric degrees, and lower bounds for
heights that hold on average over the components. The combination is a computation with real
numbers and uses nothing about the cover beyond these inequalities.

% before: cv-deep-cover
For $j<m$ the lower bound for slopes of integral cusp forms on the moduli scheme of $U(W_j)$ is
proved on a cover of deeper level. The cover carries the same kind of data as the moduli scheme it
covers: a coefficient field, a level map to the moduli scheme of $U(W_j)$, and uniformizations at
the good embeddings with an adelic description. Its level is deepened by a power of $15$, at primes
that split in $\mathbf Q(\sqrt{-d})$.

% before: cv-uniformization-model
The uniformizations of a deep cover are written in coordinates in which the Hermitian matrix of $W_j$
is fixed. Each lattice of a PEL datum is placed in $E^r$ by a map that turns its alternating form into
the trace of a Hermitian form, and its arithmetic group is cut out by the level.

% before: cv-uniformization-signature
The element $\eta$ of a Hermitian model is purely imaginary, and the sign of its imaginary part at an
embedding of $\Phi$ decides whether $p$ or $q$ counts the positive eigenvalues there. This sign fixes
the shape of the frames on which the uniformization is defined.

% before: cv-uniformization-lattice
The lattice of the PEL datum of $X$ is the lattice of $W_j$, with the trace of $\xi_0$ times the
Hermitian form of $W_j$ up to a sign. The signature of $W_j$ excludes the sign $-1$, so this
identification is a Hermitian model with the Hermitian matrix of $W_j$.

% before: cv-uniformization-models
The complex points of a PEL moduli scheme involve several lattices of the same signature, one for each
component. All of them can be placed in $E^r$ with the one Hermitian matrix of $W_j$, by the
classification of Hermitian forms. The arithmetic groups so obtained have finite index, act freely on
frames at a level of at least $3$, and contain transvections.

% before: cv-uniformization-family
Over $\mathbf C$, every point of a PEL moduli scheme is a complex torus $\mathbf C^g/P(L_i)$ for a period
$P$ of one of finitely many lattices, and two periods give the same point when they differ by a linear
map and a level automorphism. The theta functions of the polarization describe its line bundle on
$\mathbf C^g$.

% before: cv-uniformization-situation
The frames for the Hermitian matrix of $W_j$ parametrize the periods of each lattice $L_i$, through the
columns of the frame at the embeddings of $\Phi$. A frame also fixes tangent vectors and, through theta
functions, a section of the polarization; together they form the framed point that the uniformization
attaches to the frame.

% before: cv-uniformization-adapted-basis
The theta functions that normalize the framed points are taken in a symplectic basis of each lattice,
and an automorphism keeps their characteristic only if the basis is chosen well modulo $8$. Modulo an
integer $N$, a symplectic basis adapted to two isotropic subgroups that span can be built by linear
algebra, and it lifts to $\mathbf Z$ because every symplectic matrix modulo $N$ lifts to an integral one.

% before: cv-uniformization-mod-eight
Since $d\equiv7\pmod{16}$, the element $w^4$ acts on the lattice of a PEL datum as an idempotent modulo
$8$. Its image and its kernel modulo $8$ are isotropic modulo $8$ and together span the lattice, so they
have an adapted symplectic basis. In such a basis every $\mathcal O_E$-linear map is block diagonal
modulo $8$.

% before: cv-uniformization-adapted-blocks
An automorphism of the lattice that is congruent to $1$ modulo $M'$ commutes with $w^4$, so in the
adapted basis its matrix is block diagonal modulo $8$. Its first diagonal block has determinant $1$
modulo $8$, because the automorphism has determinant $1$ over $E$. These are the conditions under which
the theta transformation formula below has constant $1$.

% before: cv-uniformization-period-change
Composing a period with an automorphism of the lattice changes the normalized period matrix by the action
of a symplectic matrix on the Siegel space. The matrix that acts is read off from the blocks of the
automorphism in the symplectic basis, and the normalizing matrices change by the automorphy factor.

% before: cv-uniformization-theta-transformation
The theta function of characteristic $0$ transforms under a symplectic matrix by an automorphy factor.
For matrices whose off-diagonal blocks are divisible by $8$, its square transforms by $\det(CZ+D)$, an
exponential of a quadratic form, and a sign that depends on $\det D$ modulo $4$.

% before: cv-uniformization-theta-ratio
If such a matrix carries a theta function with a half-integral characteristic to a nowhere vanishing
multiple of itself, the characteristic is fixed modulo $1$ and the squared transformation formula holds
for it too. Comparing the two at the origin gives the square of the constant.

% before: cv-uniformization-theta-change
Two frames related by an element of $\Gamma_i$ give the same point and the same uniformization, so their
theta sections are proportional. The adapted bases and the transformation formula give the square of
the constant of proportionality, which the normalization $\kappa^2=N_\xi\det h_x$ needs for the two
normalized framed points to agree up to sign.

% before: cv-uniformization-frames
The normalized framed points at frames reach every complex framed point, and two of them agree up to a
triangular change when their frames are related by $\Gamma_i$ and $B$. Sections of automorphic bundles,
read at these points, are therefore holomorphic functions on frames with the transformation rules of
forms.

% before: cv-uniformization-marking
The adelic description of the components goes through lattices: an element of $U(W_j)(\mathbf A_F)$ moves
the lattice of $W_j$ to a new lattice, which a rational element moves onto the lattice of one component.
A marking records the component and the rational element, together with the level structure.

% before: cv-uniformization-markings
Every adelic element has a marking, and its component is determined, because the lattices with level
structure in the genus of $L$ are those of the components, each once. The rational element of a marking
is determined up to the arithmetic group of the component, and it changes predictably under
translations.

% before: cv-uniformization-adelic-frames
The archimedean part of an adelic element acts on the base frame through the embeddings of $\Phi$. The
frames obtained are orthonormal, so their leading minors have absolute value $1$, and they are all the
orthonormal frames.

% before: cv-uniformization-dictionary
Choosing a marking for every adelic element gives the map from $U(W_j)(\mathbf A_F)$ to frames of
components. Its invariance under rational elements and the level subgroup, its equivariance for the
archimedean group and its surjectivity up to $\Gamma_i$ and $B$ are the properties required of a deep
cover.

% before: cv-deep-pel-facts
A deep cover is built from three facts about PEL moduli schemes at an arbitrary level. They
give the moduli scheme of the cover, its map to the moduli scheme of $U(W_j)$, and its complex
uniformizations.

% before: cv-deep-ideal
The deep level multiplies the level of $X$ by a power of $15$ and its ideal by an ideal of
$\mathcal O_E$ coprime to its conjugate. Such an ideal exists when $3$ and $5$ split in
$\mathbf Q(\sqrt{-d})$, which is a condition on $-d$ modulo $15$.

% before: cv-deep-ray-field
The coefficient field of the cover must contain the roots of unity of the deep level and the values
that make the components of the cover geometrically connected. Both lie in a ray class field of
$\mathbf Q(\sqrt{-d})$, whose degree and root discriminant are bounded in terms of its modulus and
of $d$.

% before: cv-deep-compositum
The coefficient field of the cover is the compositum of $F_c$ with this ray class field. Its degree
and root discriminant are bounded by the products of those of the two fields.

% before: cv-components-conjugation
Lemma~\ref{cv-deep-components} below makes the coefficient field algebraically closed in the function
fields of the components by comparing complex points of the moduli scheme with their conjugates under
the automorphisms of $\mathbf C$ that fix it. Alp\"oge controls these conjugates through one elliptic
curve with complex multiplication by the ring of integers of $\mathbf Q(\sqrt{-d})$, and through the way
an automorphism moves its points of finite order.

% before: cv-components-weber-curve
The Weber values of level $n$ are the $x$-coordinates of the points of order $n$ on a model of this
curve over the field generated by its $j$-invariant. An automorphism of $\mathbf C$ that fixes them can
move such a point only to its negative, and a change of sign is absorbed into the isomorphism with the
conjugate curve.

% before: cv-components-multiplier
An automorphism of $\mathbf C$ that fixes $\sqrt{-d}$ commutes with complex multiplication on the curve,
so it acts on the points of order $n$ through multiplication by an element of the ring of integers.
Through the Weil pairing, the norm of this element is the power by which the automorphism acts on the
$n$-th roots of unity.

% before: cv-components-situation
The complex points of $X'$ are described by periods of finitely many Hermitian lattices, as in
Theorem~\ref{cv-uniformization-family}. The comparison of points with their conjugates reads the points,
their abelian varieties and their level structures in this description, for one type at the embedding
that sends $\sqrt{-d}$ to the fixed root.

% before: cv-components-points
Comparing a point with its conjugate requires knowing when two points of the uniformization coincide and
which homomorphisms the CM curve has to their abelian varieties. Passing from single points to whole
components requires regular functions that are constant along the components, and this comes from the
connectedness of the domains of periods.

% before: cv-components-split
The abelian variety of a point is built from the CM curve when the lattice is the sum of the two pure
sublattices of the period. Every lattice of the uniformization must have such a period, which holds in
definite signature and, in indefinite signature of rank at least $3$, as soon as there is a lattice of
signature $(1,0)$.

% before: cv-components-cm-point
At a point whose period splits, the conjugate abelian variety is built in the same way from the
conjugate of the CM curve. Comparing the two constructions gives an isometry between the lattices,
under which conjugation of torsion points becomes multiplication by the element of
Lemma~\ref{cv-components-multiplier}.

% before: cv-components-cm-fixed
When the automorphism fixes the CM curve with its points of order $M'$, that element is $1$, and the
isometry respects the level structures. The criterion for equality of points in
Lemma~\ref{cv-components-points} then shows that the point is fixed.

% before: cv-components-permutation
Each component of the uniformization is cut out by an idempotent regular function over a number field,
and such a function is constant along the components. So an automorphism of $\mathbf C$ that fixes
$\sqrt{-d}$ moves whole components onto whole components.

% before: cv-deep-components
The flat sections of Theorem~\ref{cv-flat-degree} are used on the cover, and they need the
coefficient field to be algebraically closed in the function field of every component. The Weber
values in the ray class field give this.

% before: cv-deep-components-literature
For Shimura varieties the fields of definition of the connected components are usually found from
Deligne's reciprocity law. The comparison places the lemma just proved in that setting, where the
Weber values enter through the ray class fields they generate.

% before: cv-deep-cover-exists
The ideal of Lemma~\ref{cv-deep-ideal}, the ray class field of Lemma~\ref{cv-deep-ray-field} and the
standard facts on PEL moduli schemes give a deep cover of every exponent. The size of its
coefficient field exceeds that of $F_c$ by a multiple of the logarithms of the deep level and of $d$.

% before: cv-deep-cover-sized
At a point of a family the exponent is taken of size $x^\varepsilon$. The size conditions and the
coefficient field bound then turn the size of the cover into a multiple of $z^2$, where
$z=frx^\varepsilon$, and this is the cover on which the slopes are bounded for $j<m$.

% before: cv-uniformization-torus-isometry
An isomorphism between two uniformized abelian varieties that respects $\mathcal O_E$ and the Weil
pairings comes from a linear map of $\mathbf C^g$, and then from an isometry of the lattices. This is how
a change of level structure, which keeps the abelian variety, is seen on the lattices of the
uniformization.

% before: cv-uniformization-model-match
Isometric lattices placed in $E^r$ by Hermitian models with the same matrix and the same element differ
by a unitary matrix. The isometries of the previous lemma therefore act on frames, through the embeddings
of $\Phi$.

% before: cv-uniformization-frames-connected
A change of level structure multiplies readings by a factor that may depend on the frame. To show that
the factor is constant, the readings are compared as holomorphic functions on the whole set of frames,
which needs this set to be connected.

% before: cv-uniformization-unit-minors
At a frame whose minors have absolute value $1$, the sections of the uniformization are normalized by a
mean over a fundamental domain. Every frame can be moved to such a frame by a diagonal element of $B$,
which changes readings only by the character of the weight.

% before: cv-uniformization-modulus
At frames with unit minors the comparison factor has absolute value $1$. Two holomorphic functions
on a connected open set whose absolute values agree everywhere then differ by a constant, by the maximum
modulus principle.

% before: cv-uniformization-factor-unique
The sections of the uniformization are read in local trivializations of the line bundle of the
polarization, through factors rigidified at the origin. Two such factors with the same transition law
and the same growth along the period lattice agree, so the factors of two framed points identified by
an isomorphism can be compared.

% before: cv-uniformization-level-lift
A change of level structure on the deep cover is an automorphism of the lattice modulo the deep level. On
the adelic quotient it is seen through an element of the unitary group over the finite adeles that
stabilizes the completed lattice. The primes of the level split in $E$, and at a split place the unitary
group is a general linear group, in which a matrix that is invertible modulo the level can be prescribed.

% before: cv-uniformization-scalar-orbits
The arithmetic group of a component is described through the transformations that keep every frame in
its orbit under the group and $B$. Such a transformation agrees with an element of the group up to a
scalar: by the Baire category theorem one element of the countable group serves on an open set of frames,
and the first columns of these frames determine the matrix up to a scalar.

% before: cv-uniformization-root-congruence
The scalar of the previous lemma turns out to be a root of unity, and a congruence modulo an integer at
least $3$ forces it to be $1$. The congruence is seen on eigenvalues: a matrix congruent to $1$ modulo $M$
on a lattice has eigenvalues $l$ with $(l-1)/M$ integral, and conjugate matrices share their eigenvalues.

% before: cv-uniformization-block-free
Two frames with the same point of the moduli scheme differ, after an element of the arithmetic group, by
a matrix that preserves the span of the first $q$ columns. Such a matrix is conjugate to an upper triangular
one by a matrix of the same shape, so the free action of the arithmetic group modulo $B$ applies to it.

% before: cv-uniformization-unitary-orbit
The frames attached to the adelic group have minors of absolute value $1$, and the real unitary group moves
them along whole orbits. A Hermitian form whose values have absolute value $1$ on such an orbit is a real
multiple of the form of that unitary group, so a matrix that preserves the first form preserves the second.

% before: cv-uniformization-uniform-constant
Passing to the level of $X$ multiplies each reading by a power of a constant of absolute value $1$ at every
frame, but the constant may a priori depend on the frame. Holomorphy on the connected set of frames and the
compactness of the unit circle give one constant for all sections and frames.

% before: cv-uniformization-theta-marking
To compare the readings of two sections of weight $(0,2)$, the cover needs one of them that does not vanish at a
given frame whose minors have absolute value $1$. Alp\"oge builds it from theta functions at torsion points, read
through the uniformization: a framed point at such a frame is a complex torus with a theta function, and its level
structure is read off a marking of the lattice of the component.

% before: cv-uniformization-theta-half-basis
The nonvanishing of a theta function at torsion points is proved through classical theta functions. Those of level
$N$ are the theta series with characteristics in $N^{-1}\mathbf Z^g$ at $N$ times the period matrix, and every
function with the transformation law of level $N$ is a combination of them.

% before: cv-uniformization-theta-half-point
At any given point some classical theta function of level $N\ge2$ does not vanish. A product of translates of the
Riemann theta function has the law of level $N$ and can be made nonzero at the point; it is a combination of the
classical ones.

% before: cv-uniformization-theta-half-classical
A function with the law of level $N\ge2$ in classical form cannot vanish at all the $N$-division points. Translating
by $1/N$ times an integer vector multiplies each classical theta function by its own character, and translating by
$1/N$ times a period moves one classical theta function to another, so vanishing at all these points would force a
classical theta function to vanish where none does.

% before: cv-uniformization-theta-half-normal-form
To apply the classical statement to the theta function of a framed point, its lattice and factor of automorphy are
put in classical form. A symplectic basis of the lattice for the imaginary part of the Hermitian form gives the
period matrix, and a quadratic exponential turns the factor into the classical one on the generators.

% before: cv-uniformization-theta-half-principal
Together these give the nonvanishing for a principal polarization in any coordinates. A holomorphic function with
the factor of automorphy of the $N$-th power of a principal polarization, for $N\ge2$, does not vanish at some
$N$-division point of the lattice.

% before: cv-uniformization-theta-half-sign-liouville
The unimodular part of the factor by which the theta function of a framed point changes under the lattice is found
by comparing two trivializations of one line bundle on $\mathbf C^g$. Their ratio is an entire function without
zeros that changes under the lattice by a unimodular constant times the exponential of a Hermitian form. Such a
function is constant and the form vanishes, by Liouville's theorem along complex lines.

% before: cv-uniformization-theta-half-sign-glue
The two trivializations come from the trivializing functions of a line bundle and from an isomorphism of the line
bundle with a tensor product of its pullbacks along endomorphisms. When the endomorphisms lift to holomorphic maps
of $\mathbf C^g$, the ratio is an entire function without zeros, and its change between two points with the same
image is read off from the trivializing functions.

% before: cv-uniformization-theta-half-sign-fibre
For an $\mathcal O_E$-semi-abelian scheme the line bundle is isomorphic to a tensor product of its pullbacks along
$0$, $1$, $w_0$ and $1+w_0$. Along an analytic fibre of the universal object, the entire function of this
isomorphism changes under the lattice by the exponential of a Hermitian form, so it is constant, and its values at
the lattice points give the unimodular factor. At a point of the flag domain the Riemann form is Hermitian with the
trace form as its imaginary part, and the periods of a self-dual lattice span $\mathbf C^g$ over $\mathbf R$.

% before: cv-uniformization-theta-half-sign
The factor by which the theta function of a framed point changes under the lattice has a unimodular part. The
normalization of the line bundle through the action of $w_0$ determines it: it is $\pm1$, according to the parity of
the pairing of a lattice vector with its product by $w_0$. The uniformization at a frame of a component is an
analytic fibre once its Hermitian form is identified with the Riemann form of the frame, which the real span of the
periods of the lattice forces.

% before: cv-uniformization-theta-half
The theta function of a framed point changes by a nonzero factor under the lattice, and it does not vanish at some
point of the torsion that the level structure sees. These two facts make the value of the section below nonzero.

% before: cv-uniformization-theta-rule-values
The section of the next lemma is defined by its values at framed points over algebras. Such a value is a
combination, with the coefficients of a harmonic polynomial, of products of Taylor coefficients at the origin of
translates of the theta section by roots of the level points, read after a faithfully flat extension that carries the
roots. The Taylor coefficients come from formal charts with a canonical
gauge, which makes them independent of the chart.

% before: cv-uniformization-theta-rule-connection-charts-frames
A chart datum starts from coordinates of the abelian scheme along its unit section at the point. The cubical
structure of the line bundle, evaluated at these coordinates, gives generators along them and the series of a formal
germ.

% before: cv-uniformization-theta-rule-connection-charts-identities
The series obtained from the cubical structure must satisfy the identities of a formal germ. These come from the
identities of the cubical structure and of the generators that express the Rosati involution.

% before: cv-uniformization-theta-rule-connection-charts-gauge
The Taylor values use a gauge, a power series whose exponentials give the series of the formal germ. Over a ring that
contains $\mathbf Q$ every formal germ has one.

% before: cv-uniformization-theta-rule-connection-charts-germs
A section of a power of the line bundle is read along the coordinates in the trivializations given by the generators.
It has a germ when these readings are the truncations of one power series, which the compatibility of the generators
with truncation provides.

% before: cv-uniformization-theta-rule-connection-charts
A canonical Taylor family needs, at every framed point, coordinates of the universal abelian scheme along its unit
section, generators of the line bundle along them and a gauge of their formal germ. Once such a chart datum exists and every
section has a germ, the Taylor values are given by a formula in the germ.

% before: cv-uniformization-theta-rule-connection
The values need one choice of Taylor coefficients for all framed points at once. The canonical gauge of a formal
chart removes the cocycles of the cubical structure, so the coefficients it gives are those of a connection.

% before: cv-uniformization-theta-rule-unique-base-change
Theta values are defined through images of framed points under homomorphisms of algebras. Images exist, and they
compose and cancel as the homomorphisms do.

% before: cv-uniformization-theta-rule-unique-translates-symmetric
A translate of the theta section uses an element of the theta group over a torsion point that is symmetric and has
order dividing that of the point. Both conditions say that a generator of a trivial line bundle, built from the
element by translations and the inversion, is $1$.

% before: cv-uniformization-theta-rule-unique-translates-lift
An element of the theta group over a point of odd order $n$ becomes symmetric, with $n$-th power $1$, after
multiplication by a constant. On an abelian scheme over any base the two defects of the element are constants, as
the global functions are those of the base, and they satisfy one relation, from which the constant is read off.

% before: cv-uniformization-theta-rule-unique-translates-theta
A translate of the generator is built from two pieces of the theta group of a power of the line bundle: an
isomorphism of the pullback under multiplication by $N$ with the $N^2$-th power, and a symmetric element over the root.
For odd $N$ both are unique once normalized, since an invertible function on an abelian scheme comes from the base.

% before: cv-uniformization-theta-rule-unique-translates-rosati
Two translates over the same point may come from different roots, which differ by a point killed by $N$. The
uniqueness of level translates needs a condition on the roots, which is stated here through a pair of endomorphisms
of the abelian scheme that are adjoint for the polarization.

% before: cv-uniformization-theta-rule-unique-translates
The theta values use translates of the generator by roots of the level points. Such translates exist over every
root of odd order, they do not depend on the root once it is killed by the square of the level ideal, and they are
compatible with base change.

% before: cv-uniformization-theta-rule-unique-chart
A canonical Taylor family is computed in one chart datum at each framed point, chosen among many. The combinations of
its Taylor values by a harmonic polynomial are the same for any two chart data at the same framed point, so the proofs
below may change the chart datum.

% before: cv-uniformization-theta-rule-unique-taylor-chart
To compare Taylor values at a framed point and at its image, a chart datum at the framed point is carried to the
image by base change. The Taylor values of the carried chart datum are the images of the original ones, since the
derivations and the gauge are only mapped coefficient by coefficient.

% before: cv-uniformization-theta-rule-unique-taylor
The Taylor values of a canonical family are computed in a chart datum chosen at each framed point. Their
combinations by a polynomial that satisfies the conditions on harmonic polynomials commute with base change.

% before: cv-uniformization-theta-rule-unique-compare
Two extensions that carry level translates give the same combination of Taylor values once both are mapped to a
common algebra. Applied to the two maps into a tensor product, this lets a theta value descend.

% before: cv-uniformization-theta-rule-unique-cover
Level translates exist after a faithfully flat extension. Multiplication by $M'$ is finite and flat, so roots of the
level points exist after one, and they can be chosen killed by the square of the level ideal.

% before: cv-uniformization-theta-rule-unique
A theta value is defined after a faithfully flat extension, and it has to descend to the algebra of the framed point.
Its independence of the extension, of the roots and of the formal chart gives exactly one value.

% before: cv-uniformization-theta-rule-weight-natural
A rational section must be natural under homomorphisms of test rings. Theta values are natural under all
homomorphisms of algebras, by the comparison of extensions.

% before: cv-uniformization-theta-rule-weight-commute
The Taylor values are built from the derivations of a formal group law that are invariant under translation. These
derivations commute with one another, as the group law is commutative, so a monomial in them does not depend on the
order of its factors.

% before: cv-uniformization-theta-rule-weight-substitution
A triangular change of the tangent vectors replaces each row of letter variables by a combination of that row and the
rows of smaller index. A polynomial that satisfies the row conditions of a harmonic polynomial changes under such a
substitution only by the diagonal entries raised to the row degrees.

% before: cv-uniformization-theta-rule-weight-operators
Over one formal germ, a combination of Taylor values of several power series is the value of a polynomial in
commuting derivations, one copy for each power series. A triangular change of the vectors that define the derivations
is then a substitution of the kind just treated.

% before: cv-uniformization-theta-rule-weight-taylor
A rational section of a given weight changes by a character under triangular changes of the tangent vectors and of
the generator. For the Taylor values the character comes from the row degrees of the harmonic polynomial.

% before: cv-uniformization-theta-rule-weight-triangular
The generator enters the theta values through the level translates, which are linear in it. Together with the change
of the Taylor values this gives the change of theta values under a triangular change.

% before: cv-uniformization-theta-rule-weight
Theta values are natural under homomorphisms of algebras, and a triangular change of the tangent vectors and of the
generator multiplies them by the character of the weight. These are the two conditions on a rational section.

% before: cv-uniformization-theta-rule-complex-charts
To descend a framed point to a finitely generated algebra, the base is covered by charts of the moduli scheme. Over
each chart there is a framed point, and the charts are spectra of test rings.

% before: cv-uniformization-theta-rule-complex-universal
Over a chart, every framed point differs from the framed point of the chart by invertible matrices on the bases and a
unit on the generator. Adjoining these entries to the test ring of the chart gives a test ring with a universal framed
point.

% before: cv-uniformization-theta-rule-complex-glue-injective
Two framed points with the same point differ by invertible matrices on the tangent bases and a unit on the
generator. These are seen after any injective base change, so framed points whose images agree there are equal.

% before: cv-uniformization-theta-rule-complex-glue-local
The same comparison works over a cover of the spectrum by principal open sets. An element of the ring that is $1$ on
every set of the cover is $1$, so framed points that agree on the cover are equal.

% before: cv-uniformization-theta-rule-complex-glue-cover
Framed points over a principal open cover that agree on the overlaps glue. Tangent vectors are morphisms from the
dual numbers and glue as morphisms do, and the generators of the sections glue as sections of the line bundle.

% before: cv-uniformization-theta-rule-complex-glue
The descended framed points are first found over localizations of the finitely generated algebra. They agree on the
overlaps, since their images in the localizations of the original algebra do, so they glue.

% before: cv-uniformization-theta-rule-complex-descent
A rational section is defined on framed points over test rings, and its value elsewhere is read through them. Every
framed point comes from a framed point over a finitely generated subalgebra.

% before: cv-uniformization-theta-rule-complex-tested
A finitely generated reduced algebra is isomorphic to a test ring. So framed points over reduced algebras, such as
complex framed points, are images of framed points over test rings.

% before: cv-uniformization-theta-rule-complex
The reading of a section at a frame is its value at a complex framed point, which is reached through a framed point
over a test ring. Naturality carries the theta value from the test ring to the complex point.

% before: cv-uniformization-theta-rule-translate
At a complex framed point the theta value becomes a value of theta series. A Heisenberg operator at a vector of
$E^r$ translates the theta series of a Schwartz--Bruhat function by the periods of that vector, up to an exponential
factor.

% before: cv-uniformization-theta-rule-product
When a product of functions is a multiple of a theta series, the translation law moves the identity to the
translated functions. The factors that appear are those of translates of theta functions.

% before: cv-uniformization-theta-rule-fibre-reading-generator
Each piece of a translate is an isomorphism between pullbacks of powers of the line bundle, and through the analytic
fibre it is read as a holomorphic function times the trivializing functions. When the Riemann forms of the pullbacks
cancel, this function has a periodic absolute value, so it is constant, and the multipliers of the fibre cancel too.

% before: cv-uniformization-theta-rule-fibre-reading-sign
A translate is read through the analytic fibre one piece at a time, and the reading of the symmetric element of the
theta group depends on the multiplier of the fibre on the lattice. This multiplier is a sign, given by the imaginary
part of the Riemann form at a lattice vector and its product by $w_0$.

% before: cv-uniformization-theta-rule-fibre-reading-pullback
The first piece of a translate is the pullback of the theta section by multiplication by $N$, carried to the
$N^2$-th power of the line bundle by the normalized isomorphism. Through the analytic fibre it is read as the theta
reading at $N$ times the variable.

% before: cv-uniformization-theta-rule-fibre-reading-symmetric
The second piece of a translate is a symmetric element of the theta group over the root, the image of a vector of
$\mathbf C^g$. When an odd multiple of the vector is twice a period of the lattice, the element acts on theta readings
by translation by that vector, with an exponential factor.

% before: cv-uniformization-theta-rule-fibre-reading
Over $\mathbf C$ the analytic fibre reads sections of powers of the line bundle as functions on $\mathbf C^g$. A
translate by a root $\operatorname{ex}(p_x(a))$ is read as a translate of the theta reading, at $N$ times the variable.

% before: cv-uniformization-theta-rule-fibre-taylor-chart
At a complex framed point the analytic fibre gives a chart datum of its own: the coordinates of $\mathbf C^g$ are a
formal chart with the additive group law, and the gauge is $0$. In it the germ of a section is the Taylor series of
its theta reading.

% before: cv-uniformization-theta-rule-fibre-taylor
The Taylor values of a canonical family are defined through formal charts. At a complex framed point their harmonic
combinations are derivatives at the origin of the theta readings of the sections.

% before: cv-uniformization-theta-rule-fibre
Over $\mathbf C$ the level translates are read through the analytic fibre as translates of the theta reading, and
their Taylor coefficients at the origin become derivatives of these functions. With a vacuum of the tensor lattice,
the product of the translates is a constant times a theta series of a Heisenberg translate of the vacuum.

% before: cv-uniformization-theta-rule-one
The section of the next lemma is a theta rule for the constant polynomial $1$. This polynomial satisfies the
conditions on harmonic polynomials with all row degrees $0$.

% before: cv-uniformization-theta-rule
The section is a theta section of points of the lattice of the datum. At a frame whose minors have absolute value
$1$ its reading is a product of values of the theta function of the framed point at lifts of the corresponding
torsion points.

% before: cv-uniformization-theta-section
Choosing all the points in the class of one point at which the theta function does not vanish gives a section of
weight $(0,n)$ whose reading at the frame is a product of nonzero numbers. With $n=2$ this is the section that the
comparison of readings on the cover needs.

% before: cv-uniformization-fibre-count-tangent
Two level changes that give the same framed point have isomorphisms of abelian varieties that agree on a basis of
the tangent space. In characteristic $0$ this forces them to agree, since the doubling map acts on the cotangent
space of the origin by $2$ and on its higher powers by higher powers of $2$.

% before: cv-uniformization-fibre-count
Passing from the cover down to $X$ multiplies the volumes of the components by the number of framed points over a
framed point. The level changes that are trivial modulo the level of $X$ move these framed points simply
transitively, so this number does not depend on the framed point.

% before: cv-cover-uniformization
Cusp forms and Petersson norms on the cover are read through its uniformizations. Changing level
structures, and passing from the cover down to $X$, move the components of these uniformizations by
unitary transformations, and the next lemmas compare readings, arithmetic groups and measures through
these transformations.

% before: cv-cover-level-automorphisms
On the deep cover a finite group of level automorphisms acts on cusp forms by changing level
structures. The descent between weights averages over a subgroup of this group and multiplies forms
by an auxiliary form, so it needs the group to be small, its action to preserve cusp forms,
integrality and Petersson norms, and products of cusp forms to be cusp forms.

% before: cv-cover-pullback
The lattices on the cover are compared with those of the Shimura variety by pulling forms back along
the level map. The comparison of slopes in Lemma~\ref{cv-lattices-different-slope} needs the pullback
to be semilinear, injective on independent families, isometric for the Petersson norms and
compatible with integrality, and it needs integrality to be stable under the Galois group of $F_Y$
over $F_c$.

% before: cv-cover-type
The forms of the prescribed type on the Shimura variety are eigenforms of the level subgroup. On the
cover, the span of their pullbacks is cut out by the same condition for the preimage of the level
subgroup in the larger group of the cover, on whose kernel the condition is invariance. This turns
the lattice of the cover into an eigenlattice, to which the averaging of the descent applies.

% before: cv-volume-forms-weights
The descent below compares forms whose weights differ by multiples of the Hodge weight and by the
anchor. Each of these weights is given by a family of integers along the flag and a power of the
determinant, and the arguments for their integral forms use that the family decreases, with a strict
drop where the signature changes from $p$ to $q$.

% before: cv-volume-forms-finite
Finite generation of the integral forms of a weight presupposes that the space of forms of that weight
is finite-dimensional. Through the readings at a good embedding, this comes down to cusp forms of a
fixed weight on the components of a uniformization.

% before: cv-volume-generators-functionals
Finite generation of the integral forms is reduced to linear forms that take integral values on them.
It suffices that every nonzero form is seen by one such linear form; values at a point over a number
field, multiplied by a fixed denominator, will be used.

% before: cv-volume-generators-power
That the integral forms span is reduced to a power of a form, which can be moved into a weight where
every form has an integral multiple. An integral multiple of the power gives one of the form itself,
because the valuations at integral points add under products.

% before: cv-volume-twisting-section
In the deep cover, the integral forms of two weights that differ by the anchor weights and the entry
$n$ are compared on each component. The comparison multiplies by a section of the difference weight whose
norm at the frames is constant and whose valuations at the integral points are bounded. The degree of such
a twisting section measures how much it moves the degree of a lattice of forms.

% before: cv-volume-twisting-products
The twisting section of the comparison is assembled from two simpler ones. Products and powers of
twisting sections are again twisting sections, and their degrees add.

% before: cv-volume-twisting-points
A twisting section asks for frames and integral points on its component. Both exist on every component
when the level is good at every prime.

% before: cv-volume-twisting-coset
The section of the next lemmas is built from values of the theta function at points of finite order. Such
values help only if the theta function does not vanish on a whole coset of the $n$-torsion, and the
irreducibility of the theta group rules this out.

% before: cv-volume-twisting-combination
A product over a finite group of linear combinations is nonzero when each factor is. Integer coefficients of
bounded size make all the factors nonzero at once.

% before: cv-volume-twisting-level-structures
The theta null values are attached to points of finite order of the abelian varieties of $X$. A full level
structure labels these points by the lattice of the PEL datum, so that expressions in the labels can be
compared from one point of $X$ to another.

% before: cv-volume-twisting-levels
Full level structures exist only after a faithfully flat extension, and two of them differ by an automorphism
of the labels. With these facts an expression in level structures that does not change under relabelling
descends to a rational section.

% before: cv-volume-twisting-null-values
The null values are the values at the origin of translates of the pulled-back theta function by points of
order dividing $81$. The normalized pullback isomorphism and the symmetric elements of the theta group make
them well defined.

% before: cv-volume-twisting-null-facts
The section built from the null values must not depend on choices, and it must have a weight. For this the
null values commute with base change and scale like the generator of sections; they are also nonzero when the
theta function does not vanish at the corresponding point.

% before: cv-volume-twisting-eta
A product over all relabellings of integer combinations of products of null values does not depend on the
level structure. Faithfully flat descent turns it into a rational section, whose weight is given by its
degree in the generator of sections.

% before: cv-volume-twisting-nonvanishing
The readings of a section of a weight with negative entries are not covered by the compatibility conditions
of the moduli problem. They are reached by dividing by a section of dominant weight, and this needs a section
of dominant weight that does not vanish at a given frame.

% before: cv-volume-twisting-holomorphic
Multiplying a section by powers of the key section and of a section that does not vanish at a frame gives a
section of dominant weight. Near that frame the reading of the first section is then a quotient of two
forms, so it is a form itself.

% before: cv-volume-twisting-cusp
A twisting section must carry cusp forms to cusp forms. This holds for a section with bounded norm whose
weight has entries that do not increase, once the moduli scheme has a key section.

% before: cv-volume-twisting-flat
The part of the twisting section given by the anchor weights comes from the flat sections of
Definition~\ref{cv-constant-sections}. A product of their powers has constant norm and constant
valuations on a component, and its degree is computed from the degrees of the flat sections and the
defects.

% before: cv-volume-twisting-colength
The part of the twisting section given by the entry $n$ is a fourth root of a section built from the key
formula. Its valuations at an integral point depend on how far the eigenvectors of $\mathcal O_E$ in the
Lie algebra over a discrete valuation ring are from a basis, and the different of $E$ bounds this.

% before: cv-volume-twisting-key-unit
The key section is built from the values of the key rule on families over the coordinate rings of charts.
Over a noetherian integrally closed base these values are units, as they are over every discrete valuation
ring by the key formula.

% before: cv-volume-twisting-key-theta
The fourth root of the key section is found from the complex points of the uniformization. At its normalized
framed points the key rule is a constant, by the normalization of the classical theta functions in the key
formula.

% before: cv-volume-key-fixed-field
A fourth root of the key rule is found first over $\mathbf C$ and then brought down to the coefficient field.
Bringing it down uses that the automorphisms of $\mathbf C$ over a subfield fix no element outside that
subfield.

% before: cv-volume-key-descent-spaces
A subspace of a complexified vector space that is stable under the automorphisms of $\mathbf C$ over $K$ is
spanned by vectors defined over $K$. Applied to the line spanned by a complex root, this gives a root over $K$
up to a constant.

% before: cv-volume-key-roots-descent
An $n$-th root of $c\otimes u$ in the complexification determines an $n$-th root of $u$ up to a scalar of the
base field. The line it spans is stable under the automorphisms of $\mathbf C$ over $K$, since such an
automorphism changes the root by a constant.

% before: cv-volume-key-connected
A continuous choice of $n$-th roots on a connected space differs from any other continuous choice by a
constant root of unity. The spaces of complex points used here are those of integral varieties, which are
connected.

% before: cv-volume-key-continuity
The fourth roots of the key rule are known to vary continuously along holomorphic families of complex points.
On a smooth variety the holomorphic charts are such families, so this gives continuity on all the complex
points.

% before: cv-volume-key-roots-complex
A unit with a continuous $n$-th root on the complex points of a smooth complex variety is an $n$-th power.
Adjoining an $n$-th root gives a domain whose complex points are connected, and the continuous root leaves
room for only one root over each point of the variety.

% before: cv-volume-key-roots
Over a subfield $K$ of $\mathbf C$, the root over $\mathbf C$ given by Lemma~\ref{cv-volume-key-roots-complex}
comes down to $K$ up to a constant by Lemma~\ref{cv-volume-key-roots-descent}. On each chart of the moduli
scheme this turns the continuous fourth roots of the key rule into a fourth root of the unit given by the key
rule.

% before: cv-volume-key-theta-point
The continuous fourth roots of the key rule come from the uniformization rebuilt from theta data. A complex
framed point is reached from a normalized framed point by a triangular change $t$, and the number $\chi(t)$ is
the fourth root in question. A theta point carries the theta function of a normalized framed point with another
scalar in front, so its change can be read off from the two scalars.

% before: cv-volume-key-root-values
A complex framed point is reached from normalized framed points along many changes, one for each frame in an
orbit of $\Gamma_i$ and $B$. They all have the same number $\chi$, since $\chi$ is the quotient of the characters
of two dominant weights on which $\Gamma_i$ and $B$ act in the same way. At a theta point this number is given
by its scalar.

% before: cv-volume-key-adapted-periods
To follow root values along a holomorphic family of framed points, the family is first presented by periods.
Composing the periods with a matrix in each fibre makes the bases of the framed points the images of the unit
vectors, and this matrix depends holomorphically on the parameter, as the bases come from a framed point over an
algebra.

% before: cv-volume-key-frame-lift
A period determines its frame only up to the columns and rows on which it does not depend. Frames near a given
one with the nearby periods are found by correcting the given frame in those columns and rows, and this gives a
choice of frames that is continuous at a point.

% before: cv-volume-key-theta-scalar
Along such a choice of frames every framed point of the family is a theta point, and only its scalar remains.
The scalar is the quotient of a value of the section of the family, which is holomorphic in the parameter, by a
value of the theta function and its factor, which is continuous.

% before: cv-volume-key-root-continuity
At a theta point the root value is a continuous expression in the frame and the scalar. So the root values of a
holomorphic family of framed points vary continuously, which is the hypothesis of
Lemma~\ref{cv-volume-key-roots} on the charts of the moduli scheme.

% before: cv-volume-key-sections
The charts of the moduli scheme carry framed points over rings, which are specialized along homomorphisms. The
generator of a framed point must remain a generator after specialization. On an abelian scheme with a principal
polarization this follows from the behaviour of the sections of the powers of the line bundle under base change.

% before: cv-volume-key-generator
With one generating section that survives every base change near each point, generating sections are recognized
fibre by fibre. A section of the line bundle of the polarization generates exactly when it does not vanish
identically on any fibre, and this condition passes to pullbacks.

% before: cv-volume-key-lie
The bases of a framed point live in the eigenspaces of the Lie algebra. Over the coefficient field the action of
$\mathcal O_E$ splits the Lie algebra into these eigenspaces, and as the abelian scheme is smooth the Lie algebra
is locally free and compatible with base change, which gives local eigenbases and their dimensions.

% before: cv-volume-key-base-change
A framed point has an image under every homomorphism. Its bases go to bases by the facts on the Lie algebra, and
its generator goes to a generator, as it does not vanish identically on any fibre and this condition survives base
change.

% before: cv-volume-key-generic-points
Values of sections at framed points are compared through one framed point on each component. It lies over the
generic point of the component, with its bases moved by a unipotent change whose entries are independent
variables, so that a section vanishing there vanishes at every framed point of the component.

% before: cv-volume-key-chart
A chart is a smooth algebra inside the field of the generic framed point, with a framed point over it that
specializes to every framed point of a part of the component. The two conditions on it say that this
specialization is determined by the framed point and that the part is open.

% before: cv-volume-key-local-frames
Charts are built from affine neighbourhoods over which the eigenspaces of the Lie algebra have bases and the line
bundle of the polarization has a generating section. Over a smooth base such a neighbourhood can be taken to be the
spectrum of a domain.

% before: cv-volume-key-reframe
Two framed points with the same point differ by invertible matrices on the eigenspaces of the Lie algebra. Over a
field each such matrix factors, after a permutation of its rows, into a unitriangular and an upper triangular
one, which separates the part of the change that a chart must absorb from a triangular change.

% before: cv-volume-key-cells
A chart is attached to a cell of matrices: those that are a matrix of a fixed unitriangular shape, after a
permutation of the rows, times an upper triangular matrix. Membership in the cell is the invertibility of certain
minors, an open condition, and the coordinates on the cell are unique.

% before: cv-volume-key-cell-change
The chart for a cell must contain the generic framed point. The change from a framed point over the generic point to
the generic framed point lies in every cell, and its coordinates in the cell are new independent variables that
generate the same field.

% before: cv-volume-key-cell-charts
For each permutation of the rows, the unitriangular parts of these factorizations form an affine cell. A
polynomial ring in coordinates for this cell over the ring of a framed affine neighbourhood, embedded in the field
of the generic framed point, gives a chart that reaches every framed point whose change factors through this
cell.

% before: cv-volume-key-atlas
The charts for the cells over all framed affine neighbourhoods reach every framed point over a field. A chart that
reaches the generic framed point of a component lives on that component, since the generic point of a chart
specializes to every point it reaches.

% before: cv-volume-key-chart-compat
A section is glued from candidate values on the charts, which must agree where two charts reach the same framed
point. Two charts that reach one framed point over a field lie on the same component, and the values they give
there differ by the character of the triangular change between the two images.

% before: cv-volume-key-local-values
The value that a chart gives at a framed point it reaches is a local value. It is unique over a field by the
compatibility of charts, it is carried along base changes and triangular changes, and at the generic framed point
of a component it is the prescribed value.

% before: cv-volume-key-representatives
Over an arbitrary algebra the local values at the field points must come from one element of the algebra. The
elements that work after multiplication by some factor form an ideal, and the chart property, which extends a chart
from a field point to the locus where some $g$ is invertible, puts a power of $g$ in it.

% before: cv-volume-key-gluing
Choosing such an element at every framed point over a test ring defines a rational section. Since test rings are
reduced, an element is determined by its values at field points, which gives the compatibility with base change
and the transformation law.

% before: cv-volume-key-xi
The normalized points are rebuilt from Hermitian models of the lattices of the uniformization family. These need
an element $\xi_0$ with $\bar\xi_0=-\xi_0$ whose imaginary part has a fixed sign at the CM type, and Hermitian
coordinates for the lattice of the datum, obtained by writing its alternating form as a trace.

% before: cv-volume-key-models-half
The level structure of $X$ fixes a lattice modulo $M$ only on the part $\bar{\mathfrak M}L$. For automorphisms that
preserve the perfect form this already forces the congruence on the whole lattice, since $\mathfrak M$ and
$\bar{\mathfrak M}$ are coprime; Minkowski's lemma can then be applied.

% before: cv-volume-key-models-volume
The Hasse norm theorem needs the quotient of the determinants of two Hermitian matrices to be a local norm at every
finite place. Alp\"oge obtains this from ideals generated by determinants: when a lattice of row vectors is the dual of
another for a matrix $A$, these ideals multiplied by $\det A$ give a power of the ideal in which the pairing takes its
values.

% before: cv-volume-key-models-dual
The lattice of a Hermitian model is its own dual for the pairing $\operatorname{Tr}_{E/\mathbf Q}(\eta\,h_G)$. This
follows from the perfectness of the alternating form, and it is the duality that the determinant identity requires.

% before: cv-volume-key-models-ideal
Two Hermitian models with the same element $\eta$ have self-dual lattices for pairings that differ only through their
matrices. The determinant identity for each of them shows that the quotient of the two determinants is a quotient of
norms of ideals from $E$.

% before: cv-volume-key-models-generator
To pass from norms of ideals to norms of elements at one prime, each ideal is replaced by one of its elements. Nakayama's
lemma gives an element that generates the ideal after multiplication by a factor congruent to $1$ modulo a given ideal,
and such a factor lies outside every prime that contains this ideal.

% before: cv-volume-key-models-unit
At a prime $\mathfrak p$ of $E^+$ these elements turn an identity between norms of ideals into one between norms of
elements, up to a factor. The factor is fixed by complex conjugation and prime to $\mathfrak p$, so it is an integer of
$E^+$ prime to $\mathfrak p$.

% before: cv-volume-key-models-binary
This factor is a norm locally because $E$ is unramified over $E^+$. A binary quadratic form whose discriminant is a unit
at $\mathfrak p$ represents every element of the finite residue field, and a contraction argument in the completion
lifts such a representation of a unit.

% before: cv-volume-key-models-units
Norms from $E=E^+(\sqrt{-d})$ are values of the form $a^2+db^2$ over $E^+$. As $\mathfrak p$ is unramified in $E$, an
integer $w$ of $E$ with $w-\bar w$ prime to $\mathfrak p$ gives a binary form of unit discriminant with the same values,
and it represents the integers of $E^+$ prime to $\mathfrak p$ in the completion.

% before: cv-volume-key-models-local
The factorization at each prime and the local representation of units combine into the local condition of the Hasse
norm theorem at the finite places. An element of $E^+$ whose ideal is a quotient of norms of ideals from $E$ is a norm
from $E\otimes_{E^+}E^+_v$ at every finite place $v$.

% before: cv-volume-key-models
The lattices of the uniformization family at a good embedding share the signature of the datum. With the Hasse norm
theorem and Landherr's theorem they all have Hermitian models for one matrix $G$, and the congruence groups of these
models at the level of $X$ have the properties needed for the frames.

% before: cv-volume-key-rebuilt
The models define frames, periods and, through theta functions with characteristics, framed points whose sections
are normalized by the determinant of the change to a symmetric period matrix. Their key values can be computed
from the theta functions, which is why they replace the framed points of the given uniformization.

% before: cv-volume-key-normalized-tangent
A framed point needs a basis of each eigenspace of the Lie algebra. The uniformization identifies $\mathbf C^g$ with
the Lie algebra through the derivative of its map at the origin, and the action of $\mathcal O_E$ through the type
sorts the unit vectors into the eigenspaces.

% before: cv-volume-key-normalized-theta
The section of a normalized framed point is a theta function times a nowhere vanishing factor. It generates the
sections of the line bundle as soon as it is nonzero at one point, and the constant Fourier coefficient of the theta
function provides such a point.

% before: cv-volume-key-normalized-gram
The periods of a frame are built from the Hermitian form $\tau(G)$ at each embedding $\tau\in\Phi$. A unitriangular
change of the columns of a frame makes them orthogonal for this form, with signs fixed by the minors.

% before: cv-volume-key-normalized-periods
Theta data need the periods of the frames to form a holomorphic family of periods of the lattices, and the cover of
the framed points needs every period to come from a frame. Both follow from the orthogonal columns of the last lemma
and from the trace form, whose imaginary parts at the embeddings of $\Phi$ add up to the alternating form.

% before: cv-volume-key-normalized-riemann
Theorem~\ref{cv-uniformization-family}(4) attaches theta functions to a family of periods once the period matrix is
brought to the form of a unit matrix and a matrix $Z$. Riemann's relations show that $Z$ lies in the Siegel upper half
space and that it depends holomorphically on the frame.

% before: cv-volume-key-normalized
Theta data come from a symplectic basis adapted modulo $8$ and from the theta functions that the uniformization
family attaches to holomorphic families of periods. With them a normalized framed point exists at every frame.

% before: cv-volume-key-cover
Every complex framed point is reached from a normalized framed point by a triangular change. The point lies in the
uniformization family, and a triangular change of coordinates on $\mathbf C^g$ turns its tangent vectors into the
images of the unit vectors at the period of a frame.

% before: cv-volume-twisting-key-section
The rule of the key formula of Theorem~\ref{cv-definite-key-formula}, evaluated at framed points, is a
section of weight $\kappa(0,-4)$ on the moduli scheme. On each component it is a fourth power up to a
constant, and its norm at the frames is a constant controlled by $2\pi$.

% before: cv-volume-twisting-key-valuation
At an integral point the section of the key formula is the fourth power of the determinant that compares
the frame with a basis of the Lie algebra over the discrete valuation ring, times a unit.
Lemma~\ref{cv-volume-twisting-colength} bounds this determinant by the different of $E$.

% before: cv-volume-twisting-key
A fourth root of the section of the key formula gives the part of the twisting section of weight
$\kappa(0,1)$. Its degree is bounded below through the norms of the key section at the frames and the
different of $E$, which contribute $\frac12\log2\pi$ and $\log\operatorname{rd}_E$ per dimension.

% before: cv-volume-twisting-unramified
The functions of the twisting sections take integer values away from the bad primes only if every
component has integral points there with integer valuations. Such points come from CM points defined over
the compositum of the coefficient field with a ray class field of $\mathbf Q(\sqrt{-d})$, which is
unramified over the coefficient field at some prime above the given place.

% before: cv-volume-forms-cm-frame
Values of a form at points over number fields are reached through a frame at which the form does not
vanish. Near a frame where its pointwise norm is close to the maximum, Lemma~\ref{cv-cusp-frame} gives such
a frame adapted to a basis of the lattice of the component, the shape of frame at which CM points are found.

% before: cv-volume-forms-triangular
As the frame varies over one point, the value of a form is a polynomial in the change of frame divided by a
power of its determinant. Equivariance under upper triangular matrices decides how many factors of the
determinant the polynomial contains.

% before: cv-volume-forms-frame-polynomial
The value of a rational section at the generic change of frame of a fixed framed point is a polynomial
divided by a power of the determinant. With the transformation law of the weight, the polynomial can be
taken homogeneous, with a power of the determinant fixed by the weight.

% before: cv-volume-forms-frame-nonzero
A frame at which a form does not vanish is first found over the complex numbers. The values of the form at
the frames over that point are given by a nonzero polynomial, so some frame defined over the number field
of the point also gives a nonzero value.

% before: cv-volume-forms-denominator
At a CM point, integral forms are integral only up to a denominator. If the polynomials that give their
values at the frames over the point have a common denominator, evaluating them at one frame adds a
denominator that depends only on that frame.

% before: cv-volume-generators-exponents
To move a form to a weight where every form has an integral multiple, a power of it is multiplied by
sections whose weights are given by a constant family and by the entry $n$. The strict drop of a regular
family where the signature changes leaves room for a constant shift that makes the total weight dominant.

% before: cv-volume-generators-integers
Short generators of a lattice of forms are built from short integers of the coefficient field. The
transference bound applied to the trace form gives a basis of the field over $\mathbf Q$ made of integers
whose conjugates are bounded by the degree and the root discriminant.

% before: cv-volume-generators-short
A family of short forms that spans the space over the coefficient field yields generators of the integral
forms over $\mathbf Z$. Reducing a basis of the lattice modulo the short family keeps the generators short,
at the cost of factors given by the degree, the rank and the root discriminant.

% before: cv-volume-generators-key-valuation-level
The valuation bound of Lemma~\ref{cv-volume-twisting-key-valuation} needs a level whose prime-to-$\ell$ part is at least
$3$ for every prime $\ell$. Every level has a finer one of this kind with the same full level: multiplying
$\mathfrak M$ by an element $\alpha$ of norm $m$ prime to $M$ gives the integer $Mm$, and no prime divides both $M$
and $m$.

% before: cv-volume-generators-key-valuation-lift
A point at which a valuation is computed must then be lifted to the moduli scheme at the finer level. The level map is
finite and surjective once the coefficient field contains enough roots of unity, so the lift exists over a finite
extension of the fraction field. The integral closure of the complete discrete valuation ring in this extension is
again a complete discrete valuation ring, whose valuation extends the old one up to the ramification index.

% before: cv-volume-generators-key-valuation
The valuation bound for the section of the key formula assumed a level whose prime-to-$\ell$ part is at least
$3$ for every prime $\ell$. Its lower half holds at every level. At a prime where the level is too small,
Alp\"oge passes to a finer level, where the value of the key formula is a unit, and descends.

% before: cv-volume-generators-key-inverse
Moving a form to a larger weight multiplies it by a power of a section of weight $\kappa(0,1)$, which must
have an integral inverse up to a constant. The inverse of the key section is such a section: on each
component the fourth power of the key section is a constant times the section of the key formula, which has
nonnegative valuations at the integral points.

% before: cv-volume-theta-index-ideal
A theta series of a lattice becomes integral after multiplication by a scalar that compares the determinants
of the lattice with those of its intersection with $\mathcal O_E^n$. Over the ring of integers these
determinant modules differ by an ideal whose norm is the index of the two lattices.

% before: cv-volume-theta-small-scalar
The scalar must also be small at every archimedean place, so that the integral multiple of a theta series
keeps a bounded norm. Minkowski's theorem for the inverse of the comparison ideal gives such a scalar, at a
cost controlled by the degree and the discriminant of the coefficient field.

% before: cv-arithmetic-volume-step
The lower bound for slopes of integral automorphic forms is where the arithmetic volume
estimate enters. For the last lift the group is definite and the bound comes from CM
points alone; for the others it comes from a deep cover and a descent between weights.

% before: cv-covolume-bound-from-literature
The steps above use a list of published results from arithmetic geometry. Assuming them,
the covolume bound follows by combining the six properties of
Proposition~\ref{cv-point-conditions} through the cut of Lemma~\ref{cv-cut}.

% before: cv-bound-as-used
A one-sided bound over families of points is only useful if it is applied to points that
exist and satisfy its hypotheses. Both conditions can be checked at the place where the
bound is used, and the error term is absorbed there into a budget of the same order.
